\input{preamble}

% OK, start here.
%
\begin{document}

\title{Le complexe cotangent}


\maketitle

\phantomsection
\label{section-phantom}

\tableofcontents

\section{Introduction}
\label{section-introduction}

\noindent
Le but de ce chapitre est de construire le complexe cotangent d'un
homomorphisme d'anneaux, d'un morphisme de schémas et d'un morphisme d'espaces
algébriques. On pourra consulter les notes \cite{quillenhomology}, l'article
\cite{quillencohomology}, ainsi que les livres \cite{Andre} et
\cite{cotangent}.




\section{Conseils au lecteur}
\label{section-advice-reader}

\noindent
En rédigeant ce chapitre, nous avons cherché à réduire au minimum l'emploi
des techniques simpliciales. Nous considérons le choix d'une
{\it résolution} $P_\bullet$ d'un anneau $B$ sur un anneau $A$ comme un outil
pour calculer l'{\it homologie} des faisceaux abéliens sur la catégorie
$\mathcal{C}_{B/A}$ ; voir la remarque \ref{remark-resolution}. Ce rôle est
analogue à celui d'un ``bon recouvrement'' dans le calcul de la cohomologie au
moyen du complexe de {\v C}ech. Pour quelques notions d'homologie sur les
catégories, voir Cohomologie sur les sites, section
\ref{sites-cohomology-section-homology}. Le foncteur dérivé gauche $L\pi_!$
joue pour l'homologie le rôle que $R\Gamma(\mathcal{C}_{B/A}, -)$ joue pour
la cohomologie. La catégorie $\mathcal{C}_{B/A}$, étudiée dans la section
\ref{section-compute-L-pi-shriek}, est l'opposée de la catégorie des
factorisations $A \to P \to B$, où $P$ est une algèbre polynomiale sur $A$.
Cette catégorie est munie d'homomorphismes de faisceaux d'anneaux
$$
\underline{A} \longrightarrow \mathcal{O} \longrightarrow \underline{B}
$$
où, sur l'objet $U =  (P \to B)$, on a $\mathcal{O}(U) = P$.
On obtient alors le complexe cotangent de $B$ sur $A$ sous la forme
$$
L_{B/A} =
L\pi_!(\Omega_{\mathcal{O}/\underline{A}} \otimes_\mathcal{O} \underline{B})
$$
voir le lemme \ref{lemma-compute-cotangent-complex}. Nous nous sommes efforcés
d'adopter systématiquement ce point de vue pour démontrer les propriétés
élémentaires des complexes cotangents d'homomorphismes d'anneaux. Tous les
résultats peuvent notamment se démontrer sans recourir à l'existence de
résolutions standard, bien que nous ne l'ayons pas fait. La théorie est tout à
fait satisfaisante, à ceci près que la démonstration du triangle fondamental
(proposition \ref{proposition-triangle}) utilise peut-être un peu plus de
théorie des foncteurs dérivés gauches. Afin d'offrir une autre voie au
lecteur, nous donnons dans les remarques \ref{remark-triangle} et
\ref{remark-explicit-map} une esquisse assez complète d'une approche fondée
sur des propriétés simples des résolutions standard.

\medskip\noindent
Pour traiter le complexe cotangent des morphismes de topos annelés, des
morphismes de schémas, des morphismes d'espaces algébriques, etc., notre
méthode consiste à déduire autant que possible du cas des ``simples
homomorphismes d'anneaux'' étudié ci-dessus.





\section{Le complexe cotangent d'un homomorphisme d'anneaux}
\label{section-cotangent-ring-map}

\noindent
Soit $A$ un anneau. Notons $\textit{Alg}_A$ la catégorie des $A$-algèbres.
Considérons le couple de foncteurs adjoints $(U, V)$ où
$V : \textit{Alg}_A \to \textit{Ensembles}$ est le foncteur d'oubli et où
$U : \textit{Ensembles} \to \textit{Alg}_A$ associe à un ensemble $E$ l'algèbre
polynomiale $A[E]$ sur $E$ au-dessus de $A$. Soit $X_\bullet$ l'objet
simplicial de $\text{Fun}(\textit{Alg}_A, \textit{Alg}_A)$ construit dans
Objets simpliciaux, section \ref{simplicial-section-standard}.

\medskip\noindent
Considérons une $A$-algèbre $B$. Notons $P_\bullet = X_\bullet(B)$ la
$A$-algèbre simpliciale ainsi obtenue. Rappelons que $P_0 = A[B]$,
$P_1 = A[A[B]]$, et ainsi de suite. En particulier, chaque terme $P_n$ est
une $A$-algèbre polynomiale. Rappelons aussi qu'il existe une augmentation
$$
\epsilon : P_\bullet \longrightarrow B
$$
où l'on considère $B$ comme une $A$-algèbre simpliciale constante.

\begin{definition}
\label{definition-standard-resolution}
Soit $A \to B$ un homomorphisme d'anneaux. La {\it résolution standard de
$B$ sur $A$} est l'augmentation $\epsilon : P_\bullet \to B$ dont les termes
sont
$$
P_0 = A[B],\quad P_1 = A[A[B]],\quad \ldots
$$
et dont les applications sont construites dans Objets simpliciaux, exemple
\ref{simplicial-example-polynomial-algebra-maps}.
\end{definition}

\noindent
Nous verrons que la résolution standard permet, dans certaines situations,
de calculer les foncteurs dérivés gauches.

\begin{definition}
\label{definition-cotangent-complex-ring-map}
Le {\it complexe cotangent} $L_{B/A}$ d'un homomorphisme d'anneaux $A \to B$
est le complexe de $B$-modules associé au $B$-module simplicial
$$
\Omega_{P_\bullet/A} \otimes_{P_\bullet, \epsilon} B
$$
où $\epsilon : P_\bullet \to B$ est la résolution standard de $B$ sur $A$.
\end{definition}

\noindent
Dans Objets simpliciaux, section \ref{simplicial-section-complexes}, nous
associons un complexe de chaînes à un module simplicial ; ici, toutefois, nous
travaillons avec des complexes de cochaînes. Ainsi, le terme $L_{B/A}^{-n}$
en degré $-n$ est le $B$-module
$\Omega_{P_n/A} \otimes_{P_n, \epsilon_n} B$, et $L_{B/A}^m = 0$ pour
$m > 0$.

\begin{remark}
\label{remark-variant-cotangent-complex}
Soit $A \to B$ un homomorphisme d'anneaux. Soit $\mathcal{A}$ la catégorie
des flèches $\psi : C \to B$ de $A$-algèbres, et soit $\mathcal{S}$ la
catégorie des applications $E \to B$, où $E$ est un ensemble. On dispose de
foncteurs adjoints $V : \mathcal{A} \to \mathcal{S}$ (le foncteur d'oubli) et
$U : \mathcal{S} \to \mathcal{A}$, qui envoie $E \to B$ sur
$A[E] \to B$. Soit $X_\bullet$ l'objet simplicial de
$\text{Fun}(\mathcal{A}, \mathcal{A})$ construit dans Objets simpliciaux,
section \ref{simplicial-section-standard}. Le diagramme
$$
\xymatrix{
\mathcal{A} \ar[d] \ar[r] & \mathcal{S} \ar@<1ex>[l] \ar[d] \\
\textit{Alg}_A \ar[r] & \textit{Ensembles} \ar@<1ex>[l]
}
$$
est commutatif. Il en résulte que $X_\bullet(\text{id}_B : B \to B)$ est égal
à la résolution standard de $B$ sur $A$.
\end{remark}

\begin{lemma}
\label{lemma-colimit-cotangent-complex}
Soit $A_i \to B_i$ un système d'homomorphismes d'anneaux indexé par un
ensemble ordonné filtrant $I$. Alors
$\colim L_{B_i/A_i} = L_{\colim B_i/\colim A_i}$.
\end{lemma}

\begin{proof}
Cela est vrai parce que le foncteur d'oubli
$V : A\textit{-Alg} \to \textit{Ensembles}$ et son adjoint
$U : \textit{Ensembles} \to A\textit{-Alg}$ commutent aux colimites filtrantes.
Il en va de même du foncteur $B/A \mapsto \Omega_{B/A}$
(Algèbre, lemme \ref{algebra-lemma-colimit-differentials}).
\end{proof}





\section{Résolutions simpliciales et foncteur dérivé gauche}
\label{section-compute-L-pi-shriek}

\noindent
Soit $A \to B$ un homomorphisme d'anneaux. Considérons la catégorie dont les
objets sont les homomorphismes de $A$-algèbres $\alpha : P \to B$, où $P$ est
une algèbre polynomiale sur $A$ (en un certain ensemble de
variables\footnote{Il suffit de considérer des ensembles dont le cardinal est
au plus celui de $B$.}), et dont les morphismes
$s : (\alpha : P \to B) \to (\alpha' : P' \to B)$ sont les homomorphismes de
$A$-algèbres $s : P \to P'$ tels que $\alpha' \circ s = \alpha$.
Notons $\mathcal{C} = \mathcal{C}_{B/A}$ la catégorie {\bf opposée} de cette
catégorie. Ce passage à l'opposée vient de ce que nous voulons considérer les
objets $(P, \alpha)$ comme correspondant au diagramme de schémas affines
$$
\xymatrix{
\Spec(B) \ar[d] \ar[r] & \Spec(P) \ar[ld] \\
\Spec(A)
}
$$
Munissons $\mathcal{C}$ de la topologie chaotique
(Sites, exemple \ref{sites-example-indiscrete}), c'est-à-dire munissons
$\mathcal{C}$ de la structure de site dans laquelle les recouvrements sont
donnés par les identités, de
sorte que tout préfaisceau est un faisceau. Munissons en outre $\mathcal{C}$
de deux faisceaux d'anneaux. Le premier est le faisceau $\mathcal{O}$ qui
associe $P$ à l'objet $(P, \alpha)$. Le second est le faisceau constant $B$,
que nous noterons $\underline{B}$. On obtient le diagramme suivant de
morphismes de topos annelés
\begin{equation}
\label{equation-pi}
\vcenter{
\xymatrix{
(\Sh(\mathcal{C}), \underline{B}) \ar[r]_i \ar[d]_\pi &
(\Sh(\mathcal{C}), \mathcal{O}) \\
(\Sh(*), B)
}
}
\end{equation}
Le morphisme $i$ est l'identité sur les topos sous-jacents, et
$i^\sharp : \mathcal{O} \to \underline{B}$ est l'homomorphisme évident.
L'application $\pi$ est celle de Cohomologie sur les sites, exemple
\ref{sites-cohomology-example-category-to-point}. Les foncteurs dérivés
suivants joueront un rôle important :
$
Li^* : D(\mathcal{O}) \longrightarrow D(\underline{B})
$
adjoint à gauche de $Ri_* = i_* : D(\underline{B}) \to D(\mathcal{O})$, et
$
L\pi_! : D(\underline{B}) \longrightarrow D(B)
$
adjoint à gauche de $\pi^* = \pi^{-1} : D(B) \to D(\underline{B})$.

\begin{lemma}
\label{lemma-identify-pi-shriek}
Avec les notations précédentes, soit $P_\bullet$ une $A$-algèbre simpliciale
munie d'une augmentation $\epsilon : P_\bullet \to B$.
Supposons que chaque $P_n$ soit une algèbre polynomiale sur $A$ et que
$\epsilon$ soit une fibration de Kan triviale sur les ensembles simpliciaux
sous-jacents. Alors
$$
L\pi_!(\mathcal{F}) = \mathcal{F}(P_\bullet, \epsilon)
$$
dans $D(\textit{Ab})$, resp.\ $D(B)$, fonctoriellement en $\mathcal{F}$ dans
$\textit{Ab}(\mathcal{C})$, resp.\ $\textit{Mod}(\underline{B})$.
\end{lemma}

\begin{proof}
Nous utiliserons pour cela le critère de Cohomologie sur les sites, lemme
\ref{sites-cohomology-lemma-compute-by-cosimplicial-resolution}.
Étant donné un objet $U = (Q, \beta)$ de $\mathcal{C}$, il faut montrer que
$$
S_\bullet = \Mor_\mathcal{C}((P_\bullet, \epsilon), (Q, \beta))
$$
est homotopiquement équivalent à un singleton.
Écrivons $Q = A[E]$ pour un certain ensemble $E$ (ce qui est possible par
notre choix de la catégorie $\mathcal{C}$). On voit que
$$
S_\bullet = \Mor_{\textit{Ensembles}}((E, \beta|_E), (P_\bullet, \epsilon))
$$
Soit $*$ l'ensemble simplicial constant de valeur un singleton. Pour $b \in B$,
soit $F_{b, \bullet}$ l'ensemble simplicial défini par le diagramme cartésien
$$
\xymatrix{
F_{b, \bullet} \ar[r] \ar[d] & P_\bullet \ar[d]_\epsilon \\
{*} \ar[r]^b & B
}
$$
Avec cette notation, $S_\bullet = \prod_{e \in E} F_{\beta(e), \bullet}$.
Comme $\epsilon$ est par hypothèse une fibration de Kan triviale,
$F_{b, \bullet} \to *$ est une fibration de Kan triviale
(Objets simpliciaux, lemme
\ref{simplicial-lemma-trivial-kan-base-change}). Ainsi,
$S_\bullet \to *$ est une fibration de Kan triviale (Objets simpliciaux,
lemme \ref{simplicial-lemma-product-trivial-kan}). Par conséquent,
$S_\bullet$ est homotopiquement équivalent à $*$ (Objets simpliciaux, lemme
\ref{simplicial-lemma-trivial-kan-homotopy}).
\end{proof}

\noindent
En particulier, on peut employer la résolution standard de $B$ sur $A$ pour
calculer le foncteur dérivé gauche correspondant.

\begin{lemma}
\label{lemma-pi-shriek-standard}
Soit $A \to B$ un homomorphisme d'anneaux. Soit
$\epsilon : P_\bullet \to B$ la résolution standard de $B$ sur $A$.
Soit $\pi$ comme dans (\ref{equation-pi}). Alors
$$
L\pi_!(\mathcal{F}) = \mathcal{F}(P_\bullet, \epsilon)
$$
dans $D(\textit{Ab})$, resp.\ $D(B)$, fonctoriellement en $\mathcal{F}$ dans
$\textit{Ab}(\mathcal{C})$, resp.\ $\textit{Mod}(\underline{B})$.
\end{lemma}

\begin{proof}[Première démonstration]
Appliquons le lemme \ref{lemma-identify-pi-shriek}.
Puisque les termes $P_n$ sont des algèbres polynomiales, la première hypothèse
de ce lemme est satisfaite. Démontrons la seconde. D'après Objets simpliciaux,
lemme \ref{simplicial-lemma-standard-simplicial-homotopy}, l'application
$\epsilon$ est une équivalence d'homotopie des ensembles simpliciaux
sous-jacents. D'après Objets simpliciaux, lemme
\ref{simplicial-lemma-homotopy-equivalence}, il en résulte que $\epsilon$
induit un quasi-isomorphisme des complexes de groupes abéliens associés.
D'après Objets simpliciaux, lemme
\ref{simplicial-lemma-qis-simplicial-abelian-groups}, cela implique que
$\epsilon$ est une fibration de Kan triviale des ensembles simpliciaux
sous-jacents.
\end{proof}

\begin{proof}[Seconde démonstration]
Nous utiliserons le critère de Cohomologie sur les sites, lemme
\ref{sites-cohomology-lemma-compute-by-cosimplicial-resolution}.
Soit $U = (Q, \beta)$ un objet de $\mathcal{C}$.
Il nous faut montrer que
$$
S_\bullet = \Mor_\mathcal{C}((P_\bullet, \epsilon), (Q, \beta))
$$
est homotopiquement équivalent à un singleton. Écrivons $Q = A[E]$ pour un
certain ensemble $E$ (ce qui est possible par notre choix de la catégorie
$\mathcal{C}$). Avec les notations de la remarque
\ref{remark-variant-cotangent-complex}, on voit que
$$
S_\bullet = \Mor_\mathcal{S}((E \to B), i(P_\bullet \to B))
$$
D'après Objets simpliciaux, lemme
\ref{simplicial-lemma-standard-simplicial-homotopy}, l'application
$i(P_\bullet \to B) \to i(B \to B)$ est une équivalence d'homotopie dans
$\mathcal{S}$. Par conséquent, $S_\bullet$ est homotopiquement équivalent à
$$
\Mor_\mathcal{S}((E \to B), (B \to B)) = \{*\}
$$
comme voulu.
\end{proof}

\begin{lemma}
\label{lemma-compute-cotangent-complex}
Soit $A \to B$ un homomorphisme d'anneaux. Soient $\pi$ et $i$ comme dans
(\ref{equation-pi}). Il existe un isomorphisme canonique
$$
L_{B/A} = L\pi_!(Li^*\Omega_{\mathcal{O}/A}) =
L\pi_!(i^*\Omega_{\mathcal{O}/A}) =
L\pi_!(\Omega_{\mathcal{O}/A} \otimes_\mathcal{O} \underline{B})
$$
dans $D(B)$.
\end{lemma}

\begin{proof}
Pour un objet $\alpha : P \to B$ de la catégorie $\mathcal{C}$, le module
$\Omega_{P/A}$ est un $P$-module libre. Ainsi,
$\Omega_{\mathcal{O}/A}$ est un $\mathcal{O}$-module plat. Par conséquent,
$Li^*\Omega_{\mathcal{O}/A} = i^*\Omega_{\mathcal{O}/A}$ est le faisceau de
$\underline{B}$-modules qui associe à $\alpha : P \to A$ le $B$-module
$\Omega_{P/A} \otimes_{P, \alpha} B$. D'après le lemme
\ref{lemma-pi-shriek-standard}, le membre de droite se calcule en évaluant ce
faisceau sur la résolution standard ; c'est précisément notre définition du
membre de gauche (définition \ref{definition-cotangent-complex-ring-map}).
\end{proof}

\begin{lemma}
\label{lemma-pi-lower-shriek-constant-sheaf}
Si $A \to B$ est un homomorphisme d'anneaux, alors
$L\pi_!(\pi^{-1}M) = M$, où $\pi$ est comme dans (\ref{equation-pi}).
\end{lemma}

\begin{proof}
Cela résulte du lemme \ref{lemma-identify-pi-shriek}, qui affirme que
$L\pi_!(\pi^{-1}M)$ se calcule par $(\pi^{-1}M)(P_\bullet, \epsilon)$, objet
simplicial constant de valeur $M$.
\end{proof}

\begin{lemma}
\label{lemma-identify-H0}
Si $A \to B$ est un homomorphisme d'anneaux, alors
$H^0(L_{B/A}) = \Omega_{B/A}$.
\end{lemma}

\begin{proof}
Nous allons le démontrer par un calcul direct, en utilisant l'identification
du lemme \ref{lemma-compute-cotangent-complex}. Il existe manifestement un
homomorphisme de $\Omega_{\mathcal{O}/A} \otimes \underline{B}$ vers le
faisceau constant de valeur $\Omega_{B/A}$. Il induit donc un homomorphisme
$$
H^0(L_{B/A}) = H^0(L\pi_!(\Omega_{\mathcal{O}/A} \otimes \underline{B}))
= \pi_!(\Omega_{\mathcal{O}/A} \otimes \underline{B}) \to \Omega_{B/A}
$$
En choisissant un objet $P \to B$ de $\mathcal{C}_{B/A}$ tel que $P \to B$
soit surjectif, on voit que cet homomorphisme est surjectif (d'après Algèbre,
lemme \ref{algebra-lemma-differential-surjective}). Pour montrer son
injectivité, supposons que $P \to B$ soit un objet de $\mathcal{C}_{B/A}$ et
que $\xi \in \Omega_{P/A} \otimes_P B$ soit un élément dont l'image dans
$\Omega_{B/A}$ est nulle. Choisissons d'abord une factorisation
$P \to P' \to B$ telle que $P' \to B$ soit surjectif et que $P'$ soit une
algèbre polynomiale sur $A$. Nous pouvons remplacer $P$ par $P'$. Si
$B = P/I$, alors le noyau de
$\Omega_{P/A} \otimes_P B \to \Omega_{B/A}$ est l'image de $I/I^2$ (Algèbre,
lemme \ref{algebra-lemma-differential-seq}). Disons que $\xi$ est l'image de
$f \in I$. Considérons alors les deux homomorphismes
$a, b : P' = P[x] \to P$ : le premier envoie $x$ sur $0$ et le second envoie
$x$ sur $f$ (dans les deux cas, $P[x] \to B$ envoie $x$ sur zéro). On voit que
$\xi$ et $0$ sont les images de $\text{d}x \otimes 1$ dans
$\Omega_{P'/A} \otimes_{P'} B$. Ainsi, $\xi$ et $0$ ont la même image dans la
colimite (voir Cohomologie sur les sites, exemple
\ref{sites-cohomology-example-category-to-point})
$\pi_!(\Omega_{\mathcal{O}/A} \otimes \underline{B})$, comme voulu.
\end{proof}

\begin{lemma}
\label{lemma-pi-lower-shriek-polynomial-algebra}
Si $B$ est une algèbre polynomiale sur l'anneau $A$, alors, pour $\pi$ comme
dans (\ref{equation-pi}), le foncteur $\pi_!$ est exact et
$\pi_!\mathcal{F} = \mathcal{F}(B \to B)$.
\end{lemma}

\begin{proof}
Cela résulte du lemme \ref{lemma-identify-pi-shriek}, qui affirme que
l'algèbre simpliciale constante de valeur $B$ peut servir à calculer $L\pi_!$.
\end{proof}

\begin{lemma}
\label{lemma-cotangent-complex-polynomial-algebra}
Si $B$ est une algèbre polynomiale sur l'anneau $A$, alors $L_{B/A}$ est
quasi-isomorphe à $\Omega_{B/A}[0]$.
\end{lemma}

\begin{proof}
Cela résulte immédiatement des lemmes
\ref{lemma-compute-cotangent-complex} et
\ref{lemma-pi-lower-shriek-polynomial-algebra}.
\end{proof}





\section{Construction d'une résolution}
\label{section-polynomial}

\noindent
Dans le cas noethérien de type fini, on peut construire une ``petite''
résolution simpliciale des homomorphismes d'anneaux de type fini.

\begin{lemma}
\label{lemma-polynomial}
Soit $A$ un anneau noethérien. Soit $A \to B$ un homomorphisme d'anneaux de
type fini. Soit $\mathcal{A}$ la catégorie des homomorphismes de $A$-algèbres
$C \to B$. Soit $n \geq 0$, et soit $P_\bullet$ un objet simplicial de
$\mathcal{A}$ tel que :
\begin{enumerate}
\item $P_\bullet \to B$ est une fibration de Kan triviale d'ensembles
simpliciaux ;
\item $P_k$ est de type fini sur $A$ pour $k \leq n$ ;
\item $P_\bullet = \text{cosk}_n \text{sk}_n P_\bullet$ comme objets
simpliciaux de $\mathcal{A}$.
\end{enumerate}
Alors $P_{n + 1}$ est une $A$-algèbre de type fini.
\end{lemma}

\begin{proof}
Bien que la démonstration que nous allons donner soit directe, elle est un peu
laborieuse. Pour en éclairer l'idée, expliquons d'abord ce qui se passe pour
les petites valeurs de $n$, avant de traiter le cas général. Si, par exemple,
$n = 0$, alors (3) signifie que $P_1 = P_0 \times_B P_0$. Comme
l'homomorphisme d'anneaux $P_0 \to B$ est surjectif, cette algèbre est de type
fini sur $A$ d'après Compléments d'algèbre, lemme
\ref{more-algebra-lemma-fibre-product-finite-type}.

\medskip\noindent
Si $n = 1$, alors (3) signifie que
$$
P_2 = \{(f_0, f_1, f_2) \in P_1^3 \mid
d_0f_0 = d_0f_1,\ d_1f_0 = d_0f_2,\ d_1f_1 = d_1f_2 \}
$$
où les égalités ont lieu dans $P_0$. Remarquons que le triplet
$$
(d_0f_0, d_1f_0, d_1f_1) = (d_0f_1, d_0f_2, d_1f_2)
$$
est un élément du produit fibré $P_0 \times_B P_0 \times_B P_0$ au-dessus de
$B$, car les applications $d_i : P_1 \to P_0$ sont des morphismes au-dessus
de $B$. On obtient donc une application
$$
\psi : P_2 \longrightarrow P_0 \times_B P_0 \times_B P_0
$$
La fibre de $\psi$ au-dessus d'un élément
$(g_0, g_1, g_2) \in P_0 \times_B P_0 \times_B P_0$
est l'ensemble des triplets $(f_0, f_1, f_2)$ de $1$-simplexes tels que
$(d_0, d_1)(f_0) = (g_0, g_1)$, $(d_0, d_1)(f_1) = (g_0, g_2)$ et
$(d_0, d_1)(f_2) = (g_1, g_2)$. Puisque $P_\bullet \to B$ est une fibration
de Kan triviale, l'application $(d_0, d_1) : P_1 \to P_0 \times_B P_0$ est
surjective. Ainsi, $P_2$ s'insère dans le diagramme cartésien
$$
\xymatrix{
P_2 \ar[d] \ar[r] & P_1^3 \ar[d] \\
P_0 \times_B P_0 \times_B P_0 \ar[r] & (P_0 \times_B P_0)^3
}
$$
On conclut par Compléments d'algèbre, lemme
\ref{more-algebra-lemma-formal-consequence}. Le cas général est analogue, mais
demande un peu plus de notation.

\medskip\noindent
Traitons le cas $n > 1$. D'après Objets simpliciaux, lemme
\ref{simplicial-lemma-cosk-above-object}, la condition
$P_\bullet = \text{cosk}_n \text{sk}_n P_\bullet$ implique que la même égalité
vaut dans la catégorie des $A$-algèbres simpliciales, puis dans celle des
ensembles (car le foncteur d'oubli des $A$-algèbres vers les ensembles commute
aux limites). Ainsi,
$$
P_{n + 1} =
\Mor(\Delta[n + 1], P_\bullet) =
\Mor(\text{sk}_n \Delta[n + 1], \text{sk}_n P_\bullet)
$$
d'après Objets simpliciaux, lemme \ref{simplicial-lemma-simplex-map}, et
l'équation (\ref{simplicial-equation-cosk}). Nous allons démontrer par
récurrence sur $1 \leq k < m \leq n + 1$ que l'anneau
$$
Q_{k, m} = \Mor(\text{sk}_k \Delta[m], \text{sk}_k P_\bullet)
$$
est de type fini sur $A$. Le cas $k = 1$, $1 < m \leq n + 1$, est entièrement
analogue à la discussion ci-dessus du cas $n = 1$. Plus précisément, on a un
diagramme cartésien
$$
\xymatrix{
Q_{1, m} \ar[d] \ar[r] & P_1^N \ar[d] \\
P_0 \times_B \ldots \times_B P_0 \ar[r] & (P_0 \times_B P_0)^N
}
$$
où $N = {m + 1 \choose 2}$. On conclut comme précédemment.

\medskip\noindent
Soit $1 \leq k_0 \leq n$, et supposons que $Q_{k, m}$ soit de type fini sur
$A$ pour tous $1 \leq k \leq k_0$ et $k < m \leq n + 1$.
Pour $k_0 + 1 < m \leq n + 1$, nous affirmons qu'il existe un carré cartésien
$$
\xymatrix{
Q_{k_0 + 1, m} \ar[d] \ar[r] & P_{k_0 + 1}^N \ar[d] \\
Q_{k_0, m} \ar[r] & Q_{k_0, k_0 + 1}^N
}
$$
où $N$ est le nombre de $(k_0 + 1)$-simplexes non dégénérés de $\Delta[m]$.
En effet, considérons un élément de $Q_{k_0 + 1, m}$ comme une application $f$
du $(k_0 + 1)$-squelette de $\Delta[m]$ vers $P_\bullet$. On peut restreindre
$f$ au $k_0$-squelette, ce qui donne l'application verticale gauche du
diagramme. Sa
restriction à chaque $(k_0 + 1)$-simplexe non dégénéré donne la flèche
horizontale supérieure. De plus, se donner un tel $f$ revient à se donner sa
restriction au $k_0$-squelette et à chaque $(k_0 + 1)$-face non dégénérée,
pourvu que ces restrictions coïncident sur leurs intersections ; c'est
exactement le contenu du diagramme. Enfin, le fait que
$P_\bullet \to B$ soit une fibration de Kan triviale implique que
l'application
$$
P_{k_0} \to Q_{k_0, k_0 + 1} = \Mor(\partial \Delta[k_0 + 1], P_\bullet)
$$
est surjective, car toute application $\partial \Delta[k_0 + 1] \to B$ se
prolonge en une application $\Delta[k_0 + 1] \to B$ pour $k_0 \geq 1$ (nous
omettons un petit argument sur les ensembles simpliciaux constants). Par
l'hypothèse de récurrence, les anneaux $Q_{k_0, m}$,
$Q_{k_0, k_0 + 1}$ sont des $A$-algèbres de type fini ; il en va donc de même
de $Q_{k_0 + 1, m}$, encore une fois d'après Compléments d'algèbre, lemme
\ref{more-algebra-lemma-formal-consequence}.
\end{proof}

\begin{proposition}
\label{proposition-polynomial}
Soit $A$ un anneau noethérien. Soit $A \to B$ un homomorphisme d'anneaux de
type fini. Il existe une $A$-algèbre simpliciale $P_\bullet$ munie d'une
augmentation $\epsilon : P_\bullet \to B$ telle que chaque $P_n$ soit une
algèbre polynomiale de type fini sur $A$ et que $\epsilon$ soit une fibration
de Kan triviale d'ensembles simpliciaux.
\end{proposition}

\begin{proof}
Soit $\mathcal{A}$ la catégorie des homomorphismes de $A$-algèbres $C \to B$.
Dans cette démonstration, les objets simpliciaux ainsi que les foncteurs
squelette et cosquelette seront pris dans cette catégorie.

\medskip\noindent
Choisissons une algèbre polynomiale $P_0$ de type fini sur $A$ et une
surjection $P_0 \to B$. En première approximation, prenons
$P_\bullet = \text{cosk}_0(P_0)$. Autrement dit, $P_\bullet$ est la
$A$-algèbre simpliciale dont les termes sont
$P_n = P_0 \times_A \ldots \times_A P_0$. (Dans le dernier paragraphe de la
démonstration, cet objet simplicial sera noté $P^0_\bullet$.) D'après Objets
simpliciaux, lemme \ref{simplicial-lemma-cosk-minus-one-equivalence},
l'application $P_\bullet \to B$ est une fibration de Kan triviale d'ensembles
simpliciaux. Remarquons aussi que
$P_\bullet = \text{cosk}_0 \text{sk}_0 P_\bullet$.

\medskip\noindent
Supposons que, pour un certain $n \geq 0$, nous ayons construit $P_\bullet$
(qui sera noté $P^n_\bullet$ dans le dernier paragraphe de la démonstration)
de telle sorte que :
\begin{enumerate}
\item[(a)] $P_\bullet \to B$ est une fibration de Kan triviale d'ensembles
simpliciaux ;
\item[(b)] $P_k$ est une algèbre polynomiale de type fini pour
$0 \leq k \leq n$ ;
\item[(c)] $P_\bullet = \text{cosk}_n \text{sk}_n P_\bullet$
\end{enumerate}
D'après le lemme \ref{lemma-polynomial}, il existe une algèbre polynomiale de
type fini $Q$ sur $A$ et une surjection $Q \to P_{n + 1}$. Comme $P_n$ est une
algèbre polynomiale, les homomorphismes de $A$-algèbres
$s_i : P_n \to P_{n + 1}$ se relèvent en des homomorphismes
$s'_i : P_n \to Q$. Définissons $d'_j : Q \to P_n$ comme la composée de
$Q \to P_{n + 1}$ et de $d_j : P_{n + 1} \to P_n$. On obtient un objet
simplicial tronqué $P'_\bullet$ de $\mathcal{A}$ en posant $P'_k = P_k$ pour
$k \leq n$ et $P'_{n + 1} = Q$, en prenant les morphismes $d'_i = d_i$ et
$s'_i = s_i$ aux degrés $k \leq n - 1$, puis les morphismes $d'_j$ et $s'_i$
au degré $n$. Prolongeons-le en un objet simplicial complet $P'_\bullet$ de
$\mathcal{A}$ au moyen de $\text{cosk}_{n + 1}$. Par fonctorialité des
foncteurs cosquelette, il existe un morphisme d'objets simpliciaux
$P'_\bullet \to P_\bullet$ qui prolonge le morphisme donné d'objets simpliciaux
$(n + 1)$-tronqués. (Ce morphisme sera noté
$P^{n + 1}_\bullet \to P^n_\bullet$ dans le dernier paragraphe de la
démonstration.)

\medskip\noindent
Notons que les conditions (b) et (c) sont satisfaites par $P'_\bullet$ après
remplacement de $n$ par $n + 1$. Nous affirmons que l'application
$P'_\bullet \to P_\bullet$ satisfait aux hypothèses (1), (2), (3) et (4) de
Objets simpliciaux, lemme \ref{simplicial-lemma-section}, avec $n + 1$ à la
place de $n$. Les conditions (1) et (2) sont vraies par construction. D'après
Objets simpliciaux, lemme \ref{simplicial-lemma-cosk-above-object}, on a
$P_\bullet = \text{cosk}_{n + 1}\text{sk}_{n + 1}P_\bullet$
et
$P'_\bullet = \text{cosk}_{n + 1}\text{sk}_{n + 1}P'_\bullet$
non seulement dans $\mathcal{A}$, mais aussi dans la catégorie des
$A$-algèbres, puis dans celle des ensembles (car le foncteur d'oubli des
$A$-algèbres vers les ensembles commute à toutes les limites). Cela démontre
(3) et (4). Le lemme s'applique donc, et
$P'_\bullet \to P_\bullet$ est une fibration de Kan triviale. D'après Objets
simpliciaux, lemme \ref{simplicial-lemma-trivial-kan-composition}, on en déduit
que $P'_\bullet \to B$ est une fibration de Kan triviale ; la condition (a)
est donc également satisfaite.

\medskip\noindent
Pour achever la démonstration, prenons la limite projective
$P_\bullet = \lim P^n_\bullet$ de la suite d'algèbres simpliciales
$$
\ldots \to P^2_\bullet \to P^1_\bullet \to P^0_\bullet
$$
construite ci-dessus. L'application $P_\bullet \to B$ est une fibration de Kan
triviale d'après Objets simpliciaux, lemme
\ref{simplicial-lemma-limit-trivial-kan}. Or la construction ci-dessus se
stabilise, en chaque degré, en une algèbre polynomiale de type fini fixée,
comme voulu.
\end{proof}

\begin{lemma}
\label{lemma-pi-shriek-finite}
Soit $A$ un anneau noethérien. Soit $A \to B$ un homomorphisme d'anneaux de
type fini. Soient $\pi$, $\underline{B}$ comme dans (\ref{equation-pi}).
Si $\mathcal{F}$ est un $\underline{B}$-module tel que
$\mathcal{F}(P, \alpha)$ soit un $B$-module fini pour tout
$\alpha : P = A[x_1, \ldots, x_n] \to B$, alors les modules de cohomologie de
$L\pi_!(\mathcal{F})$ sont des $B$-modules finis.
\end{lemma}

\begin{proof}
D'après le lemme \ref{lemma-identify-pi-shriek} et la proposition
\ref{proposition-polynomial}, on peut calculer $L\pi_!(\mathcal{F})$ au moyen
d'un complexe construit à partir des valeurs de $\mathcal{F}$ sur les algèbres
polynomiales de type fini.
\end{proof}

\begin{lemma}
\label{lemma-cotangent-finite}
Soit $A$ un anneau noethérien. Soit $A \to B$ un homomorphisme d'anneaux de
type fini. Alors $H^n(L_{B/A})$ est un $B$-module fini pour tout
$n \in \mathbf{Z}$.
\end{lemma}

\begin{proof}
Appliquer les lemmes \ref{lemma-compute-cotangent-complex} et
\ref{lemma-pi-shriek-finite}.
\end{proof}

\begin{remark}[Resolutions]
\label{remark-resolution}
Soit $A \to B$ un homomorphisme d'anneaux quelconque. Appelons une
$A$-algèbre simpliciale augmentée $\epsilon : P_\bullet \to B$ une
{\it résolution de $B$ sur $A$} si chaque $P_n$ est une algèbre polynomiale
et si $\epsilon$ est une fibration de Kan triviale d'ensembles simpliciaux.
Si $P_\bullet \to B$ est une augmentation d'une $A$-algèbre simpliciale telle
que chaque $P_n$ soit une algèbre polynomiale se surjectant sur $B$, alors les
conditions suivantes sont équivalentes :
\begin{enumerate}
\item $\epsilon : P_\bullet \to B$ est une résolution de $B$ sur $A$ ;
\item $\epsilon : P_\bullet \to B$ est un quasi-isomorphisme sur les
complexes associés ;
\item $\epsilon : P_\bullet \to B$ induit une équivalence d'homotopie
d'ensembles simpliciaux.
\end{enumerate}
Pour le voir, utiliser Objets simpliciaux, lemmes
\ref{simplicial-lemma-trivial-kan-homotopy},
\ref{simplicial-lemma-homotopy-equivalence} et
\ref{simplicial-lemma-qis-simplicial-abelian-groups}.
Une résolution $P_\bullet$ de $B$ sur $A$ fournit un objet cosimplicial
$U_\bullet$ de $\mathcal{C}_{B/A}$ comme dans Cohomologie sur les sites, lemme
\ref{sites-cohomology-lemma-compute-by-cosimplicial-resolution}
et il s'ensuit que
$$
L\pi_!\mathcal{F} = \mathcal{F}(P_\bullet)
$$
fonctoriellement en $\mathcal{F}$ ; voir le lemme
\ref{lemma-identify-pi-shriek}. La partie formelle de la démonstration de la
proposition \ref{proposition-polynomial} montre l'existence de résolutions.
Nous avons aussi vu, dans la première démonstration du lemme
\ref{lemma-pi-shriek-standard}, que la résolution standard de $B$ sur $A$ est
une résolution (de sorte que cette terminologie ne crée pas de conflit).
Cependant, l'argument de la démonstration de la proposition
\ref{proposition-polynomial} établit l'existence de résolutions sans faire
appel aux calculs simpliciaux d'Objets simpliciaux, section
\ref{simplicial-section-standard}. De plus, pour {\it tout} choix de
résolution, on dispose d'un isomorphisme canonique
$$
L_{B/A} = \Omega_{P_\bullet/A} \otimes_{P_\bullet, \epsilon} B
$$
dans $D(B)$ d'après le lemme \ref{lemma-compute-cotangent-complex}. La liberté
de choisir une résolution arbitraire peut être fort utile.
\end{remark}

\begin{lemma}
\label{lemma-O-homology-B-homology}
Soit $A \to B$ un homomorphisme d'anneaux. Soient $\pi$, $\mathcal{O}$,
$\underline{B}$ comme dans (\ref{equation-pi}). Pour tout $\mathcal{O}$-module
$\mathcal{F}$, on a
$$
L\pi_!(\mathcal{F}) = L\pi_!(Li^*\mathcal{F}) =
L\pi_!(\mathcal{F} \otimes_\mathcal{O}^\mathbf{L} \underline{B})
$$
dans $D(\textit{Ab})$.
\end{lemma}

\begin{proof}
Il suffit de vérifier que les hypothèses de Cohomologie sur les sites, lemme
\ref{sites-cohomology-lemma-O-homology-qis}
sont satisfaites pour $\mathcal{O} \to \underline{B}$ sur
$\mathcal{C}_{B/A}$. Nous utiliserons sans plus le signaler les résultats de
la remarque \ref{remark-resolution}. Choisissons une résolution $P_\bullet$
de $B$ sur $A$ afin d'obtenir un objet cosimplicial convenable $U_\bullet$ de
$\mathcal{C}_{B/A}$. Puisque $P_\bullet \to B$ induit un quasi-isomorphisme
sur les complexes associés de groupes abéliens, on voit que
$L\pi_!\mathcal{O} = B$. D'autre part, $L\pi_!\underline{B}$ est calculé par
$\underline{B}(U_\bullet) = B$. Cela vérifie la seconde hypothèse de
Cohomologie sur les sites, lemme
\ref{sites-cohomology-lemma-O-homology-qis}
et achève la démonstration.
\end{proof}

\begin{lemma}
\label{lemma-apply-O-B-comparison}
Soit $A \to B$ un homomorphisme d'anneaux. Soient $\pi$, $\mathcal{O}$,
$\underline{B}$ comme dans (\ref{equation-pi}). On a
$$
L\pi_!(\mathcal{O}) = L\pi_!(\underline{B}) = B
\quad\text{et}\quad
L_{B/A} = L\pi_!(\Omega_{\mathcal{O}/A} \otimes_\mathcal{O} \underline{B}) =
L\pi_!(\Omega_{\mathcal{O}/A})
$$
dans $D(\textit{Ab})$.
\end{lemma}

\begin{proof}
C'est une application immédiate du lemme
\ref{lemma-O-homology-B-homology} (et la première égalité à droite est le
lemme \ref{lemma-compute-cotangent-complex}).
\end{proof}

\noindent
Voici un cas particulier du triangle fondamental dont la démonstration est
facile.

\begin{lemma}
\label{lemma-special-case-triangle}
Soient $A \to B \to C$ des homomorphismes d'anneaux. Si $B$ est une algèbre
polynomiale sur $A$, alors il existe un triangle distingué 
$L_{B/A} \otimes_B^\mathbf{L} C \to L_{C/A} \to L_{C/B} \to
L_{B/A} \otimes_B^\mathbf{L} C[1]$ dans $D(C)$.
\end{lemma}

\begin{proof}
Nous utiliserons sans plus le signaler les observations de la remarque
\ref{remark-resolution}. Choisissons une résolution
$\epsilon : P_\bullet \to C$ de $C$ sur $B$ (par exemple la résolution
standard). Puisque $B$ est une algèbre polynomiale sur $A$, on voit que
$P_\bullet$ est aussi une résolution de $C$ sur $A$. Ainsi, $L_{C/A}$ est
calculé par
$\Omega_{P_\bullet/A} \otimes_{P_\bullet, \epsilon} C$
et $L_{C/B}$ est calculé par
$\Omega_{P_\bullet/B} \otimes_{P_\bullet, \epsilon} C$.
Comme, pour chaque $n$, on a la suite exacte courte
$0 \to \Omega_{B/A} \otimes_B P_n \to \Omega_{P_n/A} \to \Omega_{P_n/B} \to 0$
(Algèbre, lemme \ref{algebra-lemma-ses-formally-smooth}) et comme
$L_{B/A} = \Omega_{B/A}[0]$ (lemme
\ref{lemma-cotangent-complex-polynomial-algebra}), on obtient le résultat.
\end{proof}

\begin{example}
\label{example-resolution-length-two}
Soit $A \to B$ un homomorphisme d'anneaux. Dans cet exemple, nous allons 
construire une résolution ``explicite'' $P_\bullet$ de $B$ sur $A$ de longueur
$2$. Pour ce faire, nous suivons la procédure de la démonstration de la
proposition \ref{proposition-polynomial} ; voir aussi la discussion de la
remarque \ref{remark-resolution}.

\medskip\noindent
Choisissons une surjection $P_0 = A[u_i] \to B$, où les $u_i$ forment un
ensemble de variables. Choisissons des générateurs $f_t \in P_0$, $t \in T$,
de l'idéal $\Ker(P_0 \to B)$. Prenons $P_1 = A[u_i, x_t]$, avec pour
applications de face $d_0$ et $d_1$ les uniques homomorphismes de $A$-algèbres
tels que $d_j(u_i) = u_i$, $d_0(x_t) = 0$ et $d_1(x_t) = f_t$.
L'application $s_0 : P_0 \to P_1$ est l'unique homomorphisme de $A$-algèbres
tel que $s_0(u_i) = u_i$. Il est clair que
$$
P_1 \xrightarrow{d_0 - d_1} P_0 \to B \to 0
$$
est exacte ; en particulier, l'application
$(d_0, d_1) : P_1 \to P_0 \times_B P_0$ est surjective. Ainsi, si
$P_\bullet$ désigne l'algèbre simpliciale $1$-tronquée sur $A$ donnée par $P_0$,
$P_1$, $d_0$, $d_1$ et $s_0$, alors l'augmentation
$\text{cosk}_1(P_\bullet) \to B$ est une fibration de Kan triviale. L'étape
suivante de la procédure de la démonstration de la proposition
\ref{proposition-polynomial} consiste à choisir une algèbre polynomiale $P_2$
et une surjection
$$
P_2 \longrightarrow \text{cosk}_1(P_\bullet)_2
$$
Rappelons que
$$
\text{cosk}_1(P_\bullet)_2 =
\{(g_0, g_1, g_2) \in P_1^3 \mid d_0(g_0) = d_0(g_1),
d_1(g_0) = d_0(g_2), d_1(g_1) = d_1(g_2)\}
$$
En considérant $g_i \in P_1$ comme un polynôme en les $x_t$, les conditions
sont
$$
g_0(0) = g_1(0),\quad
g_0(f_t) = g_2(0),\quad
g_1(f_t) = g_2(f_t)
$$
Ainsi, $\text{cosk}_1(P_\bullet)_2$ contient les éléments
$y_t = (x_t, x_t, f_t)$ et $z_t = (0, x_t, x_t)$. Tout élément $G$ de
$\text{cosk}_1(P_\bullet)_2$ est de la forme $G = H + (0, 0, g)$, où $H$
appartient à l'image de
$A[u_i, y_t, z_t] \to \text{cosk}_1(P_\bullet)_2$. Ici, $g \in P_1$ est un
polynôme de terme constant nul tel que $g(f_t) = 0$ dans $P_0$. Observons que
\begin{enumerate}
\item $g = x_t x_{t'} - f_t x_{t'}$ ;
\item $g = \sum r_t x_t$ avec $r_t \in P_0$ si $\sum r_t f_t = 0$ dans $P_0$,
\end{enumerate}
sont des éléments de $P_1$ de la forme voulue. Posons
$$
Rel = \Ker(\bigoplus\nolimits_{t \in T} P_0 \longrightarrow P_0),\quad
(r_t) \longmapsto \sum r_tf_t
$$
Prenons $P_2 = A[u_i, y_t, z_t, v_r, w_{t, t'}]$, où
$r = (r_t) \in Rel$, muni de l'application
$$
P_2 \longrightarrow \text{cosk}_1(P_\bullet)_2
$$
définie par $y_t  \mapsto (x_t, x_t, f_t)$,
$z_t \mapsto (0, x_t, x_t)$,
$v_r \mapsto (0, 0, \sum r_t x_t)$ et
$w_{t, t'} \mapsto (0, 0, x_t x_{t'} - f_t x_{t'})$. Un calcul (omis) montre
que cette application est surjective. Notre choix de l'application affichée
ci-dessus détermine les applications $d_0, d_1, d_2 : P_2 \to P_1$.
Enfin, la procédure de la démonstration de la proposition
\ref{proposition-polynomial} nous dit de choisir les applications
$s_0, s_1 : P_1 \to P_2$ relevant les deux applications
$P_1 \to \text{cosk}_1(P_\bullet)_2$. Il est clair que l'on peut prendre pour
$s_i$ les uniques homomorphismes de $A$-algèbres déterminés par
$s_0(x_t) = y_t$ et $s_1(x_t) = z_t$.
\end{example}








\section{Fonctorialité}
\label{section-functoriality}

\noindent
Dans cette section, nous considérons un carré commutatif
\begin{equation}
\label{equation-commutative-square}
\vcenter{
\xymatrix{
B \ar[r] & B' \\
A \ar[u] \ar[r] & A' \ar[u]
}
}
\end{equation}
d'homomorphismes d'anneaux. Nous affirmons qu'il existe une application
canonique $B$-linéaire de complexes
$$
L_{B/A} \longrightarrow L_{B'/A'}
$$
associée à ce diagramme. En effet, si $P_\bullet \to B$ est la résolution
standard de $B$ sur $A$ et si $P'_\bullet \to B'$ est la résolution standard
de $B'$ sur $A'$, alors il existe une application canonique
$P_\bullet \to P'_\bullet$
de $A$-algèbres simpliciales compatible avec les augmentations
$P_\bullet \to B$ et $P'_\bullet \to B'$. Cela se voit au moyen de la
construction des résolutions standard dans Objets simpliciaux, section
\ref{simplicial-section-standard}, mais, dans le cas particulier qui nous
occupe, il suffit probablement de dire simplement que les applications
$$
P_0 = A[B] \longrightarrow A'[B'] = P'_0,\quad
P_1 = A[A[B]] \longrightarrow A'[A'[B']] = P'_1,
$$
et ainsi de suite sont données par les homomorphismes prescrits $A \to A'$ et
$B \to B'$. L'application souhaitée $L_{B/A} \to L_{B'/A'}$ provient alors
des applications associées $\Omega_{P_n/A} \to \Omega_{P'_n/A'}$.

\medskip\noindent
On peut décrire autrement l'application de fonctorialité comme suit. Soient
$\mathcal{C} = \mathcal{C}_{B/A}$ et $\mathcal{C}' = \mathcal{C}_{B'/A}'$
les catégories considérées dans la section
\ref{section-compute-L-pi-shriek}. Il existe un foncteur
$$
u : \mathcal{C} \longrightarrow \mathcal{C}',\quad
(P, \alpha) \longmapsto (P \otimes_A A', c \circ (\alpha \otimes 1))
$$
où $c : B \otimes_A A' \to B'$ est l'application évidente. Comme il est
expliqué dans Cohomologie sur les sites, exemple
\ref{sites-cohomology-example-morphism-categories}
on obtient un morphisme de topos $g : \Sh(\mathcal{C}) \to \Sh(\mathcal{C}')$
et un diagramme commutatif d'applications de topos annelés
\begin{equation}
\label{equation-double-square}
\vcenter{
\xymatrix{
(\Sh(\mathcal{C}'), \underline{B}) \ar[d]_{\pi'} &
(\Sh(\mathcal{C}'), \underline{B'}) \ar[d]_{\pi'} \ar[l]^h &
(\Sh(\mathcal{C}), \underline{B'}) \ar[d]_\pi \ar[l]^g \\
(\Sh(*), B) &
(\Sh(*), B') \ar[l]_f &
(\Sh(*), B') \ar[l]
}
}
\end{equation}
Ici, $h$ est l'identité sur les topos sous-jacents et est donné, sur les
faisceaux d'anneaux, par l'homomorphisme d'anneaux $B \to B'$. 
D'après Cohomologie sur les sites, remarque
\ref{sites-cohomology-remark-morphism-fibred-categories}
étant donnés $\mathcal{F}$ sur $\mathcal{C}$ et $\mathcal{F}'$ sur
$\mathcal{C}'$, ainsi qu'une transformation
$t : \mathcal{F} \to g^{-1}\mathcal{F}'$, on obtient une application canonique
$L\pi_!(\mathcal{F}) \to L\pi'_!(\mathcal{F}')$. Appliquons ceci aux faisceaux
$$
\mathcal{F} : (P, \alpha) \mapsto \Omega_{P/A} \otimes_P B,\quad
\mathcal{F}' : (P', \alpha') \mapsto \Omega_{P'/A'} \otimes_{P'} B',
$$
et à la transformation $t$ donnée par les applications canoniques
$$
\Omega_{P/A} \otimes_P B \longrightarrow
\Omega_{P \otimes_A A'/A'} \otimes_{P \otimes_A A'} B'
$$
pour obtenir une application canonique
$$
L\pi_!(\Omega_{\mathcal{O}/A} \otimes_\mathcal{O} \underline{B})
\longrightarrow
L\pi'_!(\Omega_{\mathcal{O}'/A'} \otimes_{\mathcal{O}'} \underline{B'})
$$
D'après le lemme \ref{lemma-compute-cotangent-complex}, cela fournit
$L_{B/A} \to L_{B'/A'}$. Nous omettons de vérifier que cette application
coïncide avec celle définie ci-dessus à l'aide de résolutions simpliciales.

\begin{lemma}
\label{lemma-flat-base-change}
Supposons que (\ref{equation-commutative-square}) induise un quasi-isomorphisme
$B \otimes_A^\mathbf{L} A' = B'$. Alors, avec les notations de
(\ref{equation-double-square}) et pour
$\mathcal{F}' \in \textit{Ab}(\mathcal{C}')$, on a
$L\pi_!(g^{-1}\mathcal{F}') = L\pi'_!(\mathcal{F}')$.
\end{lemma}

\begin{proof}
Nous utiliserons sans plus le signaler les résultats de la remarque
\ref{remark-resolution}. Appliquons Cohomologie sur les sites, lemme
\ref{sites-cohomology-lemma-get-it-now}. Soit $P_\bullet \to B$ une
résolution. Il suffit de montrer que
$u(P_\bullet) = P_\bullet \otimes_A A' \to B'$ est un quasi-isomorphisme. Le
complexe de $A$-modules $s(P_\bullet)$ associé à $P_\bullet$ (considéré comme
un $A$-module simplicial) est une résolution par des $A$-modules libres de
$B$. En effet, $P_n$ est un $A$-module libre et
$s(P_\bullet) \to B$ est un quasi-isomorphisme. Ainsi,
$B \otimes_A^\mathbf{L} A'$ est calculé par
$s(P_\bullet) \otimes_A A' = s(P_\bullet \otimes_A A')$. L'hypothèse du
lemme signifie donc que $\epsilon' : P_\bullet \otimes_A A' \to B'$ est un
quasi-isomorphisme.
\end{proof}

\noindent
Le lemme suivant s'applique en particulier lorsque $A \to A'$ est plat et que
$B' = B \otimes_A A'$ (changement de base plat).

\begin{lemma}
\label{lemma-flat-base-change-cotangent-complex}
Si (\ref{equation-commutative-square}) induit un quasi-isomorphisme
$B \otimes_A^\mathbf{L} A' = B'$, alors l'application de fonctorialité induit
un isomorphisme
$$
L_{B/A} \otimes_B^\mathbf{L} B' \longrightarrow L_{B'/A'}
$$
\end{lemma}

\begin{proof}
Nous utiliserons les notations introduites dans l'équation
(\ref{equation-double-square}). On a
$$
L_{B/A} \otimes_B^\mathbf{L} B' =
L\pi_!(\Omega_{\mathcal{O}/A} \otimes_\mathcal{O} \underline{B})
\otimes_B^\mathbf{L} B' =
L\pi_!(Lh^*(\Omega_{\mathcal{O}/A} \otimes_\mathcal{O} \underline{B}))
$$
où la première égalité résulte du lemme
\ref{lemma-compute-cotangent-complex} et la seconde de Cohomologie sur les
sites, lemme \ref{sites-cohomology-lemma-change-of-rings}. Comme
$\Omega_{\mathcal{O}/A}$ est un $\mathcal{O}$-module plat, on voit que
$\Omega_{\mathcal{O}/A} \otimes_\mathcal{O} \underline{B}$ est un
$\underline{B}$-module plat. Ainsi,
$Lh^*(\Omega_{\mathcal{O}/A} \otimes_\mathcal{O} \underline{B}) =
\Omega_{\mathcal{O}/A} \otimes_\mathcal{O} \underline{B'}$
qui est égal à
$g^{-1}(\Omega_{\mathcal{O}'/A'} \otimes_{\mathcal{O}'} \underline{B'})$
par inspection. On conclut par le lemme \ref{lemma-flat-base-change} et par le
fait que $L_{B'/A'}$ est calculé par
$L\pi'_!(\Omega_{\mathcal{O}'/A'} \otimes_{\mathcal{O}'} \underline{B'})$.
\end{proof}

\begin{remark}
\label{remark-homotopy-triangle}
Supposons donné un carré (\ref{equation-commutative-square}) tel qu'il existe
une flèche $\kappa : B \to A'$ qui rende le diagramme commutatif :
$$
\xymatrix{
B \ar[r]_\beta \ar[rd]_\kappa & B' \\
A \ar[u] \ar[r]^\alpha & A' \ar[u]
}
$$
Dans ce cas, nous affirmons que l'application de fonctorialité
$P_\bullet \to P'_\bullet$ est homotope à la composée
$P_\bullet \to B \to A' \to P'_\bullet$. En effet, au moyen de $\kappa$,
l'application de fonctorialité se factorise sous la forme
$$
P_\bullet \to P_{A'/A', \bullet} \to P'_\bullet
$$
où $P_{A'/A', \bullet}$ est la résolution standard de $A'$ sur $A'$. Comme
$A'$ est l'algèbre polynomiale sur l'ensemble vide sur $A'$, on déduit
d'Objets simpliciaux, lemme
\ref{simplicial-lemma-standard-simplicial-homotopy}, que l'augmentation
$\epsilon_{A'/A'} : P_{A'/A', \bullet} \to A'$ est une équivalence d'homotopie
d'anneaux simpliciaux. Observons que l'application inverse à homotopie près
$c : A' \to P_{A'/A', \bullet}$ construite dans la démonstration de ce lemme
n'est autre que le morphisme structural ; on obtient donc le résultat voulu,
car les deux composées
$$
\xymatrix{
P_\bullet \ar[r] &
P_{A'/A', \bullet} \ar@<1ex>[rr]^{\text{id}}
\ar@<-1ex>[rr]_{c \circ \epsilon_{A'/A'}} & &
P_{A'/A', \bullet} \ar[r] &
P'_\bullet
}
$$
sont les deux applications considérées ci-dessus, et elles sont homotopes
(Objets simpliciaux, remarque
\ref{simplicial-remark-homotopy-pre-post-compose}). Puisque la seconde
application $P_\bullet \to P'_\bullet$ induit l'application nulle
$\Omega_{P_\bullet/A} \to \Omega_{P'_\bullet/A'}$, on conclut que, dans ce
cas, l'application de fonctorialité $L_{B/A} \to L_{B'/A'}$ est homotope à
zéro.
\end{remark}

\begin{lemma}
\label{lemma-cotangent-complex-product}
Soient $A \to B$ et $A \to C$ des homomorphismes d'anneaux. Alors
l'application $L_{B \times C/A} \to L_{B/A} \oplus L_{C/A}$ est un
isomorphisme dans $D(B \times C)$.
\end{lemma}

\begin{proof}
Bien que ce lemme puisse se déduire du triangle fondamental, nous allons en
donner maintenant une démonstration directe et élémentaire. Factorisons
l'homomorphisme d'anneaux $A \to B \times C$ sous la forme
$A \to A[x] \to B \times C$, où $x \mapsto (1, 0)$. D'après le lemme
\ref{lemma-special-case-triangle}, on a un triangle distingué
$$
L_{A[x]/A} \otimes_{A[x]}^\mathbf{L} (B \times C) \to L_{B \times C/A} \to
L_{B \times C/A[x]} \to L_{A[x]/A} \otimes_{A[x]}^\mathbf{L} (B \times C)[1]
$$
dans $D(B \times C)$. De même, on a les triangles distingués
$$
\begin{matrix}
L_{A[x]/A} \otimes_{A[x]}^\mathbf{L} B \to L_{B/A} \to L_{B/A[x]}
\to L_{A[x]/A} \otimes_{A[x]}^\mathbf{L} B[1] \\
L_{A[x]/A} \otimes_{A[x]}^\mathbf{L} C \to L_{C/A} \to L_{C/A[x]}
\to L_{A[x]/A} \otimes_{A[x]}^\mathbf{L} C[1]
\end{matrix}
$$
Il suffit donc de démontrer le résultat pour $B \times C$ sur $A[x]$.
Notons que $A[x] \to A[x, x^{-1}]$ est plat, que
$(B \times C) \otimes_{A[x]} A[x, x^{-1}] = B \otimes_{A[x]} A[x, x^{-1}]$,
et que $C \otimes_{A[x]} A[x, x^{-1}] = 0$. Par changement de base (lemme
\ref{lemma-flat-base-change-cotangent-complex}), l'application
$L_{B \times C/A[x]} \to L_{B/A[x]} \oplus L_{C/A[x]}$ devient un
isomorphisme après inversion de $x$. De la même manière, on montre que
l'application devient un isomorphisme après inversion de $x - 1$. Cela
démontre le lemme.
\end{proof}




\section{Le triangle fondamental}
\label{section-triangle}

\noindent
Dans cette section, nous considérons une suite d'homomorphismes d'anneaux
$A \to B \to C$. Notre but est de montrer que ce triangle donne naissance
à un triangle distingué
\begin{equation}
\label{equation-triangle}
L_{B/A} \otimes_B^\mathbf{L} C \to L_{C/A} \to L_{C/B} \to
L_{B/A} \otimes_B^\mathbf{L} C[1]
\end{equation}
dans $D(C)$. Cela sera démontré dans la proposition
\ref{proposition-triangle}. Pour une autre approche, voir la remarque
\ref{remark-triangle}.

\medskip\noindent
Considérons la catégorie $\mathcal{C}_{C/B/A}$, qui est la catégorie
{\bf opposée} de la catégorie dont les objets sont les
$(P \to B, Q \to C)$ tels que
\begin{enumerate}
\item $P$ soit une algèbre polynomiale sur $A$,
\item $P \to B$ soit un homomorphisme de $A$-algèbres,
\item $Q$ soit une algèbre polynomiale sur $P$, et
\item $Q \to C$ soit un homomorphisme de $P$-algèbres.
\end{enumerate}
Nous prenons la catégorie opposée parce que nous voulons voir
$(P \to B, Q \to C)$ comme correspondant au diagramme commutatif
$$
\xymatrix{
\Spec(C) \ar[d] \ar[r] & \Spec(Q) \ar[d] \\
\Spec(B) \ar[d] \ar[r] & \Spec(P) \ar[dl] \\
\Spec(A)
}
$$
Soient $\mathcal{C}_{B/A}$, $\mathcal{C}_{C/A}$ et $\mathcal{C}_{C/B}$
les catégories considérées dans la section
\ref{section-compute-L-pi-shriek}. Il existe des foncteurs
$$
\begin{matrix}
u_1 : \mathcal{C}_{C/B/A} \to \mathcal{C}_{B/A}, &
(P \to B, Q \to C) \mapsto (P \to B) \\
u_2 : \mathcal{C}_{C/B/A} \to \mathcal{C}_{C/A}, &
(P \to B, Q \to C) \mapsto (Q \to C) \\
u_3 : \mathcal{C}_{C/B/A} \to \mathcal{C}_{C/B}, &
(P \to B, Q \to C) \mapsto (Q \otimes_P B \to C)
\end{matrix}
$$
Ces foncteurs induisent des morphismes de topos correspondants $g_i$.
Posons $\mathcal{O}_i = g_i^{-1}\mathcal{O}$ ; nous obtenons ainsi des
morphismes de topos annelés
\begin{equation}
\label{equation-three-maps}
\begin{matrix}
g_1 : (\Sh(\mathcal{C}_{C/B/A}), \mathcal{O}_1)
\longrightarrow (\Sh(\mathcal{C}_{B/A}), \mathcal{O}) \\
g_2 : (\Sh(\mathcal{C}_{C/B/A}), \mathcal{O}_2)
\longrightarrow (\Sh(\mathcal{C}_{C/A}), \mathcal{O}) \\
g_3 : (\Sh(\mathcal{C}_{C/B/A}), \mathcal{O}_3)
\longrightarrow (\Sh(\mathcal{C}_{C/B}), \mathcal{O})
\end{matrix}
\end{equation}
Notons
$\pi : \Sh(\mathcal{C}_{C/B/A}) \to \Sh(*)$,
$\pi_1 : \Sh(\mathcal{C}_{B/A}) \to \Sh(*)$,
$\pi_2 : \Sh(\mathcal{C}_{C/A}) \to \Sh(*)$, et
$\pi_3 : \Sh(\mathcal{C}_{C/B}) \to \Sh(*)$,
de sorte que $\pi = \pi_i \circ g_i$. Nous obtiendrons notre triangle
distingué à partir de l'identification du complexe cotangent dans le lemme
\ref{lemma-compute-cotangent-complex} et des lemmes suivants.

\begin{lemma}
\label{lemma-triangle-ses}
Avec les notations de (\ref{equation-three-maps}), posons
$$
\begin{matrix}
\Omega_1 = \Omega_{\mathcal{O}/A} \otimes_\mathcal{O} \underline{B}
\text{ sur }\mathcal{C}_{B/A} \\
\Omega_2 = \Omega_{\mathcal{O}/A} \otimes_\mathcal{O} \underline{C}
\text{ sur }\mathcal{C}_{C/A} \\
\Omega_3 = \Omega_{\mathcal{O}/B} \otimes_\mathcal{O} \underline{C}
\text{ sur }\mathcal{C}_{C/B}
\end{matrix}
$$
Alors nous avons une suite exacte courte canonique de faisceaux de
$\underline{C}$-modules
$$
0 \to g_1^{-1}\Omega_1 \otimes_{\underline{B}} \underline{C} \to
g_2^{-1}\Omega_2 \to
g_3^{-1}\Omega_3 \to 0
$$
sur $\mathcal{C}_{C/B/A}$.
\end{lemma}

\begin{proof}
Rappelons que $g_i^{-1}$ s'obtient simplement par précomposition avec $u_i$.
Pour un objet $U = (P \to B, Q \to C)$, nous avons une suite exacte courte
scindée
$$
0 \to \Omega_{P/A} \otimes Q \to \Omega_{Q/A} \to \Omega_{Q/P} \to 0
$$
par exemple d'après Algèbre, lemme
\ref{algebra-lemma-ses-formally-smooth}. En tensorisant par $C$ sur $Q$,
nous obtenons une suite exacte courte
$$
0 \to \Omega_{P/A} \otimes C \to \Omega_{Q/A} \otimes C \to
\Omega_{Q/P} \otimes C \to 0
$$
Nous avons $\Omega_{P/A} \otimes C = \Omega_{P/A} \otimes B \otimes C$ ;
c'est donc la valeur de
$g_1^{-1}\Omega_1 \otimes_{\underline{B}} \underline{C}$
en $U$. Le module $\Omega_{Q/A} \otimes C$ est la valeur de
$g_2^{-1}\Omega_2$ en $U$. Nous avons
$\Omega_{Q/P} \otimes C = \Omega_{Q \otimes_P B/B} \otimes C$
d'après Algèbre, lemme \ref{algebra-lemma-differentials-base-change} ;
c'est donc la valeur de $g_3^{-1}\Omega_3$ en $U$. La suite exacte courte
du lemme s'obtient ainsi en associant à $U$ la dernière suite exacte courte
affichée.
\end{proof}

\begin{lemma}
\label{lemma-polynomial-on-top}
Avec les notations de (\ref{equation-three-maps}), supposons que $C$ soit
une algèbre polynomiale sur $B$. Alors
$L\pi_!(g_3^{-1}\mathcal{F}) = L\pi_{3, !}\mathcal{F} = \pi_{3, !}\mathcal{F}$
pour tout faisceau abélien $\mathcal{F}$ sur $\mathcal{C}_{C/B}$.
\end{lemma}

\begin{proof}
Écrivons $C = B[E]$ pour un certain ensemble $E$. Choisissons une résolution
$P_\bullet \to B$ de $B$ sur $A$. Pour tout $n$, considérons l'objet
$U_n = (P_n \to B, P_n[E] \to C)$ de $\mathcal{C}_{C/B/A}$. Alors
$U_\bullet$ est un objet cosimplicial de $\mathcal{C}_{C/B/A}$. Notons que
$u_3(U_\bullet)$ est l'objet cosimplicial constant de
$\mathcal{C}_{C/B}$ de valeur $(C \to C)$. Nous allons démontrer que l'objet
$U_\bullet$ de $\mathcal{C}_{C/B/A}$ satisfait aux hypothèses de
Cohomologie sur les sites, lemme
\ref{sites-cohomology-lemma-compute-by-cosimplicial-resolution}.
Cela implique le lemme, car on voit ainsi que $L\pi_!(g_3^{-1}\mathcal{F})$
est calculé par le groupe abélien simplicial constant
$\mathcal{F}(C \to C)$, qui est la valeur de
$L\pi_{3, !}\mathcal{F} = \pi_{3, !}\mathcal{F}$ d'après le lemme
\ref{lemma-pi-lower-shriek-polynomial-algebra}.

\medskip\noindent
Soit $U = (\beta : P \to B, \gamma : Q \to C)$ un objet de
$\mathcal{C}_{C/B/A}$. Nous pouvons écrire $P = A[S]$ et
$Q = A[S \amalg T]$, par définition de notre catégorie
$\mathcal{C}_{C/B/A}$. Nous devons montrer que
$$
\Mor_{\mathcal{C}_{C/B/A}}(U_\bullet, U)
$$
est homotopiquement équivalent au singleton simplicial $*$. Observons que
cet ensemble simplicial est le produit
$$
\prod\nolimits_{s \in S} F_s \times \prod\nolimits_{t \in T} F'_t
$$
où $F_s$ est l'ensemble simplicial correspondant à
$U_s = (A[\{s\}] \to B, A[\{s\}] \to C)$
et $F'_t$ est l'ensemble simplicial correspondant à
$U_t = (A \to B, A[\{t\}] \to C)$. En effet, l'objet $U$ est le produit
$\prod U_s \times \prod U_t$ dans $\mathcal{C}_{C/B/A}$. Il suffit que
chaque $F_s$ et chaque $F'_t$ soit homotopiquement équivalent à $*$ ; voir
Simplicial, lemme \ref{simplicial-lemma-products-homotopy}. Le cas de $F_s$
résulte de ce que $P_\bullet \to B$ est une fibration de Kan triviale
(puisqu'il s'agit d'une résolution) et que $F_s$ est la fibre de cette
application au-dessus de $\beta(s)$. (Utiliser Simplicial, lemmes
\ref{simplicial-lemma-trivial-kan-base-change} et
\ref{simplicial-lemma-trivial-kan-homotopy}).
Le cas de $F'_t$ est plus intéressant. Il s'agit ici d'affirmer que la
fibre de
$$
P_\bullet[E] \longrightarrow C = B[E]
$$
au-dessus de $\gamma(t) \in C$ est homotopiquement équivalente à un point.
En fait, nous allons montrer que cette application est une fibration de Kan
triviale. En effet, $P_\bullet \to B$ est une fibration de Kan triviale.
Pour tout anneau $R$, nous avons
$$
R[E] =
\colim_{\Sigma \subset \text{Map}(E, \mathbf{Z}_{\geq 0})\text{ finite}}
\prod\nolimits_{I \in \Sigma} R
$$
(colimite filtrante). Ainsi, l'application affichée d'ensembles simpliciaux
est une colimite filtrante de fibrations de Kan triviales, donc une fibration
de Kan triviale d'après Simplicial, lemme
\ref{simplicial-lemma-filtered-colimit-trivial-kan}.
\end{proof}

\begin{lemma}
\label{lemma-triangle-compute-lower-shriek}
Avec les notations de (\ref{equation-three-maps}), nous avons
$Lg_{i, !} \circ g_i^{-1} = \text{id}$ for $i = 1, 2, 3$
et par conséquent aussi $L\pi_! \circ g_i^{-1} = L\pi_{i, !}$ pour
$i = 1, 2, 3$.
\end{lemma}

\begin{proof}
Démonstration pour $i = 1$. Nous affirmons que le foncteur
$\mathcal{C}_{C/B/A}$ est une catégorie fibrée sur $\mathcal{C}_{B/A}$.
En effet, supposons donnés $(P \to B, Q \to C)$ et un morphisme
$(P' \to B) \to (P \to B)$ de $\mathcal{C}_{B/A}$. Rappelons que cela
signifie que nous avons un homomorphisme de $A$-algèbres $P \to P'$
compatible aux applications vers $B$. Posons alors $Q' = Q \otimes_P P'$,
muni de l'application induite vers $C$ ; le morphisme
$$
(P' \to B, Q' \to C) \longrightarrow (P \to B, Q \to C)
$$
dans $\mathcal{C}_{C/B/A}$ (noter à nouveau le renversement des flèches)
est fortement cartésien dans $\mathcal{C}_{C/B/A}$ au-dessus de
$\mathcal{C}_{B/A}$. Observons en outre que la catégorie fibre de $u_1$
au-dessus de $P \to B$ est la catégorie $\mathcal{C}_{C/P}$. Soit
$\mathcal{F}$ un faisceau abélien sur $\mathcal{C}_{B/A}$. Puisque nous
avons une catégorie fibrée, nous pouvons appliquer Cohomologie sur les sites,
lemme \ref{sites-cohomology-lemma-compute-left-derived-pi-shriek}. Ainsi,
$L_ng_{1, !}g_1^{-1}\mathcal{F}$ est le (pré)faisceau qui associe à
$U \in \Ob(\mathcal{C}_{B/A})$ la $n$-ième homologie de
$g_1^{-1}\mathcal{F}$ restreint à la catégorie fibre au-dessus de $U$.
Comme ces restrictions sont constantes, le résultat voulu découle du lemme
\ref{lemma-pi-lower-shriek-constant-sheaf}, grâce aux identifications des
catégories fibres données ci-dessus.

\medskip\noindent
Le cas $i = 2$. Nous affirmons que $\mathcal{C}_{C/B/A}$ est une catégorie
fibrée sur $\mathcal{C}_{C/A}$. En effet, supposons donnés
$(P \to B, Q \to C)$ et un morphisme $(Q' \to C) \to (Q \to C)$ de
$\mathcal{C}_{C/A}$. Rappelons que cela signifie que nous avons un
homomorphisme de $B$-algèbres $Q \to Q'$ compatible aux applications vers
$C$. Alors
$$
(P \to B, Q' \to C) \longrightarrow (P \to B, Q \to C)
$$
est fortement cartésien dans $\mathcal{C}_{C/B/A}$ au-dessus de
$\mathcal{C}_{C/A}$. Notons que la catégorie fibre de $u_2$ au-dessus de
$Q \to C$ possède un objet final (attention au renversement des flèches),
à savoir $(A \to B, Q \to C)$. Soit $\mathcal{F}$ un faisceau abélien sur
$\mathcal{C}_{C/A}$. Puisque nous avons une catégorie fibrée, nous pouvons
appliquer Cohomologie sur les sites, lemme
\ref{sites-cohomology-lemma-compute-left-derived-pi-shriek}. Ainsi,
$L_ng_{2, !}g_2^{-1}\mathcal{F}$ est le (pré)faisceau qui associe à
$U \in \Ob(\mathcal{C}_{C/A})$ la $n$-ième homologie de
$g_1^{-1}\mathcal{F}$ restreint à la catégorie fibre au-dessus de $U$.
Comme ces restrictions sont constantes, le résultat voulu découle de
Cohomologie sur les sites, lemme
\ref{sites-cohomology-lemma-initial-final}, puisque les catégories fibres
possèdent toutes des objets finals.

\medskip\noindent
Le cas $i = 3$. Dans ce cas, nous appliquerons Cohomologie sur les sites,
lemme
\ref{sites-cohomology-lemma-compute-left-derived-g-shriek}
à $u = u_3 : \mathcal{C}_{C/B/A} \to \mathcal{C}_{C/B}$ et à
$\mathcal{F}' = g_3^{-1}\mathcal{F}$, où $\mathcal{F}$ est un faisceau
abélien sur $\mathcal{C}_{C/B}$. Supposons que
$U = (\overline{Q} \to C)$ soit un objet de $\mathcal{C}_{C/B}$. Alors
$\mathcal{I}_U = \mathcal{C}_{\overline{Q}/B/A}$ (attention encore au
renversement des flèches). Le faisceau $\mathcal{F}'_U$ est donné par la
règle $(P \to B, Q \to \overline{Q}) \mapsto \mathcal{F}(Q \otimes_P B \to C)$.
Autrement dit, ce faisceau est l'image inverse d'un faisceau sur
$\mathcal{C}_{\overline{Q}/C}$ par le morphisme
$\Sh(\mathcal{C}_{\overline{Q}/B/A}) \to \Sh(\mathcal{C}_{\overline{Q}/B})$.
Le lemme \ref{lemma-polynomial-on-top} montre donc que
$H_n(\mathcal{I}_U, \mathcal{F}'_U) = 0$ for $n > 0$
et est égal à $\mathcal{F}(\overline{Q} \to C)$ pour $n = 0$. Le lemme de
Cohomologie sur les sites déjà cité,
\ref{sites-cohomology-lemma-compute-left-derived-g-shriek}
implique que $Lg_{3, !}(g_3^{-1}\mathcal{F}) = \mathcal{F}$, ce qui achève
la démonstration.
\end{proof}

\begin{proposition}
\label{proposition-triangle}
Soient $A \to B \to C$ des homomorphismes d'anneaux. Il existe un triangle
distingué canonique
$$
L_{B/A} \otimes_B^\mathbf{L} C \to L_{C/A} \to L_{C/B} \to
L_{B/A} \otimes_B^\mathbf{L} C[1]
$$
dans $D(C)$.
\end{proposition}

\begin{proof}
Considérons la suite exacte courte de faisceaux du lemme
\ref{lemma-triangle-ses} et appliquons-lui le foncteur dérivé $L\pi_!$ ;
nous obtenons un triangle distingué
$$
L\pi_!(g_1^{-1}\Omega_1 \otimes_{\underline{B}} \underline{C}) \to
L\pi_!(g_2^{-1}\Omega_2) \to
L\pi_!(g_3^{-1}\Omega_3) \to
L\pi_!(g_1^{-1}\Omega_1 \otimes_{\underline{B}} \underline{C})[1]
$$
dans $D(C)$. Les lemmes \ref{lemma-triangle-compute-lower-shriek} et
\ref{lemma-compute-cotangent-complex}
montrent que les deuxième et troisième termes s'identifient respectivement
à $L_{C/A}$ et $L_{C/B}$, tandis que le premier est égal à
$$
L\pi_{1, !}(\Omega_1 \otimes_{\underline{B}} \underline{C}) =
L\pi_{1, !}(\Omega_1) \otimes_B^\mathbf{L} C =
L_{B/A} \otimes_B^\mathbf{L} C
$$
La première égalité résulte de Cohomologie sur les sites, lemme
\ref{sites-cohomology-lemma-change-of-rings}
(et de la platitude de $\Omega_1$ comme faisceau de modules sur
$\underline{B}$), et la seconde du lemme
\ref{lemma-compute-cotangent-complex}.
\end{proof}

\begin{remark}
\label{remark-triangle}
Esquissons une autre démonstration, peut-être plus simple, de l'existence du
triangle fondamental. Soient $A \to B \to C$ des homomorphismes d'anneaux et
supposons que $B \to C$ soit injectif. Soit $P_\bullet \to B$ la résolution
standard de $B$ sur $A$, et soit $Q_\bullet \to C$ la résolution standard de
$C$ sur $B$. Représentons-les ainsi :
$$
\xymatrix{
P_\bullet : &
A[A[A[B]]] \ar[d]
\ar@<2ex>[r]
\ar@<0ex>[r]
\ar@<-2ex>[r]
&
A[A[B]] \ar[d]
\ar@<1ex>[r]
\ar@<-1ex>[r]
\ar@<1ex>[l]
\ar@<-1ex>[l]
&
A[B] \ar[d] \ar@<0ex>[l] \ar[r] &
B \\
Q_\bullet : &
A[A[A[C]]]
\ar@<2ex>[r]
\ar@<0ex>[r]
\ar@<-2ex>[r]
&
A[A[C]]
\ar@<1ex>[r]
\ar@<-1ex>[r]
\ar@<1ex>[l]
\ar@<-1ex>[l]
&
A[C] \ar@<0ex>[l] \ar[r] &
C
}
$$
Puisque $B \to C$ est injectif, observons que l'anneau $Q_n$ est une algèbre
polynomiale sur $P_n$ pour tout $n$. Nous obtenons donc un objet cosimplicial
de $\mathcal{C}_{C/B/A}$ (attention au renversement des flèches). Posons
maintenant $\overline{Q}_\bullet = Q_\bullet \otimes_{P_\bullet} B$.
Le point clef de la démonstration de la proposition
\ref{proposition-triangle} consiste à montrer que $\overline{Q}_\bullet$ est
une résolution de $C$ sur $B$. Cela résulte de Cohomologie sur les sites,
lemme
\ref{sites-cohomology-lemma-O-homology-qis}
appliqué à $\mathcal{C} = \Delta$, $\mathcal{O} = P_\bullet$,
$\mathcal{O}' = B$ et $\mathcal{F} = Q_\bullet$ (on utilise ici que $Q_n$
est plat sur $P_n$ ; pour relier les modules simpliciaux aux faisceaux, voir
Cohomologie sur les sites, remarque
\ref{sites-cohomology-remark-simplicial-modules}). Ce fait clef implique que
le triangle distingué de la proposition \ref{proposition-triangle} est celui
qui est associé à la suite exacte courte de $C$-modules simpliciaux
$$
0 \to
\Omega_{P_\bullet/A} \otimes_{P_\bullet} C \to
\Omega_{Q_\bullet/A} \otimes_{Q_\bullet} C \to
\Omega_{\overline{Q}_\bullet/B} \otimes_{\overline{Q}_\bullet} C \to 0
$$
qui se déduit des suites exactes courtes
$0 \to \Omega_{P_n/A} \otimes_{P_n} Q_n \to \Omega_{Q_n/A} \to
\Omega_{Q_n/P_n} \to 0$ of
Algèbre, lemme \ref{algebra-lemma-ses-formally-smooth}.
En effet, d'après la remarque \ref{remark-resolution} et le fait clef, le
complexe du membre de droite représente $L_{C/B}$ dans $D(C)$.

\medskip\noindent
Si $B \to C$ n'est pas injectif, nous pouvons appliquer ce qui précède pour
obtenir un triangle fondamental pour $A \to B \to B \times C$. Puisque
$L_{B \times C/B} \to L_{B/B} \oplus L_{C/B}$ et
$L_{B \times C/A} \to L_{B/A} \oplus L_{C/A}$
constituent des quasi-isomorphismes dans $D(B \times C)$
(lemme \ref{lemma-cotangent-complex-product})
cela induit le triangle distingué voulu dans $D(C)$ en tensorisant par
l'homomorphisme plat d'anneaux $B \times C \to C$.
\end{remark}

\begin{remark}
\label{remark-explicit-map}
Soient $A \to B \to C$ des homomorphismes d'anneaux, avec $B \to C$
injectif. Rappelons les notations $P_\bullet$, $Q_\bullet$ et
$\overline{Q}_\bullet$ de la remarque \ref{remark-triangle}. Soit
$R_\bullet$ la résolution standard de $C$ sur $B$. Dans cette remarque,
nous expliquons comment obtenir l'identification canonique de
$\Omega_{\overline{Q}_\bullet/B} \otimes_{\overline{Q}_\bullet} C$ à
$L_{C/B} = \Omega_{R_\bullet/B} \otimes_{R_\bullet} C$. Soit
$S_\bullet \to B$ la résolution standard de $B$ sur $B$. Notons que
l'application de fonctorialité $S_\bullet \to R_\bullet$ fait de $R_n$ une
algèbre polynomiale sur $S_n$, puisque $B \to C$ est injectif.
Par exemple, en degré $0$, nous avons l'application $B[B] \to B[C]$ ; en
degré $1$, l'application $B[B[B]] \to B[B[C]]$ ; et ainsi de suite. Par
conséquent,
$\overline{R}_\bullet = R_\bullet \otimes_{S_\bullet} B$
est elle aussi une algèbre polynomiale simpliciale sur $B$, et il résulte
(comme dans la remarque \ref{remark-triangle}) de Cohomologie sur les sites,
lemme
\ref{sites-cohomology-lemma-O-homology-qis}
que $\overline{R}_\bullet \to C$ est une résolution. Comme nous avons un
diagramme commutatif
$$
\xymatrix{
Q_\bullet \ar[r] & R_\bullet \\
P_\bullet \ar[u] \ar[r] & S_\bullet \ar[u] \ar[r] & B
}
$$
nous obtenons une application canonique
$\overline{Q}_\bullet = Q_\bullet \otimes_{P_\bullet} B \to
\overline{R}_\bullet$. Les applications
$$
L_{C/B} = \Omega_{R_\bullet/B} \otimes_{R_\bullet} C
\longrightarrow
\Omega_{\overline{R}_\bullet/B} \otimes_{\overline{R}_\bullet} C
\longleftarrow
\Omega_{\overline{Q}_\bullet/B} \otimes_{\overline{Q}_\bullet} C
$$
sont alors des quasi-isomorphismes (remarque \ref{remark-resolution}), et
composer l'une avec l'inverse de l'autre donne l'identification voulue.
\end{remark}





\section{Localisation et homomorphismes d'anneaux \'etales}
\label{section-localization}

\noindent
Dans cette section, nous étudions ce qui se passe lorsque nous localisons nos
anneaux. Soient $A \to A' \to B$ des homomorphismes d'anneaux tels que
$B = B \otimes_A^\mathbf{L} A'$. C'est le cas, par exemple, si
$A' = S^{-1}A$ est le localisé de $A$ par une partie multiplicative
$S \subset A$. Dans ce cas, pour un faisceau abélien $\mathcal{F}'$ sur
$\mathcal{C}_{B/A'}$, l'homologie de $g^{-1}\mathcal{F}'$ sur
$\mathcal{C}_{B/A}$ coïncide avec celle de $\mathcal{F}'$ sur
$\mathcal{C}_{B/A'}$ ; voir le lemme \ref{lemma-flat-base-change} pour un
énoncé précis.

\begin{lemma}
\label{lemma-localize-at-bottom}
Soient $A \to A' \to B$ des homomorphismes d'anneaux tels que
$B = B \otimes_A^\mathbf{L} A'$. Alors $L_{B/A} = L_{B/A'}$ dans $D(B)$.
\end{lemma}

\begin{proof}
D'après la discussion précédente (c'est-à-dire en utilisant le lemme
\ref{lemma-flat-base-change}) et le lemme
\ref{lemma-compute-cotangent-complex}, il nous faut montrer que le faisceau
défini sur $\mathcal{C}_{B/A}$ par la règle
$(P \to B) \mapsto \Omega_{P/A} \otimes_P B$ est l'image inverse du faisceau
défini par la règle $(P \to B) \mapsto \Omega_{P/A'} \otimes_P B$.
Le foncteur image inverse $g^{-1}$ s'obtient en précomposant avec le foncteur
$u : \mathcal{C}_{B/A} \to \mathcal{C}_{B/A'}$,
$(P \to B) \mapsto (P \otimes_A A' \to B)$. Il nous faut donc montrer que
$$
\Omega_{P/A} \otimes_P B =
\Omega_{P \otimes_A A'/A'} \otimes_{(P \otimes_A A')} B
$$
D'après Algèbre, lemme \ref{algebra-lemma-differentials-base-change}, le
membre de droite est égal à
$$
(\Omega_{P/A} \otimes_A A') \otimes_{(P \otimes_A A')} B
$$
Comme $P$ est une algèbre polynomiale sur $A$, le module $\Omega_{P/A}$ est
libre et l'égalité est évidente.
\end{proof}

\begin{lemma}
\label{lemma-derived-diagonal}
Soit $A \to B$ un homomorphisme d'anneaux tel que
$B = B \otimes_A^\mathbf{L} B$. Alors $L_{B/A} = 0$ dans $D(B)$.
\end{lemma}

\begin{proof}
Cela résulte de $L_{B/A} = L_{B/B} = 0$, d'après les lemmes
\ref{lemma-localize-at-bottom} et
\ref{lemma-cotangent-complex-polynomial-algebra}.
\end{proof}

\begin{lemma}
\label{lemma-bootstrap}
Soit $A \to B$ un homomorphisme d'anneaux tel que
$\text{Tor}^A_i(B, B) = 0$ pour $i > 0$ et tel que
$L_{B/B \otimes_A B} = 0$. Alors $L_{B/A} = 0$ dans $D(B)$.
\end{lemma}

\begin{proof}
D'après le lemme \ref{lemma-flat-base-change-cotangent-complex}, on a
$L_{B/A} \otimes_B^\mathbf{L} (B \otimes_A B) = L_{B \otimes_A B/B}$.
Utilisons maintenant le triangle distingué (\ref{equation-triangle})
$$
L_{B \otimes_A B/B} \otimes^\mathbf{L}_{(B \otimes_A B)} B \to
L_{B/B} \to L_{B/B \otimes_A B} \to
L_{B \otimes_A B/B} \otimes^\mathbf{L}_{(B \otimes_A B)} B[1]
$$
associé aux homomorphismes d'anneaux $B \to B \otimes_A B \to B$, ainsi que
l'annulation de $L_{B/B}$ (lemme
\ref{lemma-cotangent-complex-polynomial-algebra}) et celle, supposée, de
$L_{B/B \otimes_A B}$, pour obtenir
$$
0 =
L_{B \otimes_A B/B} \otimes^\mathbf{L}_{(B \otimes_A B)} B =
L_{B/A} \otimes_B^\mathbf{L} (B \otimes_A B)
\otimes^\mathbf{L}_{(B \otimes_A B)} B = L_{B/A}
$$
comme voulu.
\end{proof}

\begin{lemma}
\label{lemma-when-zero}
Le complexe cotangent $L_{B/A}$ est nul dans chacun des cas suivants :
\begin{enumerate}
\item $A \to B$ et $B \otimes_A B \to B$ sont plats, c'est-à-dire que
$A \to B$ est faiblement \'etale (Compléments sur l'algèbre, définition
\ref{more-algebra-definition-weakly-etale}),
\item $A \to B$ est un épimorphisme plat d'anneaux,
\item $B = S^{-1}A$ pour une partie multiplicative $S \subset A$,
\item $A \to B$ est non ramifié et plat,
\item $A \to B$ est \'etale,
\item $A \to B$ est une colimite filtrante d'homomorphismes d'anneaux dont le
complexe cotangent s'annule,
\item $B$ est un hensélisé d'un anneau local de $A$,
\item $B$ est un hensélisé strict d'un anneau local de $A$, et
\item ajouter ici d'autres cas.
\end{enumerate}
\end{lemma}

\begin{proof}
Dans le cas (1), nous pouvons appliquer le lemme
\ref{lemma-derived-diagonal} à l'homomorphisme d'anneaux plat surjectif
$B \otimes_A B \to B$ pour conclure que $L_{B/B \otimes_A B} = 0$, puis le
lemme \ref{lemma-bootstrap} permet de conclure. Les cas (2) -- (5) sont tous
des cas particuliers de (1). L'assertion (6) résulte du lemme
\ref{lemma-colimit-cotangent-complex}. Les assertions (7) et (8) résultent du
fait que les hensélisés (stricts) sont des colimites filtrantes d'extensions
d'anneaux \'etales de $A$ ; voir Algèbre, lemmes
\ref{algebra-lemma-henselization-different} et
\ref{algebra-lemma-strict-henselization-different}.
\end{proof}

\begin{lemma}
\label{lemma-localize-on-top}
Soient $A \to B \to C$ des homomorphismes d'anneaux tels que $L_{C/B} = 0$.
Alors $L_{C/A} = L_{B/A} \otimes_B^\mathbf{L} C$.
\end{lemma}

\begin{proof}
C'est une conséquence immédiate du triangle distingué
(\ref{equation-triangle}).
\end{proof}

\begin{lemma}
\label{lemma-localize}
Soit $A \to B$ un homomorphisme d'anneaux, et soient $S \subset A$ et
$T \subset B$ des parties multiplicatives telles que $S$ s'envoie dans $T$.
Alors $L_{T^{-1}B/S^{-1}A} = L_{B/A} \otimes_B T^{-1}B$ dans $D(T^{-1}B)$.
\end{lemma}

\begin{proof}
D'après le lemme \ref{lemma-localize-on-top},
$L_{T^{-1}B/A} = L_{B/A} \otimes_B T^{-1}B$
et, d'après le lemme \ref{lemma-localize-at-bottom},
$L_{T^{-1}B/A} = L_{T^{-1}B/S^{-1}A}$.
\end{proof}

\begin{lemma}
\label{lemma-cotangent-complex-henselization}
Soit $A \to B$ un homomorphisme local d'anneaux locaux. Notons
$A^h \to B^h$, resp.\ $A^{sh} \to B^{sh}$, les homomorphismes induits entre
hensélisés, resp.\ hensélisés stricts. Alors
$$
L_{B^h/A^h} = L_{B^h/A} = L_{B/A} \otimes_B^\mathbf{L} B^h
\quad\text{resp.}\quad
L_{B^{sh}/A^{sh}} = L_{B^{sh}/A} = L_{B/A} \otimes_B^\mathbf{L} B^{sh}
$$
dans $D(B^h)$, resp.\ $D(B^{sh})$.
\end{lemma}

\begin{proof}
Les complexes $L_{A^h/A}$, $L_{A^{sh}/A}$, $L_{B^h/B}$ et
$L_{B^{sh}/B}$ sont tous nuls d'après le lemme \ref{lemma-when-zero}. En
appliquant le triangle distingué fondamental (\ref{equation-triangle}) à
$A \to B \to B^h$, nous obtenons
$L_{B^h/A} = L_{B/A} \otimes_B^\mathbf{L} B^h$.
Le triangle fondamental appliqué à $A \to A^h \to B^h$ donne
$L_{B^h/A^h} = L_{B^h/A}$. Le raisonnement est le même pour les hensélisés
stricts.
\end{proof}




\section{Homomorphismes d'anneaux lisses}
\label{section-smooth}

\noindent
Soit $C \to B$ une surjection d'anneaux de noyau $I$. Disons qu'un tel
homomorphisme d'anneaux est « faiblement quasi régulier » si $I/I^2$ est un
$B$-module plat et si $\text{Tor}_*^C(B, B)$ est l'algèbre extérieure de
$I/I^2$. Pour généraliser aux « homomorphismes d'anneaux lisses » ce qui est
fait dans le lemme \ref{lemma-when-zero} pour les « homomorphismes d'anneaux
\'etales », il convient de considérer les homomorphismes plats $A \to B$ tels
que l'homomorphisme de multiplication $B \otimes_A B \to B$ soit faiblement
quasi régulier. Pour le moment, nous nous en tenons aux homomorphismes
d'anneaux lisses.

\begin{lemma}
\label{lemma-when-projective}
Si $A \to B$ est un homomorphisme d'anneaux lisse, alors
$L_{B/A} = \Omega_{B/A}[0]$.
\end{lemma}

\begin{proof}
L'identification en degré cohomologique $0$ est donnée par le lemme
\ref{lemma-identify-H0}. Il suffit donc de montrer que les autres groupes de
cohomologie sont nuls. Il suffit de le faire localement sur $\Spec(B)$, puisque
$L_{B_g/A} = (L_{B/A})_g$ pour $g \in B$ d'après le lemme
\ref{lemma-localize-on-top}. Nous pouvons ainsi supposer que $A \to B$ est
lisse standard (Algèbre, lemme \ref{algebra-lemma-smooth-syntomic}),
c'est-à-dire que $A \to B$ se factorise sous la forme
$A \to A[x_1, \ldots, x_n] \to B$, où
$A[x_1, \ldots, x_n] \to B$ est \'etale. Dans ce cas, les lemmes
\ref{lemma-when-zero} et \ref{lemma-localize-on-top} montrent que
$L_{B/A} = L_{A[x_1, \ldots, x_n]/A} \otimes B$
d'où la conclusion par le lemme
\ref{lemma-cotangent-complex-polynomial-algebra}.
\end{proof}





\section{Caractéristique positive}
\label{section-positive-characteristic}

\noindent
Dans cette section, nous fixons un nombre premier $p$. Si $A$ est un anneau
tel que $p = 0$ dans $A$, alors $F_A : A \to A$ désigne l'endomorphisme de
Frobenius $a \mapsto a^p$.

\begin{lemma}
\label{lemma-frobenius-homotopy}
Soit $A \to B$ un homomorphisme d'anneaux tel que $p = 0$ dans $A$. Soit
$P_\bullet$ la résolution standard de $B$ sur $A$. L'application
$P_\bullet \to P_\bullet$ induite par le diagramme
$$
\xymatrix{
B \ar[r]_{F_B} & B \\
A \ar[u] \ar[r]^{F_A} & A \ar[u]
}
$$
considéré dans la section \ref{section-functoriality} est homotope à
l'endomorphisme de Frobenius $P_\bullet \to P_\bullet$ donné par Frobenius sur
chaque $P_n$.
\end{lemma}

\begin{proof}
Soit $\mathcal{A}$ la catégorie des homomorphismes de $\mathbf{F}_p$-algèbres
$A \to B$. Soit $\mathcal{S}$ la catégorie des couples $(A, E)$, où $A$ est
une $\mathbf{F}_p$-algèbre et $E$ un ensemble. Considérons les foncteurs
adjoints
$$
V : \mathcal{A} \to \mathcal{S}, \quad (A \to B) \mapsto (A, B)
$$
et
$$
U : \mathcal{S} \to \mathcal{A}, \quad (A, E) \mapsto (A \to A[E])
$$
Soit $X$ l'objet simplicial de la catégorie des foncteurs de $\mathcal{A}$
dans $\mathcal{A}$ construit dans Simplicial, section
\ref{simplicial-section-standard}. Il est clair que
$P_\bullet = X(A \to B)$ car, si l'on fixe $A$, alors.

\medskip\noindent
Posons $Y = U \circ V$. Rappelons que $X$ est construit à partir de $Y$ et de
certaines applications, et que ses termes sont
$X_n = Y \circ \ldots \circ Y$, avec $n + 1$ facteurs ; la construction est
donnée dans Simplicial, exemple \ref{simplicial-example-godement}, et l'on
trouvera les détails dans la démonstration de Simplicial, lemme
\ref{simplicial-lemma-standard-simplicial}.

\medskip\noindent
Soit $f : \text{id}_\mathcal{A} \to \text{id}_\mathcal{A}$ l'endomorphisme de
Frobenius du foncteur identité. Autrement dit, posons
$f_{A \to B} = (F_A, F_B) : (A \to B) \to (A \to B)$. Nos deux applications
sur $X(A \to B)$ sont alors données par les transformations naturelles
$f \star 1_X$ et $1_X \star f$. Nous omettons les détails. On conclut donc par
Simplicial, lemme \ref{simplicial-lemma-godement-before-after}.
\end{proof}

\begin{lemma}
\label{lemma-frobenius-acts-as-zero}
Soit $p$ un nombre premier. Soit $A \to B$ un homomorphisme d'anneaux et
supposons que $p = 0$ dans $A$. L'application $L_{B/A} \to L_{B/A}$ de la
section \ref{section-functoriality} induite par les applications de Frobenius
$F_A$ et $F_B$ est homotope à zéro.
\end{lemma}

\begin{proof}
Soit $P_\bullet$ la résolution standard de $B$ sur $A$. D'après le lemme
\ref{lemma-frobenius-homotopy}, l'application $P_\bullet \to P_\bullet$ induite
par $F_A$ et $F_B$ est homotope à l'application
$F_{P_\bullet} : P_\bullet \to P_\bullet$ donnée par Frobenius sur chaque
terme. On obtient donc le résultat voulu, puisque $F_{P_\bullet}$ induit
manifestement l'application nulle $\Omega_{P_n/A} \to \Omega_{P_n/A}$ (la
dérivée d'une puissance $p$-ième étant nulle).
\end{proof}

\begin{lemma}
\label{lemma-perfect-zero}
Soit $p$ un nombre premier. Soit $A \to B$ un homomorphisme d'anneaux et
supposons que $p = 0$ dans $A$. Si $A$ et $B$ sont parfaits, alors $L_{B/A}$
est nul dans $D(B)$.
\end{lemma}

\begin{proof}
L'application $(F_A, F_B) : (A \to B) \to (A \to B)$ est un isomorphisme et
induit donc un isomorphisme sur $L_{B/A}$ ; d'autre part, elle induit zéro sur
$L_{B/A}$ d'après le lemme \ref{lemma-frobenius-acts-as-zero}.
\end{proof}





\section{Comparaison avec le complexe cotangent naïf}
\label{section-surjections}

\noindent
Le complexe cotangent naïf a été introduit dans Algèbre, section
\ref{algebra-section-netherlander}.

\begin{remark}
\label{remark-make-map}
Soit $A \to B$ un homomorphisme d'anneaux. Travaillons sur
$\mathcal{C}_{B/A}$ comme dans la section
\ref{section-compute-L-pi-shriek}, et notons
$\mathcal{J} \subset \mathcal{O}$ le noyau de
$\mathcal{O} \to \underline{B}$. Remarquons que
$L\pi_!(\mathcal{J}) = 0$ d'après le lemme
\ref{lemma-apply-O-B-comparison}. Posons
$\Omega =  \Omega_{\mathcal{O}/A} \otimes_\mathcal{O} \underline{B}$
de sorte que
$L_{B/A} = L\pi_!(\Omega)$ d'après le lemme
\ref{lemma-compute-cotangent-complex}. Il s'ensuit que
$L\pi_!(\mathcal{J} \to \Omega) = L\pi_!(\Omega) = L_{B/A}$. Ainsi, pour tout
objet $U = (P \to B)$ de $\mathcal{C}_{B/A}$, nous obtenons une application
\begin{equation}
\label{equation-comparison-map-A}
(J \to \Omega_{P/A} \otimes_P B) \longrightarrow L_{B/A}
\end{equation}
où $J = \Ker(P \to B)$ dans $D(A)$ ; voir Cohomologie sur les sites, remarque
\ref{sites-cohomology-remark-map-evaluation-to-derived}. Poursuivant de cette
façon, remarquons que
$L\pi_!(\mathcal{J} \otimes_\mathcal{O}^\mathbf{L} \underline{B}) =
L\pi_!(\mathcal{J}) = 0$ by
lemme \ref{lemma-O-homology-B-homology}.
Comme $\text{Tor}_0^\mathcal{O}(\mathcal{J}, \underline{B}) =
\mathcal{J}/\mathcal{J}^2$, la suite spectrale
$$
H_p(\mathcal{C}_{B/A}, \text{Tor}_q^\mathcal{O}(\mathcal{J}, \underline{B}))
\Rightarrow 
H_{p + q}(\mathcal{C}_{B/A},
\mathcal{J} \otimes_\mathcal{O}^\mathbf{L} \underline{B}) = 0
$$
(duale de Catégories dérivées, lemme
\ref{derived-lemma-two-ss-complex-functor}) implique que
$H_0(\mathcal{C}_{B/A}, \mathcal{J}/\mathcal{J}^2) = 0$
et $H_1(\mathcal{C}_{B/A}, \mathcal{J}/\mathcal{J}^2) = 0$.
Il s'ensuit que le complexe de $\underline{B}$-modules
$\mathcal{J}/\mathcal{J}^2 \to \Omega$ vérifie
$\tau_{\geq -1}L\pi_!(\mathcal{J}/\mathcal{J}^2 \to \Omega) =
\tau_{\geq -1}L_{B/A}$.
Ainsi, pour tout objet $U = (P \to B)$ de $\mathcal{C}_{B/A}$, nous obtenons
une application
\begin{equation}
\label{equation-comparison-map}
(J/J^2 \to \Omega_{P/A} \otimes_P B) \longrightarrow \tau_{\geq -1}L_{B/A}
\end{equation}
dans $D(B)$ ; voir Cohomologie sur les sites, remarque
\ref{sites-cohomology-remark-map-evaluation-to-derived}.
\end{remark}

\noindent
Le premier cas est celui d'une surjection d'anneaux.

\begin{lemma}
\label{lemma-surjection}
\begin{slogan}
La cohomologie du complexe cotangent d'un homomorphisme d'anneaux surjectif est
nulle en degré zéro ; en degré $-1$, elle est le noyau modulo son carré.
\end{slogan}
Soit $A \to B$ un homomorphisme d'anneaux surjectif de noyau $I$. Alors
$H^0(L_{B/A}) = 0$ et $H^{-1}(L_{B/A}) = I/I^2$. Cet isomorphisme provient de
l'application (\ref{equation-comparison-map}) pour l'objet $(A \to B)$ de
$\mathcal{C}_{B/A}$.
\end{lemma}

\begin{proof}
Nous allons montrer ci-dessous (en utilisant la surjectivité de $A \to B$)
qu'il existe une suite exacte courte
$$
0 \to \pi^{-1}(I/I^2) \to \mathcal{J}/\mathcal{J}^2 \to \Omega \to 0
$$
de faisceaux sur $\mathcal{C}_{B/A}$. En appliquant $L\pi_!$, puis la suite
exacte longue d'homologie associée, et en utilisant l'annulation de
$H_1(\mathcal{C}_{B/A}, \mathcal{J}/\mathcal{J}^2)$ et de
$H_0(\mathcal{C}_{B/A}, \mathcal{J}/\mathcal{J}^2)$ établie dans la remarque
\ref{remark-make-map}, nous obtenons le résultat voulu grâce au lemme
\ref{lemma-pi-lower-shriek-constant-sheaf}.

\medskip\noindent
Il reste à vérifier l'énoncé local mentionné ci-dessus. Pour tout objet
$U = (P \to B)$ de $\mathcal{C}_{B/A}$, nous pouvons choisir un isomorphisme
$P = A[E]$ tel que l'homomorphisme $P \to B$ envoie chaque $e \in E$ sur zéro.
Alors
$J = \mathcal{J}(U) \subset P = \mathcal{O}(U)$
est égal à $J = IP + (e; e \in E)$. La valeur en $U$ de la suite exacte courte
de faisceaux ci-dessus est la suite
$$
0 \to I/I^2 \to J/J^2 \to \Omega_{P/A} \otimes_P B \to 0
$$
Nous omettons la vérification (indication : le seul point délicat est l'égalité
$IP \cap J^2 = IJ$, qui résulte par exemple de Compléments sur l'algèbre,
lemme \ref{more-algebra-lemma-conormal-sequence-H1-regular}).
\end{proof}

\begin{lemma}
\label{lemma-relation-with-naive-cotangent-complex}
Soit $A \to B$ un homomorphisme d'anneaux. Alors
$\tau_{\geq -1}L_{B/A}$ est canoniquement quasi-isomorphe au complexe
cotangent naïf.
\end{lemma}

\begin{proof}
Considérons $P = A[B] \to B$, de noyau $I$. Le complexe cotangent naïf
$\NL_{B/A}$ de $B$ sur $A$ est le complexe
$I/I^2 \to \Omega_{P/A} \otimes_P B$ ; voir Algèbre, définition
\ref{algebra-definition-naive-cotangent-complex}. Remarquons que nous avons
déjà construit dans (\ref{equation-comparison-map}) une application canonique
$$
c : \NL_{B/A} \longrightarrow \tau_{\geq -1}L_{B/A}
$$
Considérons le triangle distingué (\ref{equation-triangle})
$$
L_{P/A} \otimes_P^\mathbf{L} B \to L_{B/A} \to L_{B/P} \to 
(L_{P/A} \otimes_P^\mathbf{L} B)[1]
$$
associé aux homomorphismes d'anneaux $A \to A[B] \to B$. Nous savons que
$L_{P/A} = \Omega_{P/A}[0] = \NL_{P/A}$ in $D(P)$
(lemme \ref{lemma-cotangent-complex-polynomial-algebra} et Algèbre, lemme
\ref{algebra-lemma-NL-polynomial-algebra}) et que
$\tau_{\geq -1}L_{B/P} = I/I^2[1] = \NL_{B/P}$ in $D(B)$
(lemme \ref{lemma-surjection} et Algèbre, lemme
\ref{algebra-lemma-NL-surjection}). Pour montrer que $c$ est un
quasi-isomorphisme, il suffit, d'après Algèbre, lemme
\ref{algebra-lemma-exact-sequence-NL}, et la suite exacte longue de cohomologie
associée au triangle distingué, de montrer que les applications
$L_{P/A} \to L_{B/A} \to L_{B/P}$ sont compatibles sur les groupes de
cohomologie avec les applications correspondantes
$\NL_{P/A} \to \NL_{B/A} \to \NL_{B/P}$
du complexe cotangent naïf. Nous omettons la vérification.
\end{proof}

\begin{remark}
\label{remark-explicit-comparison-map}
Nous pouvons expliciter comme suit l'application de comparaison du lemme
\ref{lemma-relation-with-naive-cotangent-complex}. Soit $P_\bullet$ la
résolution standard de $B$ sur $A$. Posons $I = \Ker(A[B] \to B)$. Rappelons
que $P_0 = A[B]$. L'application du lemme est donnée par le diagramme
commutatif
$$
\xymatrix{
L_{B/A} \ar[d] & \ldots \ar[r] &
\Omega_{P_2/A} \otimes_{P_2} B
\ar[r] \ar[d] &
\Omega_{P_1/A} \otimes_{P_1} B
\ar[r] \ar[d] &
\Omega_{P_0/A} \otimes_{P_0} B
\ar[d] \\
\NL_{B/A} & \ldots \ar[r] &
0 \ar[r] & 
I/I^2 \ar[r] &
\Omega_{P_0/A} \otimes_{P_0} B
}
$$
Nous construisons la flèche descendante de but $I/I^2$ en envoyant
$\text{d}f \otimes b$ sur la classe de $(d_0(f) - d_1(f))b$ dans $I/I^2$.
Ici, $d_i : P_1 \to P_0$, $i = 0, 1$, sont les deux applications de face de
la structure simpliciale. Cela a un sens puisque $d_0 - d_1$ envoie $P_1$
dans $I = \Ker(P_0 \to B)$. Nous omettons de vérifier que cette règle est bien
définie. Notre application est compatible avec la différentielle
$\Omega_{P_1/A} \otimes_{P_1} B \to \Omega_{P_0/A} \otimes_{P_0} B$
car celle-ci envoie $\text{d}f \otimes b$ sur
$\text{d}(d_0(f) - d_1(f)) \otimes b$. En outre, la différentielle
$\Omega_{P_2/A} \otimes_{P_2} B \to \Omega_{P_1/A} \otimes_{P_1} B$
envoie $\text{d}f \otimes b$ sur
$\text{d}(d_0(f) - d_1(f) + d_2(f)) \otimes b$, qui est annulé par notre
flèche descendante. Nous obtenons ainsi une application de complexes. Nous
omettons de vérifier qu'elle coïncide avec l'application du lemme
\ref{lemma-relation-with-naive-cotangent-complex}.
\end{remark}

\begin{remark}
\label{remark-surjection}
Reprenons les notations de la remarque \ref{remark-make-map}. Les arguments
qui y sont donnés montrent que la différentielle
$$
H_2(\mathcal{C}_{B/A}, \mathcal{J}/\mathcal{J}^2)
\longrightarrow
H_0(\mathcal{C}_{B/A}, \text{Tor}_1^\mathcal{O}(\mathcal{J}, \underline{B}))
$$
de la suite spectrale est un isomorphisme. Notons $\mathcal{C}'_{B/A}$ la
sous-catégorie pleine de $\mathcal{C}_{B/A}$ formée des homomorphismes
surjectifs $P \to B$. L'identification du complexe cotangent avec le complexe
cotangent naïf (lemme \ref{lemma-relation-with-naive-cotangent-complex}) montre
que nous avons une suite exacte de faisceaux
$$
0 \to \underline{H_1(L_{B/A})} \to
\mathcal{J}/\mathcal{J}^2 \xrightarrow{\text{d}} \Omega \to
\underline{H_2(L_{B/A})} \to 0
$$
sur $\mathcal{C}'_{B/A}$. Il s'ensuit que $\Ker(d)$ et $\Coker(d)$ sur la
catégorie tout entière $\mathcal{C}_{B/A}$ ont leurs groupes d'homologie
supérieure nuls, puisque ceux-ci sont calculés par les groupes d'homologie de
groupes abéliens simpliciaux constants d'après le lemme
\ref{lemma-identify-pi-shriek}. Nous en concluons que
$$
H_n(\mathcal{C}_{B/A}, \mathcal{J}/\mathcal{J}^2) \to H_n(L_{B/A})
$$
est un isomorphisme pour tout $n \geq 2$. Avec la remarque précédente, ceci
donne la formule
$H_2(L_{B/A}) =
H_0(\mathcal{C}_{B/A}, \text{Tor}_1^\mathcal{O}(\mathcal{J}, \underline{B}))$.
\end{remark}





\section{Une suite spectrale de Quillen}
\label{section-spectral-sequence}

\noindent
Dans cette section, nous étudions une suite spectrale qui relie le produit
tensoriel dérivé au complexe cotangent.

\begin{lemma}
\label{lemma-vanishing-symmetric-powers}
Adoptons les notations et les hypothèses de Cohomologie sur les sites, exemple
\ref{sites-cohomology-example-category-to-point}. Supposons que $\mathcal{C}$
possède un objet cosimplicial comme dans Cohomologie sur les sites, lemme
\ref{sites-cohomology-lemma-compute-by-cosimplicial-resolution}. Soit
$\mathcal{F}$ un $\underline{B}$-module plat tel que
$H_0(\mathcal{C}, \mathcal{F}) = 0$. Alors
$H_l(\mathcal{C}, \text{Sym}_{\underline{B}}^k(\mathcal{F})) = 0$ pour
$l < k$.
\end{lemma}

\begin{proof}
Nous omettons l'indice ${}_{\underline{B}}$ dans les produits tensoriels, les
puissances extérieures et les puissances symétriques. Démontrons le lemme par
récurrence sur $k$. Les cas $k = 0, 1$ résultent des hypothèses. Si $k > 1$,
considérons le complexe exact
$$
\ldots \to
\wedge^2\mathcal{F} \otimes \text{Sym}^{k - 2}\mathcal{F} \to
\mathcal{F} \otimes \text{Sym}^{k - 1}\mathcal{F} \to
\text{Sym}^k\mathcal{F} \to 0
$$
dont les différentielles sont celles du complexe de Koszul. Considéré comme
une résolution de $\text{Sym}^k\mathcal{F}$, ce complexe donne une suite
spectrale du premier quadrant
$$
E_1^{p, q} =
H_p(\mathcal{C},
\wedge^{q + 1}\mathcal{F} \otimes \text{Sym}^{k - q - 1}\mathcal{F})
\Rightarrow
H_{p + q}(\mathcal{C}, \text{Sym}^k(\mathcal{F}))
$$
D'après Cohomologie sur les sites, lemme
\ref{sites-cohomology-lemma-eilenberg-zilber}, nous avons
$$
L\pi_!(\wedge^{q + 1}\mathcal{F} \otimes \text{Sym}^{k - q - 1}\mathcal{F}) =
L\pi_!(\wedge^{q + 1}\mathcal{F}) \otimes_B^\mathbf{L}
L\pi_!(\text{Sym}^{k - q - 1}\mathcal{F}))
$$
Il résulte de la construction des produits tensoriels dérivés que l'hypothèse
de récurrence, jointe à l'annulation de
$H_0(\mathcal{C}, \wedge^{q + 1}(\mathcal{F})) = 0$, donnera le résultat voulu.
Cette annulation est vraie parce que $\wedge^{q + 1}(\mathcal{F})$ est un
quotient de $\mathcal{F}^{\otimes q + 1}$ et que
$H_0(\mathcal{C}, \mathcal{F}^{\otimes q + 1})$ est un quotient de
$H_0(\mathcal{C}, \mathcal{F})^{\otimes q + 1}$, qui est nul.
\end{proof}

\begin{remark}
\label{remark-first-homology-symmetric-power}
Dans la situation du lemme \ref{lemma-vanishing-symmetric-powers}, on peut
montrer que
$H_k(\mathcal{C}, \text{Sym}^k(\mathcal{F})) =
\wedge^k_B(H_1(\mathcal{C}, \mathcal{F}))$.
En effet, on peut déduire de la démonstration que
$H_k(\mathcal{C}, \text{Sym}^k(\mathcal{F}))$ est le module des coinvariants
sous $S_k$ de
$$
H^{-k}(L\pi_!(\mathcal{F}) \otimes_B^\mathbf{L}
L\pi_!(\mathcal{F}) \otimes_B^\mathbf{L}
\ldots \otimes_B^\mathbf{L} L\pi_!(\mathcal{F})) =
H_1(\mathcal{C}, \mathcal{F})^{\otimes k}
$$
Notre assertion est donc que cette action est l'action usuelle de $S_k$ sur le
produit tensoriel, multipliée par le caractère signature. Pour le démontrer,
il faut examiner les conventions de signes dans la définition du complexe
total associé à un multicomplexe. Nous omettons la vérification.
\end{remark}

\begin{lemma}
\label{lemma-map-tors-zero}
Soit $A$ un anneau. Soit $P = A[E]$ un anneau de polynômes. Posons
$I = (e; e \in E) \subset P$. Les applications
$\text{Tor}_i^P(A, I^{n + 1}) \to \text{Tor}_i^P(A, I^n)$
sont nulles pour tous $i$ et $n$.
\end{lemma}

\begin{proof}
Notons $x_e \in P$ la variable correspondant à $e \in E$. Le complexe de
Koszul $K_\bullet$ associé aux $x_e$ fournit une résolution libre de $A$ sur
$P$. Ici, $K_i$ a pour base les produits extérieurs
$e_1 \wedge \ldots \wedge e_i$, $e_1, \ldots, e_i \in E$, et $d(e) = x_e$.
Ainsi, $K_\bullet \otimes_P I^n = I^nK_\bullet$ calcule
$\text{Tor}_i^P(A, I^n)$. Remarquons que tout est gradué, avec
$\deg(x_e) = 1$, $\deg(e) = 1$ et $\deg(a) = 0$ pour $a \in A$. Soit
$\xi \in I^{n + 1}K_i$ un cocycle homogène de degré $m$. On a
$m \geq i + 1 + n$. Il existe alors $\eta \in K_{i + 1}$ tel que
$\xi = \text{d}\eta$, puisque $K_\bullet$ est exact en degrés $ > 0$. (Le cas
$i = 0$ est laissé au lecteur.) Or $\deg(\eta) = m \geq i + 1 + n$. En
écrivant $\eta$ dans la base, on voit donc que ses coordonnées appartiennent à
$I^n$. Ainsi, $\xi$ s'envoie sur zéro dans l'homologie de $I^nK_\bullet$,
comme voulu.
\end{proof}

\begin{theorem}[Suite spectrale de Quillen]
\label{theorem-quillen-spectral-sequence}
Soit $A \to B$ un homomorphisme d'anneaux surjectif. Considérons le faisceau
de $\underline{B}$-modules
$\Omega = \Omega_{\mathcal{O}/A} \otimes_\mathcal{O} \underline{B}$ sur
$\mathcal{C}_{B/A}$ ; voir la section \ref{section-compute-L-pi-shriek}. Il
existe alors une suite spectrale dont la page $E_1$ est
$$
E_1^{p, q} =
H_{- p - q}(\mathcal{C}_{B/A}, \text{Sym}^p_{\underline{B}}(\Omega))
\Rightarrow \text{Tor}^A_{- p - q}(B, B)
$$
et dans laquelle $d_r$ est de bidegré $(r, -r + 1)$. De plus,
$H_i(\mathcal{C}_{B/A}, \text{Sym}^k_{\underline{B}}(\Omega)) = 0$ pour
$i < k$.
\end{theorem}

\begin{proof}
Soit $I \subset A$ le noyau de $A \to B$. Soit
$\mathcal{J} \subset \mathcal{O}$ le noyau de
$\mathcal{O} \to \underline{B}$. Alors
$I\mathcal{O} \subset \mathcal{J}$. Posons
$\mathcal{K} = \mathcal{J}/I\mathcal{O}$ et
$\overline{\mathcal{O}} = \mathcal{O}/I\mathcal{O}$.

\medskip\noindent
Pour tout objet $U = (P \to B)$ de $\mathcal{C}_{B/A}$, nous pouvons choisir
un isomorphisme $P = A[E]$ tel que l'homomorphisme $P \to B$ envoie chaque
$e \in E$ sur zéro. Alors
$J = \mathcal{J}(U) \subset P = \mathcal{O}(U)$
est égal à $J = IP + (e; e \in E)$. De plus,
$\overline{\mathcal{O}}(U) = B[E]$ et $K = \mathcal{K}(U) = (e; e \in E)$
est l'idéal engendré par les variables dans l'anneau de polynômes $B[E]$. En
particulier, il est clair que
$$
K/K^2 \xrightarrow{\text{d}} \Omega_{P/A} \otimes_P B
$$
est une bijection. Autrement dit, $\Omega = \mathcal{K}/\mathcal{K}^2$ et
$\text{Sym}_B^k(\Omega) = \mathcal{K}^k/\mathcal{K}^{k + 1}$. Remarquons que
$\pi_!(\Omega) = \Omega_{B/A} = 0$ (lemme \ref{lemma-identify-H0}), puisque
$A \to B$ est surjectif (Algèbre, lemme
\ref{algebra-lemma-trivial-differential-surjective}). D'après le lemme
\ref{lemma-vanishing-symmetric-powers}, nous en concluons que
$$
H_i(\mathcal{C}_{B/A}, \mathcal{K}^k/\mathcal{K}^{k + 1}) =
H_i(\mathcal{C}_{B/A}, \text{Sym}^k_{\underline{B}}(\Omega)) = 0
$$
pour $i < k$. Cela démontre la dernière assertion du théorème.

\medskip\noindent
Pour aborder le théorème, remarquons que
$$
B \otimes_A^\mathbf{L} B = L\pi_!(\mathcal{O}) \otimes_A^\mathbf{L} B =
L\pi_!(\mathcal{O} \otimes_{\underline{A}}^\mathbf{L} \underline{B}) =
L\pi_!(\overline{\mathcal{O}})
$$
La première égalité résulte du lemme \ref{lemma-apply-O-B-comparison}, la
deuxième de Cohomologie sur les sites, lemme
\ref{sites-cohomology-lemma-change-of-rings}, et la troisième de la platitude
de $\mathcal{O}$ sur $\underline{A}$. Le faisceau
$\overline{\mathcal{O}}$ possède une filtration
$$
\ldots \subset
\mathcal{K}^3 \subset
\mathcal{K}^2 \subset
\mathcal{K} \subset
\overline{\mathcal{O}}
$$
Celle-ci induit une filtration $F$ sur un complexe $C$ représentant
$L\pi_!(\overline{\mathcal{O}})$, où $F^pC$ représente
$L\pi_!(\mathcal{K}^p)$ (nous omettons la construction de $C$ et de $F$).
Considérons la suite spectrale de Homologie, section
\ref{homology-section-filtered-complex}, associée à $(C, F)$. Sa page $E_1$
est
$$
E_1^{p, q} = H_{- p - q}(\mathcal{C}_{B/A}, \mathcal{K}^p/\mathcal{K}^{p + 1})
\quad\Rightarrow\quad
H_{- p - q}(\mathcal{C}_{B/A}, \overline{\mathcal{O}}) = 
\text{Tor}_{- p - q}^A(B, B)
$$
et ses différentielles sont
$E_r^{p, q} \to E_r^{p + r, q - r + 1}$. Pour montrer la convergence, nous
allons établir que, pour tout $k$, il existe un $c$ tel que
$H_i(\mathcal{C}_{B/A}, \mathcal{K}^n) = 0$
pour $i < k$ et $n > c$\footnote{A posteriori, on pourra en déduire
l'annulation « correcte » $H_i(\mathcal{C}_{B/A}, \mathcal{K}^n) = 0$ pour
$i < n$.}.

\medskip\noindent
Étant donné $k \geq 0$, posons $c = k^2$. Nous affirmons que
$$
H_i(\mathcal{C}_{B/A}, \mathcal{K}^{n + c}) \to
H_i(\mathcal{C}_{B/A}, \mathcal{K}^n)
$$
est nulle pour $i < k$ et tout $n \geq 0$. Remarquons que
$\mathcal{K}^n/\mathcal{K}^{n + c}$ possède une filtration finie dont les
quotients successifs $\mathcal{K}^m/\mathcal{K}^{m + 1}$,
$n \leq m < n + c$, vérifient
$H_i(\mathcal{C}_{B/A}, \mathcal{K}^m/\mathcal{K}^{m + 1}) = 0$ pour $i < n$
(voir ci-dessus). L'assertion implique donc
$H_i(\mathcal{C}_{B/A}, \mathcal{K}^{n + c}) = 0$ pour $i < k$ et tout
$n \geq k$, ce qu'il fallait montrer.

\medskip\noindent
Démonstration de l'assertion. Rappelons que, pour tout $\mathcal{O}$-module
$\mathcal{F}$, l'application
$\mathcal{F} \to \mathcal{F} \otimes_\mathcal{O}^\mathbf{L} B$ induit un
isomorphisme après application de $L\pi_!$ ; voir le lemme
\ref{lemma-O-homology-B-homology}. Considérons l'application
$$
\mathcal{K}^{n + k} \otimes_\mathcal{O}^\mathbf{L} B
\longrightarrow
\mathcal{K}^n \otimes_\mathcal{O}^\mathbf{L} B
$$
Nous affirmons que cette application induit l'application nulle sur les
faisceaux de cohomologie en degrés $0, -1, \ldots, - k + 1$. Si cette seconde
assertion est vraie, la composée de $k$ applications
$$
\mathcal{K}^{n + c} \otimes_\mathcal{O}^\mathbf{L} B
\longrightarrow
\mathcal{K}^n \otimes_\mathcal{O}^\mathbf{L} B
$$
se factorise par
$\tau_{\leq -k}\mathcal{K}^n \otimes_\mathcal{O}^\mathbf{L} B$ et induit donc
zéro sur $H_i(\mathcal{C}_{B/A}, -) = L_i\pi_!( - )$ pour $i < k$ ; voir
Catégories dérivées, lemme
\ref{derived-lemma-trick-vanishing-composition}. D'après la remarque
précédente, il en va de même de
$H_i(\mathcal{C}_{B/A}, \mathcal{K}^{n + c}) \to
H_i(\mathcal{C}_{B/A}, \mathcal{K}^n)$
ce qui démontre la première assertion.

\medskip\noindent
Démonstration de la seconde assertion. L'énoncé est local ; nous pouvons donc
travailler au-dessus d'un objet $U = (P \to B)$ comme ci-dessus. Il faut
montrer que les applications
$$
\text{Tor}_i^P(B, K^{n + k}) \to \text{Tor}_i^P(B, K^n)
$$
sont nulles pour $i < k$. Il existe une suite spectrale
$$
\text{Tor}_a^P(P/IP, \text{Tor}_b^{P/IP}(B, K^n))
\Rightarrow
\text{Tor}_{a + b}^P(B, K^n),
$$
voir Compléments sur l'algèbre, exemple
\ref{more-algebra-example-tor-change-rings}. Il suffit donc de montrer que les
applications
$$
\text{Tor}_i^{P/IP}(B, K^{n + 1}) \to \text{Tor}_i^{P/IP}(B, K^n)
$$
sont nulles pour tout $i$. C'est le lemme \ref{lemma-map-tors-zero}.
\end{proof}

\begin{remark}
\label{remark-elucidate-ss}
Dans la situation du théorème \ref{theorem-quillen-spectral-sequence}, posons
$I = \Ker(A \to B)$. Alors
$H^{-1}(L_{B/A}) = H_1(\mathcal{C}_{B/A}, \Omega) = I/I^2$, d'après le
lemme \ref{lemma-surjection}.
Ainsi,
$H_k(\mathcal{C}_{B/A}, \text{Sym}^k(\Omega)) = \wedge^k_B(I/I^2)$ d'après la
remarque \ref{remark-first-homology-symmetric-power}. La page $E_1$ se présente
donc sous la forme
$$
\begin{matrix}
B \\
0 \\
0 & I/I^2 \\
0 & H^{-2}(L_{B/A}) \\
0 & H^{-3}(L_{B/A}) & \wedge^2(I/I^2) \\
0 & H^{-4}(L_{B/A}) & H_3(\mathcal{C}_{B/A}, \text{Sym}^2(\Omega)) \\
0 & H^{-5}(L_{B/A}) & H_4(\mathcal{C}_{B/A}, \text{Sym}^2(\Omega)) &
\wedge^3(I/I^2)
\end{matrix}
$$
avec une différentielle horizontale. Nous obtenons ainsi des morphismes de bord
$\text{Tor}_i^A(B, B) \to H^{-i}(L_{B/A})$, $i > 0$, et
$\wedge^i_B(I/I^2) \to \text{Tor}_i^A(B, B)$. Enfin,
$\text{Tor}_1^A(B, B) = I/I^2$ et il existe une suite exacte à cinq termes
$$
\text{Tor}_3^A(B, B) \to H^{-3}(L_{B/A}) \to \wedge^2_B(I/I^2) \to
\text{Tor}_2^A(B, B) \to H^{-2}(L_{B/A}) \to 0
$$
formée des termes de bas degré.
\end{remark}

\begin{remark}
\label{remark-elucidate-degree-two}
Soit $A \to B$ un homomorphisme d'anneaux. Soit $P_\bullet$ une résolution de
$B$ sur $A$ (remarque \ref{remark-resolution}). Posons
$J_n = \Ker(P_n \to B)$. Remarquons que
$$
\text{Tor}_2^{P_n}(B, B) = 
\text{Tor}_1^{P_n}(J_n, B) =
\Ker(J_n \otimes_{P_n} J_n \to J_n^2).
$$
Ainsi, $H_2(L_{B/A})$ est canoniquement égal à
$$
\Coker(\text{Tor}_2^{P_1}(B, B) \to \text{Tor}_2^{P_0}(B, B))
$$
d'après la remarque \ref{remark-surjection}. Pour expliciter ceci, choisissons
$P_2$, $P_1$, $P_0$ comme dans l'exemple
\ref{example-resolution-length-two}. Nous affirmons que
$$
\text{Tor}_2^{P_1}(B, B) =
\wedge^2(\bigoplus\nolimits_{t \in T} B)\ \oplus
\ \bigoplus\nolimits_{t \in T} J_0\ \oplus
\ \text{Tor}_2^{P_0}(B, B)
$$
Plus précisément, les éléments de base $x_t \wedge x_{t'}$ du premier facteur
direct correspondent à l'élément
$x_t \otimes x_{t'} - x_{t'} \otimes x_t$ de $J_1 \otimes_{P_1} J_1$. Pour
$f \in J_0$, l'élément $x_t \otimes f$ du deuxième facteur direct correspond
à l'élément $x_t \otimes s_0(f) - s_0(f) \otimes x_t$ de
$J_1 \otimes_{P_1} J_1$. Enfin, l'application
$\text{Tor}_2^{P_0}(B, B) \to \text{Tor}_2^{P_1}(B, B)$ est donnée par $s_0$.
L'application
$d_0 - d_1 : \text{Tor}_2^{P_1}(B, B) \to \text{Tor}_2^{P_0}(B, B)$
est nulle sur le dernier facteur direct, envoie $x_t \otimes f$ sur
$f \otimes f_t - f_t \otimes f$, et $x_t \wedge x_{t'}$ sur
$f_t \otimes f_{t'} - f_{t'} \otimes f_t$. En définitive, nous en concluons
qu'il existe une suite exacte
$$
\wedge^2_B(J_0/J_0^2) \to \text{Tor}_2^{P_0}(B, B) \to H^{-2}(L_{B/A}) \to 0
$$
Nous obtenons ainsi une démonstration directe d'une conséquence de la suite
spectrale de Quillen examinée dans la remarque \ref{remark-elucidate-ss}.
\end{remark}






\section{Comparaison avec Lichtenbaum-Schlessinger}
\label{section-compare-higher}

\noindent
Soit $A \to B$ un homomorphisme d'anneaux. L'article
\cite{Lichtenbaum-Schlessinger} donne une détermination assez explicite de
$\tau_{\geq -2}L_{B/A}$, souvent utilisée pour calculer les espaces de
déformations verselles de singularités. En voici la construction. Choisissons
une algèbre polynomiale $P$ sur $A$ et une surjection $P \to B$ de noyau $I$.
Choisissons des générateurs $f_t$, $t \in T$, de $I$, ce qui induit une
surjection $F = \bigoplus_{t \in T} P \to I$, où $F$ est un $P$-module libre.
Soit $Rel \subset F$ le noyau de $F \to I$ ; autrement dit, $Rel$ est
l'ensemble des relations entre les $f_t$. Soit $TrivRel \subset Rel$ le
sous-module des relations triviales, c'est-à-dire le sous-module de $Rel$
engendré par les éléments
$(\ldots, f_{t'}, 0, \ldots, 0, -f_t, 0, \ldots)$. Considérons le complexe de
$B$-modules
\begin{equation}
\label{equation-lichtenbaum-schlessinger}
Rel/TrivRel \longrightarrow
F \otimes_P B \longrightarrow
\Omega_{P/A} \otimes_P B
\end{equation}
où le dernier terme est placé en degré $0$. La première application est
l'application évidente et la seconde envoie l'élément de base correspondant à
$t \in T$ sur $\text{d}f_t \otimes 1$.

\begin{definition}
\label{definition-biderivation}
Soit $A \to B$ un homomorphisme d'anneaux. Soit $M$ un $(B, B)$-bimodule sur
$A$. Une {\it $A$-bidérivation} est une application $A$-linéaire
$\lambda : B \to M$ telle que $\lambda(xy) = x\lambda(y) + \lambda(x)y$.
\end{definition}

\noindent
Pour une algèbre polynomiale, les bidérivations sont faciles à décrire.

\begin{lemma}
\label{lemma-polynomial-ring-unique}
Soit $P = A[S]$ un anneau de polynômes sur $A$. Soit $M$ un
$(P, P)$-bimodule sur $A$. Étant donnés des $m_s \in M$, pour $s \in S$, il
existe une unique $A$-bidérivation $\lambda : P \to M$ qui envoie $s$ sur
$m_s$ pour $s \in S$.
\end{lemma}

\begin{proof}
Posons
$$
\lambda(s_1 \ldots s_t) =
\sum s_1 \ldots s_{i - 1} m_{s_i} s_{i + 1} \ldots s_t
$$
dans $M$. Le prolongement par $A$-linéarité est une bidérivation.
\end{proof}

\noindent
Voici l'énoncé de comparaison. Le lecteur pourra également consulter
\cite[page 206, Proposition 12]{Andre-Homologie} ou l'article \cite{Doncel},
qui prolonge le complexe (\ref{equation-lichtenbaum-schlessinger}) d'un terme
et étend la comparaison à $\tau_{\geq -3}$.

\begin{lemma}
\label{lemma-compare-higher}
Dans la situation précédente, notons $L$ le complexe
(\ref{equation-lichtenbaum-schlessinger}). Il existe dans $D(B)$ une
application canonique $L_{B/A} \to L$ qui induit un isomorphisme
$\tau_{\geq -2}L_{B/A} \to L$ dans $D(B)$.
\end{lemma}

\begin{proof}
Soit $P_\bullet \to B$ une résolution de $B$ sur $A$ (remarque
\ref{remark-resolution}). Nous identifierons $L_{B/A}$ à
$\Omega_{P_\bullet/A} \otimes B$. Pour construire l'application, effectuons
quelques choix.

\medskip\noindent
Choisissons un homomorphisme de $A$-algèbres $\psi : P_0 \to P$ compatible
avec les homomorphismes donnés $P_0 \to B$ et $P \to B$.

\medskip\noindent
Écrivons $P_1 = A[S]$ pour un certain ensemble $S$. Pour $s \in S$, nous
pouvons écrire
$$
\psi(d_0(s) - d_1(s)) = \sum p_{s, t} f_t
$$
pour certains $p_{s, t} \in P$. Considérons
$F = \bigoplus_{t \in T} P$ comme un $(P_1, P_1)$-bimodule au moyen des
homomorphismes $(\psi \circ d_0, \psi \circ d_1)$. D'après le lemme
\ref{lemma-polynomial-ring-unique}, il existe une unique $A$-bidérivation
$\lambda : P_1 \to F$ envoyant $s$ sur le vecteur de coordonnées $p_{s, t}$.
Par construction, la composée
$$
P_1 \longrightarrow F \longrightarrow P
$$
envoie $f \in P_1$ sur $\psi(d_0(f) - d_1(f))$, car l'application
$f \mapsto \psi(d_0(f) - d_1(f))$ est une $A$-bidérivation qui coïncide avec la
composée sur les générateurs.

\medskip\noindent
Pour $g \in P_2$, nous affirmons que
$\lambda(d_0(g) - d_1(g) + d_2(g))$ appartient à $Rel$. En effet, d'après la
dernière observation du paragraphe précédent, l'image de
$\lambda(d_0(g) - d_1(g) + d_2(g))$ dans $P$ est
$$
\psi((d_0 - d_1)(d_0(g) - d_1(g) + d_2(g)))
$$
qui est nulle d'après Simplicial, section
\ref{simplicial-section-complexes}.

\medskip\noindent
Le choix de $\psi$ détermine une application
$$
\text{d}\psi \otimes 1 :
\Omega_{P_0/A} \otimes B
\longrightarrow
\Omega_{P/A} \otimes B
$$
La composée de $\lambda$ avec l'application $F \to F \otimes B$ est une
$A$-dérivation usuelle, puisque les deux structures de $P_1$-module sur
$F \otimes B$ coïncident. Ainsi, $\lambda$ détermine une application
$$
\overline{\lambda} :
\Omega_{P_1/A} \otimes B
\longrightarrow
F \otimes B
$$
Enfin, nous obtenons une application $B$-linéaire
$$
q :
\Omega_{P_2/A} \otimes B
\longrightarrow
Rel/TrivRel
$$
en envoyant $\text{d}g$ sur la classe de
$\lambda(d_0(g) - d_1(g) + d_2(g))$ dans le quotient.

\medskip\noindent
Le diagramme
$$
\xymatrix{
\Omega_{P_3/A} \otimes B \ar[r] \ar[d] &
\Omega_{P_2/A} \otimes B \ar[r] \ar[d]_q &
\Omega_{P_1/A} \otimes B \ar[r] \ar[d]_{\overline{\lambda}} &
\Omega_{P_0/A} \otimes B \ar[d]_{\text{d}\psi \otimes 1} \\
0 \ar[r] &
Rel/TrivRel \ar[r] &
F \otimes B \ar[r] &
\Omega_{P/A} \otimes B
}
$$
est commutatif (calcul omis), et nous obtenons l'application du lemme. La
remarque \ref{remark-explicit-comparison-map} et le lemme
\ref{lemma-relation-with-naive-cotangent-complex} montrent que cette
application induit des isomorphismes $H_1(L_{B/A}) \to H_1(L)$ et
$H_0(L_{B/A}) \to H_0(L)$.

\medskip\noindent
Il reste à voir que notre application $L_{B/A} \to L$ induit un isomorphisme
$H_2(L_{B/A}) \to H_2(L)$. Choisissons une résolution de $B$ sur $A$ telle que
$P_0 = P = A[u_i]$, puis $P_1$ et $P_2$ comme dans l'exemple
\ref{example-resolution-length-two}. Dans la remarque
\ref{remark-elucidate-degree-two}, nous avons construit une suite exacte
$$
\wedge^2_B(J_0/J_0^2) \to \text{Tor}_2^{P_0}(B, B) \to H^{-2}(L_{B/A}) \to 0
$$
où $P_0 = P$ et $J_0 = \Ker(P \to B) = I$. En calculant le groupe de Tor à
l'aide des suites exactes courtes $0 \to I \to P \to B \to 0$ et
$0 \to Rel \to F \to I \to 0$, nous trouvons que
$\text{Tor}_2^P(B, B) = \Ker(Rel \otimes B \to F \otimes B)$.
Sous cette identification, l'image de l'application
$\wedge^2_B(I/I^2) \to \text{Tor}_2^P(B, B)$ est exactement l'image de
$TrivRel \otimes B$. On voit donc que $H_2(L_{B/A}) \cong H_2(L)$.

\medskip\noindent
Enfin, il faut vérifier que notre application $L_{B/A} \to L$ induit bien cet
isomorphisme. Nous utiliserons sans autre mention les notations et les résultats
de l'exemple \ref{example-resolution-length-two} et des remarques
\ref{remark-elucidate-degree-two} et \ref{remark-surjection}. Choisissons un
élément $\xi$ de
$\text{Tor}_2^{P_0}(B, B) = \Ker(I \otimes_P I \to I^2)$.
Écrivons $\xi = \sum h_{t', t}f_{t'} \otimes f_t$ pour certains
$h_{t', t} \in P$. En suivant les suites exactes précédentes, on trouve que
$\xi$ correspond à l'image dans $Rel \otimes B$ de l'élément
$r \in Rel \subset F = \bigoplus_{t \in T} P$ dont la coordonnée d'indice $t$
est $r_t = \sum_{t' \in T} h_{t', t}f_{t'}$. D'autre part, $\xi$ correspond à
l'élément de $H_2(L_{B/A}) = H_2(\Omega)$ qui est l'image, par
$\text{d} : H_2(\mathcal{J}/\mathcal{J}^2) \to H_2(\Omega)$, du bord de $\xi$
pour la $2$-extension
$$
0 \to
\text{Tor}_2^\mathcal{O}(\underline{B}, \underline{B})
\to 
\mathcal{J} \otimes_\mathcal{O} \mathcal{J} \to \mathcal{J}
\to
\mathcal{J}/\mathcal{J}^2 \to 0
$$
Calculons les transgressions successives de notre élément. Tout d'abord,
$$
\xi = (d_0 - d_1)(- \sum s_0(h_{t', t} f_{t'}) \otimes x_t)
$$
puis
$$
\sum s_0(h_{t', t} f_{t'}) x_t = d_0(v_r) - d_1(v_r) + d_2(v_r)
$$
par notre choix des variables $v$ dans l'exemple
\ref{example-resolution-length-two}. Nous pouvons choisir l'application
$\lambda$ ci-dessus de telle sorte que
$\lambda(u_i) = 0$ et $\lambda(x_t) = - e_t$, où $e_t \in F$ désigne le
vecteur de base correspondant à $t \in T$. La
construction de l'application $q$ ci-dessus envoie donc $\text{d}v_r$ sur
$$
\lambda(\sum s_0(h_{t', t} f_{t'}) x_t) =
\sum\nolimits_t \left(\sum\nolimits_{t'} h_{t', t}f_{t'}\right) e_t
$$
qui coïncide avec l'image de $\xi$ dans $Rel \otimes B$ (les deux signes moins
apparus ci-dessus se compensent). Cette coïncidence achève la démonstration.
\end{proof}

\begin{remark}[Fonctorialité du complexe de Lichtenbaum-Schlessinger]
\label{remark-functoriality-lichtenbaum-schlessinger}
Considérons un carré commutatif
$$
\xymatrix{
A' \ar[r] & B' \\
A \ar[u] \ar[r] & B \ar[u]
}
$$
d'homomorphismes d'anneaux. Choisissons une factorisation
$$
\xymatrix{
A' \ar[r] & P' \ar[r] & B' \\
A \ar[u] \ar[r] & P \ar[u] \ar[r] & B \ar[u]
}
$$
où $P$ est une algèbre polynomiale sur $A$ et $P'$ une algèbre polynomiale sur
$A'$. Choisissons des générateurs $f_t$, $t \in T$, de $\Ker(P \to B)$. Pour
$t \in T$, notons $f'_t$ l'image de $f_t$ dans $P'$. Choisissons des
$f'_s \in P'$ de sorte que les éléments $f'_t$, pour
$t \in T' = T \amalg S$, engendrent le noyau de $P' \to B'$. Posons
$F = \bigoplus_{t \in T} P$ et $F' = \bigoplus_{t' \in T'} P'$. Soient
$Rel = \Ker(F \to P)$ et $Rel' = \Ker(F' \to P')$, où les applications sont
données sur les coordonnées par la multiplication par $f_t$, resp.\ $f'_t$.
Enfin, soit $TrivRel$, resp.\ $TrivRel'$, le sous-module de $Rel$, resp.\
$Rel'$, engendré par les éléments
$(\ldots, f_{t'}, 0, \ldots, 0, -f_t, 0, \ldots)$
pour $t, t' \in T$, resp.\ $T'$. Ces choix donnent un diagramme commutatif
canonique
$$
\xymatrix{
L' : &
Rel'/TrivRel' \ar[r] &
F' \otimes_{P'} B' \ar[r] &
\Omega_{P'/A'} \otimes_{P'} B' \\
L : \ar[u] &
Rel/TrivRel \ar[r] \ar[u] &
F \otimes_P B \ar[r] \ar[u] &
\Omega_{P/A} \otimes_P B \ar[u]
}
$$
De plus, en suivant les choix effectués dans la démonstration du lemme
\ref{lemma-compare-higher}, le lecteur constate que l'on obtient un diagramme
commutatif
$$
\xymatrix{
L_{B'/A'} \ar[r] & L' \\
L_{B/A} \ar[r] \ar[u] & L \ar[u]
}
$$
\end{remark}





\section{Le complexe cotangent d'une intersection complète locale}
\label{section-lci}

\noindent
Si $A \to B$ est un homomorphisme d'intersection complète locale, alors
$L_{B/A}$ est un complexe parfait. Le lemme suivant est la clef de la
démonstration.

\begin{lemma}
\label{lemma-special-case}
Soit $A = \mathbf{Z}[x_1, \ldots, x_n] \to B = \mathbf{Z}$ l'homomorphisme
d'anneaux qui envoie $x_i$ sur $0$ pour $i = 1, \ldots, n$. Soit
$I = (x_1, \ldots, x_n) \subset A$. Alors $L_{B/A}$ est quasi-isomorphe à
$I/I^2[1]$.
\end{lemma}

\begin{proof}
Il y a plusieurs façons de le démontrer. On peut, par exemple, construire
explicitement une résolution de $B$ sur $A$ et effectuer le calcul. Nous
utiliserons (\ref{equation-triangle}). Considérons en effet le triangle
distingué
$$
L_{\mathbf{Z}[x_1, \ldots, x_n]/\mathbf{Z}}
\otimes_{\mathbf{Z}[x_1, \ldots, x_n]} \mathbf{Z} \to
L_{\mathbf{Z}/\mathbf{Z}} \to
L_{\mathbf{Z}/\mathbf{Z}[x_1, \ldots, x_n]}\to
L_{\mathbf{Z}[x_1, \ldots, x_n]/\mathbf{Z}}
\otimes_{\mathbf{Z}[x_1, \ldots, x_n]} \mathbf{Z}[1]
$$
Le complexe $L_{\mathbf{Z}[x_1, \ldots, x_n]/\mathbf{Z}}$ est quasi-isomorphe à
$\Omega_{\mathbf{Z}[x_1, \ldots, x_n]/\mathbf{Z}}$ d'après le lemme
\ref{lemma-cotangent-complex-polynomial-algebra}. Le complexe
$L_{\mathbf{Z}/\mathbf{Z}}$ est nul dans $D(\mathbf{Z})$ d'après le lemme
\ref{lemma-when-zero}. On voit ainsi que $L_{B/A}$ n'a qu'un seul groupe de
cohomologie non nul, qui est celui décrit dans l'énoncé d'après le lemme
\ref{lemma-surjection}.
\end{proof}

\begin{lemma}
\label{lemma-mod-regular-sequence}
Soit $A \to B$ un homomorphisme d'anneaux surjectif dont le noyau $I$ est
engendré par une suite régulière au sens de Koszul (par exemple une suite
régulière). Alors $L_{B/A}$ est quasi-isomorphe à $I/I^2[1]$.
\end{lemma}

\begin{proof}
Soit $f_1, \ldots, f_r \in I$ une suite régulière au sens de Koszul qui
engendre $I$. Considérons l'homomorphisme d'anneaux
$\mathbf{Z}[x_1, \ldots, x_r] \to A$ qui envoie $x_i$ sur $f_i$. Comme
$x_1, \ldots, x_r$ est une suite régulière dans
$\mathbf{Z}[x_1, \ldots, x_r]$, le complexe de Koszul associé à
$x_1, \ldots, x_r$ est une résolution libre de
$\mathbf{Z} = \mathbf{Z}[x_1, \ldots, x_r]/(x_1, \ldots, x_r)$
sur $\mathbf{Z}[x_1, \ldots, x_r]$ (voir Compléments sur l'algèbre, lemme
\ref{more-algebra-lemma-regular-koszul-regular}). Ainsi, l'hypothèse que
$f_1, \ldots, f_r$ est régulière au sens de Koszul signifie exactement que
$B = A \otimes_{\mathbf{Z}[x_1, \ldots, x_r]}^\mathbf{L} \mathbf{Z}$.
Par conséquent,
$L_{B/A} = L_{\mathbf{Z}/\mathbf{Z}[x_1, \ldots, x_r]}
\otimes_\mathbf{Z}^\mathbf{L} B$ d'après les lemmes
\ref{lemma-flat-base-change-cotangent-complex} et \ref{lemma-special-case}.
\end{proof}

\begin{lemma}
\label{lemma-mod-Koszul-regular-ideal}
Soit $A \to B$ un homomorphisme d'anneaux surjectif dont le noyau $I$ est un
idéal de Koszul. Alors $L_{B/A}$ est quasi-isomorphe à $I/I^2[1]$.
\end{lemma}

\begin{proof}
Localement sur $\Spec(A)$, l'idéal $I$ est engendré par une suite régulière au
sens de Koszul ; voir Compléments sur l'algèbre, définition
\ref{more-algebra-definition-regular-ideal}. Le résultat découle donc du lemme
\ref{lemma-flat-base-change-cotangent-complex}.
\end{proof}

\begin{proposition}
\label{proposition-cotangent-complex-local-complete-intersection}
Soit $A \to B$ un homomorphisme d'intersection complète locale. Alors
$L_{B/A}$ est un complexe parfait d'amplitude de Tor contenue dans $[-1, 0]$.
\end{proposition}

\begin{proof}
Choisissons une surjection $P = A[x_1, \ldots, x_n] \to B$ de noyau $J$.
D'après le lemme \ref{lemma-relation-with-naive-cotangent-complex}, le
complexe $J/J^2 \to \bigoplus B\text{d}x_i$ est quasi-isomorphe à
$\tau_{\geq -1}L_{B/A}$. Remarquons que $J/J^2$ est projectif de type fini
(Compléments sur l'algèbre, lemme
\ref{more-algebra-lemma-quasi-regular-ideal-finite-projective}) ; ainsi,
$\tau_{\geq -1}L_{B/A}$ est un complexe parfait d'amplitude de Tor contenue
dans $[-1, 0]$. Il suffit donc de montrer que $H^i(L_{B/A}) = 0$ pour
$i \not \in [-1, 0]$. Cela résulte de (\ref{equation-triangle})
$$
L_{P/A} \otimes_P^\mathbf{L} B \to L_{B/A} \to L_{B/P} \to
L_{P/A} \otimes_P^\mathbf{L} B[1]
$$
et du lemme \ref{lemma-mod-Koszul-regular-ideal}, qui montrent que
$H^i(L_{B/P})$ est nul sauf si $i \in \{-1, 0\}$. (Pour le terme de gauche,
nous utilisons également le lemme
\ref{lemma-cotangent-complex-polynomial-algebra}.)
\end{proof}








\section{Produits tensoriels et complexe cotangent}
\label{section-tensor-product}

\noindent
Soit $R$ un anneau et soient $A$, $B$ des $R$-algèbres. Dans cette section,
nous étudions $L_{A \otimes_R B/R}$. L'essentiel des renseignements recherchés
est contenu dans le diagramme suivant
\begin{equation}
\label{equation-tensor-product}
\vcenter{
\xymatrix{
L_{A/R} \otimes_A^\mathbf{L} (A \otimes_R B) \ar[r] &
L_{A \otimes_R B/B} \ar[r] &
E \\
L_{A/R} \otimes_A^\mathbf{L} (A \otimes_R B) \ar[r] \ar@{=}[u] &
L_{A \otimes_R B/R} \ar[r] \ar[u] &
L_{A \otimes_R B/A} \ar[u] \\
 &
L_{B/R} \otimes_B^\mathbf{L} (A \otimes_R B) \ar[u] \ar@{=}[r] &
L_{B/R} \otimes_B^\mathbf{L} (A \otimes_R B) \ar[u]
}
}
\end{equation}
Explication : la ligne médiane est le triangle fondamental
(\ref{equation-triangle}) associé aux homomorphismes d'anneaux
$R \to A \to A \otimes_R B$. La colonne médiane est le triangle fondamental
(\ref{equation-triangle}) associé aux homomorphismes d'anneaux
$R \to B \to A \otimes_R B$. Ensuite, $E$ est un objet de
$D(A \otimes_R B)$ qui « complète » le coin supérieur droit, c'est-à-dire qui
fait de la ligne supérieure et de la colonne droite des triangles distingués.
Un tel $E$ existe d'après Catégories dérivées, proposition
\ref{derived-proposition-9}, appliquée au carré inférieur gauche (en plaçant
$0$ à la position manquante). Plus explicitement, nous pourrions par exemple
définir $E$ comme le cône (Catégories dérivées, définition
\ref{derived-definition-cone}) de l'application de complexes
$$
L_{A/R} \otimes_A^\mathbf{L} (A \otimes_R B) \oplus
L_{B/R} \otimes_B^\mathbf{L} (A \otimes_R B)
\longrightarrow
L_{A \otimes_R B/R}
$$
et obtenir les deux applications de but $E$ en appliquant TR3. Dans le cas
Tor-indépendant, l'objet $E$ est nul.

\begin{lemma}
\label{lemma-tensor-product-tor-independent}
Si $A$ et $B$ sont des $R$-algèbres Tor-indépendantes, alors l'objet $E$ de
(\ref{equation-tensor-product}) est nul. Dans ce cas, nous avons
$$
L_{A \otimes_R B/R} =
L_{A/R} \otimes_A^\mathbf{L} (A \otimes_R B) \oplus
L_{B/R} \otimes_B^\mathbf{L} (A \otimes_R B)
$$
qui est représenté par le complexe
$L_{A/R} \otimes_R B \oplus L_{B/R} \otimes_R A $
de $A \otimes_R B$-modules.
\end{lemma}

\begin{proof}
Les deux premières assertions résultent immédiatement du lemme
\ref{lemma-flat-base-change-cotangent-complex}. La dernière résulte du fait que
$L_{A/R}$ est un complexe de $A$-modules libres ; ainsi,
$L_{A/R} \otimes_A^\mathbf{L} (A \otimes_R B)$ est représenté par
$L_{A/R} \otimes_A (A \otimes_R B) = L_{A/R} \otimes_R B$
\end{proof}

\noindent
En général, nous pouvons dire ce qui suit de l'objet $E$.

\begin{lemma}
\label{lemma-tensor-product}
Soit $R$ un anneau et soient $A$, $B$ des $R$-algèbres. L'objet $E$ de
(\ref{equation-tensor-product}) vérifie
$$
H^i(E) =
\left\{
\begin{matrix}
0 & \text{si} & i \geq -1 \\
\text{Tor}_1^R(A, B) & \text{si} & i = -2
\end{matrix}
\right.
$$
\end{lemma}

\begin{proof}
Nous utilisons la description de $E$ comme le cône de
$L_{B/R} \otimes_B^\mathbf{L} (A \otimes_R B) \to L_{A \otimes_R B/A}$.
D'après le lemme \ref{lemma-compare-higher}, les troncations canoniques
$\tau_{\geq -2}L_{B/R}$ et $\tau_{\geq -2}L_{A \otimes_R B/A}$ sont calculées
par le complexe de Lichtenbaum-Schlessinger
(\ref{equation-lichtenbaum-schlessinger}). Ces isomorphismes sont compatibles
à la fonctorialité (remarque
\ref{remark-functoriality-lichtenbaum-schlessinger}). Nous travaillons donc
avec les complexes de Lichtenbaum-Schlessinger dans cette démonstration.

\medskip\noindent
Choisissons une algèbre polynomiale $P$ sur $R$ et une surjection $P \to B$.
Choisissons des générateurs $f_t \in P$, $t \in T$, du noyau de cette
surjection. Soit $Rel \subset F = \bigoplus_{t \in T} P$ le noyau de
l'application $F \to P$ qui envoie le vecteur de base correspondant à $t$ sur
$f_t$. Posons $P_A = A \otimes_R P$ et
$F_A = A \otimes_R F = P_A \otimes_P F$. Soit $Rel_A$ le noyau de
l'application $F_A \to P_A$. La suite exacte
$$
0 \to Rel \to F \to P \to B \to 0
$$
et les suites exactes courtes usuelles pour Tor donnent une suite exacte
$$
A \otimes_R Rel \to Rel_A \to \text{Tor}_1^R(A, B) \to 0
$$
Remarquons que $P_A \to A \otimes_R B$ est une surjection dont le noyau est
engendré par les éléments $1 \otimes f_t$ de $P_A$. Notons
$TrivRel_A \subset Rel_A$ le sous-$P_A$-module engendré par les éléments
$(\ldots, 1 \otimes f_{t'}, 0, \ldots,
0, - 1 \otimes f_t \otimes 1, 0, \ldots)$.
Comme $TrivRel \otimes_R A \to TrivRel_A$ est surjectif, nous obtenons une
suite exacte canonique
$$
A \otimes_R (Rel/TrivRel) \to Rel_A/TrivRel_A \to \text{Tor}_1^R(A, B) \to 0
$$
L'application entre les complexes de Lichtenbaum-Schlessinger est donnée par
le diagramme
$$
\xymatrix{
Rel_A/TrivRel_A \ar[r] &
F_A \otimes_{P_A} (A \otimes_R B) \ar[r] &
\Omega_{P_A/A \otimes_R B} \otimes_{P_A} (A \otimes_R B) \\
Rel/TrivRel \ar[r] \ar[u]_{-2} &
F \otimes_P B \ar[r] \ar[u]_{-1} &
\Omega_{P/A} \otimes_P B \ar[u]_0
}
$$
Remarquons que les applications verticales $-1$ et $0$ induisent un
isomorphisme après application du foncteur
$A \otimes_R - = P_A \otimes_P -$ à la source, et que l'application verticale
$-2$ donne exactement l'application dont le conoyau est le module de Tor
recherché, comme nous l'avons vu ci-dessus.
\end{proof}












\section{Déformations d'homomorphismes d'anneaux et complexe cotangent}
\label{section-deformations}

\noindent
Cette section prolonge Théorie des déformations, section
\ref{defos-section-deformations}, que nous invitons vivement le lecteur à lire
d'abord. Partons d'un homomorphisme d'anneaux surjectif $A' \to A$ dont le
noyau est un idéal $I$ de carré nul. Supposons de plus donnés un homomorphisme
d'anneaux $A \to B$, un $B$-module $N$ et une application de $A$-modules
$c : I \to N$. Nous cherchons ici si l'on peut remplacer le point
d'interrogation dans le diagramme suivant
\begin{equation}
\label{equation-to-solve}
\vcenter{
\xymatrix{
0 \ar[r] & N \ar[r] & {?} \ar[r] & B \ar[r] & 0 \\
0 \ar[r] & I \ar[u]^c \ar[r] & A' \ar[u] \ar[r] & A \ar[u] \ar[r] & 0
}
}
\end{equation}
et, si une solution existe, à quel point elle est unique. Plus précisément,
nous cherchons une surjection de $A'$-algèbres $B' \to B$ dont le noyau est
un idéal de carré nul identifié à $N$, et telle que $A' \to B'$ induise
l'application donnée $c$. Nous dirons que $B'$ est une {\it solution} de
(\ref{equation-to-solve}).

\begin{lemma}
\label{lemma-find-obstruction}
Dans la situation précédente :
\begin{enumerate}
\item il existe un élément canonique $\xi \in \Ext^2_B(L_{B/A}, N)$ dont
l'annulation est une condition nécessaire et suffisante pour l'existence d'une
solution de (\ref{equation-to-solve}) ;
\item s'il existe une solution, l'ensemble des classes d'isomorphisme de
solutions est un espace principal homogène sous $\Ext^1_B(L_{B/A}, N)$ ;
\item étant donnée une solution $B'$, l'ensemble des automorphismes de $B'$
qui s'insèrent dans (\ref{equation-to-solve}) est canoniquement isomorphe à
$\Ext^0_B(L_{B/A}, N)$.
\end{enumerate}
\end{lemma}

\begin{proof}
Au moyen des identifications $\NL_{B/A} = \tau_{\geq -1}L_{B/A}$ (lemme
\ref{lemma-relation-with-naive-cotangent-complex}) et
$H^0(L_{B/A}) = \Omega_{B/A}$ (lemme \ref{lemma-identify-H0}), nous avons déjà
établi (2) et (3) dans Théorie des déformations, lemmes
\ref{defos-lemma-huge-diagram} et \ref{defos-lemma-choices}.

\medskip\noindent
Démonstration de (1). En gros, l'assertion découle de la discussion de Théorie
des déformations, remarque \ref{defos-remark-parametrize-solutions}, en
remplaçant le complexe cotangent naïf par le complexe cotangent entier. Voici
une explication plus détaillée. D'après Théorie des déformations, lemme
\ref{defos-lemma-parametrize-solutions} et remarque
\ref{defos-remark-parametrize-solutions}, il existe un élément
$$
\xi' \in
\Ext^1_A(\NL_{A/A'}, N) =
\Ext^1_B(\NL_{A/A'} \otimes_A^\mathbf{L} B, N) =
\Ext^1_B(L_{A/A'} \otimes_A^\mathbf{L} B, N)
$$
(pour les égalités, voir Théorie des déformations, remarque
\ref{defos-remark-parametrize-solutions}, et utiliser
$\NL_{A'/A} = \tau_{\geq -1} L_{A'/A}$) tel qu'une solution existe si et
seulement si cet élément appartient à l'image de l'application
$$
\Ext^1_B(\NL_{B/A'}, N) = \Ext^1_B(L_{B/A'}, N)
\longrightarrow
\Ext^1_B(L_{A/A'} \otimes_A^\mathbf{L} B, N)
$$
Le triangle distingué (\ref{equation-triangle}) associé à
$A' \to A \to B$ donne naissance à une suite exacte longue
$$
\ldots \to
\Ext^1_B(L_{B/A'}, N) \to
\Ext^1_B(L_{A/A'} \otimes_A^\mathbf{L} B, N) \to
\Ext^2_B(L_{B/A}, N) \to \ldots
$$
Il suffit donc de prendre pour $\xi$ l'image de $\xi'$.
\end{proof}





\section{La classe d'Atiyah d'un module}
\label{section-atiyah}

\noindent
Soit $A \to B$ un homomorphisme d'anneaux. Soit $M$ un $B$-module. Soit
$P \to B$ un objet de $\mathcal{C}_{B/A}$ (section
\ref{section-compute-L-pi-shriek}). Considérons l'extension des parties
principales
$$
0 \to \Omega_{P/A} \otimes_P M \to P^1_{P/A}(M) \to M \to 0
$$
voir Algèbre, lemme \ref{algebra-lemma-sequence-of-principal-parts}. Cette
suite est fonctorielle en $P$ d'après Algèbre, remarque
\ref{algebra-remark-functoriality-principal-parts}. Nous obtenons ainsi une
suite exacte courte de faisceaux de $\mathcal{O}$-modules
$$
0 \to \Omega_{\mathcal{O}/\underline{A}} \otimes_\mathcal{O} \underline{M} \to
P^1_{\mathcal{O}/\underline{A}}(M) \to \underline{M} \to 0
$$
sur $\mathcal{C}_{B/A}$. Nous avons
$L\pi_!(\Omega_{\mathcal{O}/\underline{A}} \otimes_\mathcal{O} \underline{M})
= L_{B/A} \otimes_B M = L_{B/A} \otimes_B^\mathbf{L} M$
d'après le lemme \ref{lemma-pi-shriek-standard} et la platitude des termes de
$L_{B/A}$. De plus, $L\pi_!(\underline{M}) = M$ d'après le lemme
\ref{lemma-pi-lower-shriek-constant-sheaf}. On obtient donc un triangle
distingué
\begin{equation}
\label{equation-atiyah}
L_{B/A} \otimes_B^\mathbf{L} M \to
L\pi_!\left(P^1_{\mathcal{O}/\underline{A}}(M)\right) \to M
\to L_{B/A} \otimes_B^\mathbf{L} M [1]
\end{equation}
dans $D(B)$. Nous utilisons ici Cohomologie sur les sites, remarque
\ref{sites-cohomology-remark-O-homology-B-homology-general}, pour obtenir un
triangle distingué dans $D(B)$, et pas seulement dans $D(A)$.

\begin{definition}
\label{definition-atiyah-class}
Soit $A \to B$ un homomorphisme d'anneaux. Soit $M$ un $B$-module.
L'application $M \to L_{B/A} \otimes_B^\mathbf{L} M[1]$ de
(\ref{equation-atiyah}) est appelée la {\it classe d'Atiyah} de $M$.
\end{definition}




\section{Le complexe cotangent}
\label{section-cotangent-complex}

\noindent
Dans cette section, nous étudions le complexe cotangent d'un homomorphisme de
faisceaux d'anneaux sur un site. Dans les sections suivantes, nous
spécialiserons cette construction afin d'obtenir le complexe cotangent d'un
morphisme de topos annelés, d'espaces annelés, de schémas, d'espaces
algébriques, etc.

\medskip\noindent
Soit $\mathcal{C}$ un site et notons $\Sh(\mathcal{C})$ le topos associé. Soit
$\mathcal{A}$ un faisceau d'anneaux sur $\mathcal{C}$. Notons
$\mathcal{A}\textit{-Alg}$ la catégorie des $\mathcal{A}$-algèbres.
Considérons le couple de foncteurs adjoints $(U, V)$, où
$V : \mathcal{A}\textit{-Alg} \to \Sh(\mathcal{C})$ est le foncteur d'oubli et
où $U : \Sh(\mathcal{C}) \to \mathcal{A}\textit{-Alg}$ associe à un faisceau
d'ensembles $\mathcal{E}$ l'algèbre polynomiale
$\mathcal{A}[\mathcal{E}]$ en $\mathcal{E}$ sur $\mathcal{A}$. Soit
$X_\bullet$ l'objet simplicial de
$\text{Fun}(\mathcal{A}\textit{-Alg}, \mathcal{A}\textit{-Alg})$ construit dans
Simplicial, section \ref{simplicial-section-standard}.

\medskip\noindent
Supposons maintenant que $\mathcal{A} \to \mathcal{B}$ soit un homomorphisme
de faisceaux d'anneaux. Alors $\mathcal{B}$ est un objet de la catégorie
$\mathcal{A}\textit{-Alg}$. Notons
$\mathcal{P}_\bullet = X_\bullet(\mathcal{B})$ la $\mathcal{A}$-algèbre
simpliciale ainsi obtenue. Rappelons que
$\mathcal{P}_0 = \mathcal{A}[\mathcal{B}]$,
$\mathcal{P}_1 = \mathcal{A}[\mathcal{A}[\mathcal{B}]]$, et ainsi de suite.
Rappelons également qu'il existe une augmentation
$$
\epsilon : \mathcal{P}_\bullet \longrightarrow \mathcal{B}
$$
où $\mathcal{B}$ est regardée comme une $\mathcal{A}$-algèbre simpliciale
constante.

\begin{definition}
\label{definition-standard-resolution-sheaves-rings}
Soit $\mathcal{C}$ un site. Soit $\mathcal{A} \to \mathcal{B}$ un
homomorphisme de faisceaux d'anneaux sur $\mathcal{C}$. La {\it résolution
standard de $\mathcal{B}$ sur $\mathcal{A}$} est l'augmentation
$\epsilon : \mathcal{P}_\bullet \to \mathcal{B}$
dont les termes sont
$$
\mathcal{P}_0 = \mathcal{A}[\mathcal{B}],\quad
\mathcal{P}_1 = \mathcal{A}[\mathcal{A}[\mathcal{B}]],\quad \ldots
$$
et dont les applications sont celles construites ci-dessus.
\end{definition}

\noindent
Cette définition étant posée, le complexe cotangent d'un homomorphisme de
faisceaux d'anneaux se définit comme suit. Nous utiliserons le module des
différentielles défini dans Modules sur les sites, section
\ref{sites-modules-section-differentials}.

\begin{definition}
\label{definition-cotangent-complex-morphism-sheaves-rings}
Soit $\mathcal{C}$ un site. Soit $\mathcal{A} \to \mathcal{B}$ un
homomorphisme de faisceaux d'anneaux sur $\mathcal{C}$. Le {\it complexe
cotangent} $L_{\mathcal{B}/\mathcal{A}}$ est le complexe de
$\mathcal{B}$-modules associé au module simplicial
$$
\Omega_{\mathcal{P}_\bullet/\mathcal{A}}
\otimes_{\mathcal{P}_\bullet, \epsilon} \mathcal{B}
$$
où $\epsilon : \mathcal{P}_\bullet \to \mathcal{B}$ est la résolution standard
de $\mathcal{B}$ sur $\mathcal{A}$. Nous considérons généralement
$L_{\mathcal{B}/\mathcal{A}}$ comme un objet de $D(\mathcal{B})$.
\end{definition}

\noindent
Ces constructions possèdent une fonctorialité analogue à celle étudiée dans la
section \ref{section-functoriality}. Plus précisément, étant donné un diagramme
commutatif
\begin{equation}
\label{equation-commutative-square-sheaves}
\vcenter{
\xymatrix{
\mathcal{B} \ar[r] & \mathcal{B}' \\
\mathcal{A} \ar[u] \ar[r] & \mathcal{A}' \ar[u]
}
}
\end{equation}
de faisceaux d'anneaux sur $\mathcal{C}$, il existe une application canonique
$\mathcal{B}$-linéaire de complexes
$$
L_{\mathcal{B}/\mathcal{A}} \longrightarrow L_{\mathcal{B}'/\mathcal{A}'}
$$
construite de la manière suivante. Si $\mathcal{P}_\bullet \to \mathcal{B}$ est
la résolution standard de $\mathcal{B}$ sur $\mathcal{A}$ et si
$\mathcal{P}'_\bullet \to \mathcal{B}'$ est la résolution standard de
$\mathcal{B}'$ sur $\mathcal{A}'$, il existe une application canonique
$\mathcal{P}_\bullet \to \mathcal{P}'_\bullet$ de $\mathcal{A}$-algèbres
simpliciales, compatible avec les augmentations
$\mathcal{P}_\bullet \to \mathcal{B}$ et
$\mathcal{P}'_\bullet \to \mathcal{B}'$. Les applications
$$
\mathcal{P}_0 = \mathcal{A}[\mathcal{B}]
\longrightarrow
\mathcal{A}'[\mathcal{B}'] = \mathcal{P}'_0,
\quad
\mathcal{P}_1 = \mathcal{A}[\mathcal{A}[\mathcal{B}]]
\longrightarrow
\mathcal{A}'[\mathcal{A}'[\mathcal{B}']] = \mathcal{P}'_1
$$
et ainsi de suite sont données par les homomorphismes
$\mathcal{A} \to \mathcal{A}'$ et $\mathcal{B} \to \mathcal{B}'$. L'application
recherchée $L_{\mathcal{B}/\mathcal{A}} \to L_{\mathcal{B}'/\mathcal{A}'}$
provient alors des applications associées sur les faisceaux de
différentielles.

\begin{lemma}
\label{lemma-pullback-cotangent-morphism-topoi}
Soit $f : \Sh(\mathcal{D}) \to \Sh(\mathcal{C})$ un morphisme de topos. Soit
$\mathcal{A} \to \mathcal{B}$ un homomorphisme de faisceaux d'anneaux sur
$\mathcal{C}$. Alors
$f^{-1}L_{\mathcal{B}/\mathcal{A}} = L_{f^{-1}\mathcal{B}/f^{-1}\mathcal{A}}$.
\end{lemma}

\begin{proof}
Le diagramme
$$
\xymatrix{
\mathcal{A}\textit{-Alg} \ar[d]_{f^{-1}} \ar[r] &
\Sh(\mathcal{C}) \ar@<1ex>[l] \ar[d]^{f^{-1}} \\
f^{-1}\mathcal{A}\textit{-Alg} \ar[r] & \Sh(\mathcal{D}) \ar@<1ex>[l]
}
$$
est commutatif.
\end{proof}

\begin{lemma}
\label{lemma-compute-L-morphism-sheaves-rings}
Soit $\mathcal{C}$ un site. Soit $\mathcal{A} \to \mathcal{B}$ un
homomorphisme de faisceaux d'anneaux sur $\mathcal{C}$. Alors
$H^i(L_{\mathcal{B}/\mathcal{A}})$ est le faisceau associé au préfaisceau
$U \mapsto H^i(L_{\mathcal{B}(U)/\mathcal{A}(U)})$.
\end{lemma}

\begin{proof}
Soit $\mathcal{C}'$ le site obtenu en munissant $\mathcal{C}$ de la topologie
chaotique (les préfaisceaux sont alors des faisceaux). Il existe un morphisme
de topos $f : \Sh(\mathcal{C}) \to \Sh(\mathcal{C}')$, où $f_*$ est
l'inclusion des faisceaux dans les préfaisceaux et $f^{-1}$ la faisceautisation.
D'après le lemme \ref{lemma-pullback-cotangent-morphism-topoi}, il suffit de
démontrer le résultat pour $\mathcal{C}'$, c'est-à-dire lorsque
$\mathcal{C}$ est muni de la topologie chaotique.

\medskip\noindent
Si $\mathcal{C}$ est muni de la topologie chaotique, alors
$L_{\mathcal{B}/\mathcal{A}}(U)$ est égal à
$L_{\mathcal{B}(U)/\mathcal{A}(U)}$, car le diagramme
$$
\xymatrix{
\mathcal{A}\textit{-Alg} \ar[d]_{\text{sections sur }U} \ar[r] &
\Sh(\mathcal{C}) \ar@<1ex>[l] \ar[d]^{\text{sections sur }U} \\
\mathcal{A}(U)\textit{-Alg} \ar[r] & \textit{Ensembles} \ar@<1ex>[l]
}
$$
est commutatif.
\end{proof}

\begin{remark}
\label{remark-map-sections-over-U}
Il ressort de la démonstration du lemme
\ref{lemma-compute-L-morphism-sheaves-rings} que, pour tout
$U \in \Ob(\mathcal{C})$, il existe une application canonique
$L_{\mathcal{B}(U)/\mathcal{A}(U)} \to L_{\mathcal{B}/\mathcal{A}}(U)$
de complexes de $\mathcal{B}(U)$-modules. De plus, ces applications sont
compatibles aux applications de restriction, et le complexe
$L_{\mathcal{B}/\mathcal{A}}$ est la faisceautisation de la règle
$U \mapsto L_{\mathcal{B}(U)/\mathcal{A}(U)}$.
\end{remark}

\begin{lemma}
\label{lemma-H0-L-morphism-sheaves-rings}
Soit $\mathcal{C}$ un site. Soit $\mathcal{A} \to \mathcal{B}$ un
homomorphisme de faisceaux d'anneaux sur $\mathcal{C}$. Alors
$H^0(L_{\mathcal{B}/\mathcal{A}}) = \Omega_{\mathcal{B}/\mathcal{A}}$.
\end{lemma}

\begin{proof}
Cela résulte des lemmes \ref{lemma-compute-L-morphism-sheaves-rings} et
\ref{lemma-identify-H0}, ainsi que de Modules sur les sites, lemme
\ref{sites-modules-lemma-differentials-sheafify}.
\end{proof}

\begin{lemma}
\label{lemma-compute-L-product-sheaves-rings}
Soit $\mathcal{C}$ un site. Soient $\mathcal{A} \to \mathcal{B}$ et
$\mathcal{A} \to \mathcal{B}'$ des homomorphismes de faisceaux d'anneaux sur
$\mathcal{C}$. Alors
$$
L_{\mathcal{B} \times \mathcal{B}'/\mathcal{A}}
\longrightarrow
L_{\mathcal{B}/\mathcal{A}} \oplus L_{\mathcal{B}'/\mathcal{A}}
$$
est un isomorphisme dans $D(\mathcal{B} \times \mathcal{B}')$.
\end{lemma}

\begin{proof}
D'après le lemme \ref{lemma-compute-L-morphism-sheaves-rings}, il suffit de le
démontrer pour des homomorphismes d'anneaux. Dans le cas des anneaux, il s'agit
du lemme \ref{lemma-cotangent-complex-product}.
\end{proof}

\noindent
Le triangle fondamental pour le complexe cotangent de faisceaux d'anneaux est
une conséquence immédiate du résultat pour les homomorphismes d'anneaux.

\begin{lemma}
\label{lemma-triangle-sheaves-rings}
Soit $\mathcal{D}$ un site. Soient
$\mathcal{A} \to \mathcal{B} \to \mathcal{C}$ des homomorphismes de faisceaux
d'anneaux sur $\mathcal{D}$. Il existe un triangle distingué canonique
$$
L_{\mathcal{B}/\mathcal{A}} \otimes_\mathcal{B}^\mathbf{L} \mathcal{C}
\to L_{\mathcal{C}/\mathcal{A}} \to L_{\mathcal{C}/\mathcal{B}} \to
L_{\mathcal{B}/\mathcal{A}} \otimes_\mathcal{B}^\mathbf{L} \mathcal{C}[1]
$$
dans $D(\mathcal{C})$.
\end{lemma}

\begin{proof}
Nous utiliserons la méthode décrite dans les remarques \ref{remark-triangle} et
\ref{remark-explicit-map} pour construire le triangle, et ferons librement
usage des résultats qui y sont mentionnés. Comme dans ces remarques, nous
construisons d'abord le triangle lorsque
$\mathcal{B} \to \mathcal{C}$ est un homomorphisme injectif de faisceaux
d'anneaux. Dans ce cas, posons :
\begin{enumerate}
\item $\mathcal{P}_\bullet$ est la résolution standard de $\mathcal{B}$ sur
$\mathcal{A}$,
\item $\mathcal{Q}_\bullet$ est la résolution standard de $\mathcal{C}$ sur
$\mathcal{A}$,
\item $\mathcal{R}_\bullet$ est la résolution standard de $\mathcal{C}$ sur
$\mathcal{B}$,
\item $\mathcal{S}_\bullet$ est la résolution standard de $\mathcal{B}$ sur
$\mathcal{B}$,
\item $\overline{\mathcal{Q}}_\bullet =
\mathcal{Q}_\bullet \otimes_{\mathcal{P}_\bullet} \mathcal{B}$, et
\item $\overline{\mathcal{R}}_\bullet =
\mathcal{R}_\bullet \otimes_{\mathcal{S}_\bullet} \mathcal{B}$.
\end{enumerate}
Le triangle distingué est celui qui est associé à la suite exacte courte de
$\mathcal{C}$-modules simpliciaux
$$
0 \to
\Omega_{\mathcal{P}_\bullet/\mathcal{A}}
\otimes_{\mathcal{P}_\bullet} \mathcal{C} \to
\Omega_{\mathcal{Q}_\bullet/\mathcal{A}}
\otimes_{\mathcal{Q}_\bullet} \mathcal{C} \to
\Omega_{\overline{\mathcal{Q}}_\bullet/\mathcal{B}}
\otimes_{\overline{\mathcal{Q}}_\bullet} \mathcal{C} \to 0
$$
Les deux premiers termes sont égaux aux deux premiers termes du triangle de
l'énoncé. L'identification du dernier terme avec
$L_{\mathcal{C}/\mathcal{B}}$ utilise les quasi-isomorphismes de complexes
$$
L_{\mathcal{C}/\mathcal{B}} =
\Omega_{\mathcal{R}_\bullet/\mathcal{B}}
\otimes_{\mathcal{R}_\bullet} \mathcal{C}
\longrightarrow
\Omega_{\overline{\mathcal{R}}_\bullet/\mathcal{B}}
\otimes_{\overline{\mathcal{R}}_\bullet} \mathcal{C}
\longleftarrow
\Omega_{\overline{\mathcal{Q}}_\bullet/\mathcal{B}}
\otimes_{\overline{\mathcal{Q}}_\bullet} \mathcal{C}
$$
Toutes les constructions utilisées ci-dessus peuvent d'abord être effectuées
au niveau des préfaisceaux, puis faisceautisées. Ainsi, pour démontrer que des
suites sont exactes ou que des applications sont des quasi-isomorphismes, il
suffit de démontrer l'énoncé correspondant pour les homomorphismes d'anneaux
$\mathcal{A}(U) \to \mathcal{B}(U) \to \mathcal{C}(U)$
pour lesquels le résultat est connu. Ceci achève la démonstration lorsque
$\mathcal{B} \to \mathcal{C}$ est injectif.

\medskip\noindent
Dans le cas général, nous nous ramenons au cas où
$\mathcal{B} \to \mathcal{C}$ est injectif en remplaçant, si nécessaire,
$\mathcal{C}$ par $\mathcal{B} \times \mathcal{C}$. Cela est possible grâce à
l'argument de la remarque \ref{remark-triangle} et au lemme
\ref{lemma-compute-L-product-sheaves-rings}.
\end{proof}

\begin{lemma}
\label{lemma-stalk-cotangent-complex}
Soit $\mathcal{C}$ un site. Soit $\mathcal{A} \to \mathcal{B}$ un
homomorphisme de faisceaux d'anneaux sur $\mathcal{C}$. Si $p$ est un point de
$\mathcal{C}$, alors
$(L_{\mathcal{B}/\mathcal{A}})_p = L_{\mathcal{B}_p/\mathcal{A}_p}$.
\end{lemma}

\begin{proof}
C'est un cas particulier du lemme
\ref{lemma-pullback-cotangent-morphism-topoi}.
\end{proof}

\noindent
Pour la construction du complexe cotangent naïf et ses propriétés, nous
renvoyons à Modules sur les sites, section
\ref{sites-modules-section-netherlander}.

\begin{lemma}
\label{lemma-compare-cotangent-complex-with-naive}
Soit $\mathcal{C}$ un site. Soit $\mathcal{A} \to \mathcal{B}$ un
homomorphisme de faisceaux d'anneaux sur $\mathcal{C}$. Il existe une
application canonique
$L_{\mathcal{B}/\mathcal{A}} \to \NL_{\mathcal{B}/\mathcal{A}}$ qui identifie
le complexe cotangent naïf à la troncation
$\tau_{\geq -1}L_{\mathcal{B}/\mathcal{A}}$.
\end{lemma}

\begin{proof}
Soit $\mathcal{P}_\bullet$ la résolution standard de $\mathcal{B}$ sur
$\mathcal{A}$. Posons
$\mathcal{I} = \Ker(\mathcal{A}[\mathcal{B}] \to \mathcal{B})$. Rappelons que
$\mathcal{P}_0 = \mathcal{A}[\mathcal{B}]$. L'application du lemme est donnée
par le diagramme commutatif
$$
\xymatrix{
L_{\mathcal{B}/\mathcal{A}} \ar[d] & \ldots \ar[r] &
\Omega_{\mathcal{P}_2/\mathcal{A}} \otimes_{\mathcal{P}_2} \mathcal{B}
\ar[r] \ar[d] &
\Omega_{\mathcal{P}_1/\mathcal{A}} \otimes_{\mathcal{P}_1} \mathcal{B}
\ar[r] \ar[d] &
\Omega_{\mathcal{P}_0/\mathcal{A}} \otimes_{\mathcal{P}_0} \mathcal{B}
\ar[d] \\
\NL_{\mathcal{B}/\mathcal{A}} & \ldots \ar[r] &
0 \ar[r] & 
\mathcal{I}/\mathcal{I}^2 \ar[r] &
\Omega_{\mathcal{P}_0/\mathcal{A}} \otimes_{\mathcal{P}_0} \mathcal{B}
}
$$
Nous construisons la flèche descendante de but
$\mathcal{I}/\mathcal{I}^2$ en envoyant une section locale
$\text{d}f \otimes b$ sur la classe de $(d_0(f) - d_1(f))b$ dans
$\mathcal{I}/\mathcal{I}^2$. Ici, $d_i : \mathcal{P}_1 \to \mathcal{P}_0$,
$i = 0, 1$, sont les deux applications de face de la structure simpliciale.
Cela a un sens puisque $d_0 - d_1$ envoie $\mathcal{P}_1$ dans
$\mathcal{I} = \Ker(\mathcal{P}_0 \to \mathcal{B})$. Nous omettons de vérifier
que cette règle est bien définie. Notre application est compatible avec la
différentielle
$\Omega_{\mathcal{P}_1/\mathcal{A}} \otimes_{\mathcal{P}_1} \mathcal{B}
\to \Omega_{\mathcal{P}_0/\mathcal{A}} \otimes_{\mathcal{P}_0} \mathcal{B}$
car celle-ci envoie une section locale $\text{d}f \otimes b$ sur
$\text{d}(d_0(f) - d_1(f)) \otimes b$. En outre, la différentielle
$\Omega_{\mathcal{P}_2/\mathcal{A}} \otimes_{\mathcal{P}_2} \mathcal{B}
\to \Omega_{\mathcal{P}_1/\mathcal{A}} \otimes_{\mathcal{P}_1} \mathcal{B}$
envoie une section locale $\text{d}f \otimes b$ sur
$\text{d}(d_0(f) - d_1(f) + d_2(f)) \otimes b$, qui est annulé par notre
flèche descendante. Nous obtenons ainsi une application de complexes.

\medskip\noindent
Pour voir que notre application induit un isomorphisme sur les faisceaux de
cohomologie $H^0$ et $H^{-1}$, raisonnons comme suit. Soit $\mathcal{C}'$ le
site ayant même catégorie sous-jacente que $\mathcal{C}$, mais muni de la
topologie chaotique. Soit $f : \Sh(\mathcal{C}) \to \Sh(\mathcal{C}')$ le
morphisme de topos dont le foncteur image inverse est la faisceautisation.
Considérons l'homomorphisme donné $\mathcal{A}' \to \mathcal{B}'$ comme un
homomorphisme de faisceaux d'anneaux sur $\mathcal{C}'$. La construction
précédente donne sur $\mathcal{C}'$ une application
$L_{\mathcal{B}'/\mathcal{A}'} \to \NL_{\mathcal{B}'/\mathcal{A}'}$ dont la
valeur sur tout objet $U$ de $\mathcal{C}'$ est simplement l'application
$$
L_{\mathcal{B}(U)/\mathcal{A}(U)} \to \NL_{\mathcal{B}(U)/\mathcal{A}(U)}
$$
de la remarque \ref{remark-explicit-comparison-map}, qui induit un
isomorphisme sur $H^0$ et $H^{-1}$. Puisque
$f^{-1}L_{\mathcal{B}'/\mathcal{A}'} = L_{\mathcal{B}/\mathcal{A}}$
(lemme \ref{lemma-pullback-cotangent-morphism-topoi}) et
$f^{-1}\NL_{\mathcal{B}'/\mathcal{A}'} = \NL_{\mathcal{B}/\mathcal{A}}$
(Modules sur les sites, lemme \ref{sites-modules-lemma-pullback-NL}), le lemme
est démontré.
\end{proof}











\section{La classe d'Atiyah d'un faisceau de modules}
\label{section-atiyah-general}

\noindent
Soit $\mathcal{C}$ un site. Soit $\mathcal{A} \to \mathcal{B}$ un
homomorphisme de faisceaux d'anneaux. Soit $\mathcal{F}$ un faisceau de
$\mathcal{B}$-modules. Soit $\mathcal{P}_\bullet \to \mathcal{B}$ la
résolution standard de $\mathcal{B}$ sur $\mathcal{A}$ (section
\ref{section-cotangent-complex}). Pour tout $n \geq 0$, considérons l'extension
des parties principales
\begin{equation}
\label{equation-atiyah-extension}
0 \to
\Omega_{\mathcal{P}_n/\mathcal{A}} \otimes_{\mathcal{P}_n} \mathcal{F} \to
\mathcal{P}^1_{\mathcal{P}_n/\mathcal{A}}(\mathcal{F}) \to
\mathcal{F} \to 0
\end{equation}
voir Modules sur les sites, lemme
\ref{sites-modules-lemma-sequence-of-principal-parts}. La fonctorialité de
cette construction (Modules sur les sites, remarque
\ref{sites-modules-remark-functoriality-principal-parts}) montre que
(\ref{equation-atiyah-extension}) est la partie de degré $n$ d'une suite
exacte courte de $\mathcal{P}_\bullet$-modules simpliciaux (Cohomologie sur les
sites, section \ref{sites-cohomology-section-simplicial-modules}). En utilisant
le foncteur $L\pi_! : D(\mathcal{P}_\bullet) \to D(\mathcal{B})$ de Cohomologie
sur les sites, remarque \ref{sites-cohomology-remark-homology-augmentation}
(nous utilisons ici le fait que $\mathcal{P}_\bullet \to \mathcal{A}$ est une
résolution), nous obtenons un triangle distingué
\begin{equation}
\label{equation-atiyah-general}
L_{\mathcal{B}/\mathcal{A}} \otimes_\mathcal{B}^\mathbf{L} \mathcal{F} \to
L\pi_!\left(\mathcal{P}^1_{\mathcal{P}_\bullet/\mathcal{A}}(\mathcal{F})\right)
\to \mathcal{F} \to
L_{\mathcal{B}/\mathcal{A}} \otimes_\mathcal{B}^\mathbf{L} \mathcal{F} [1]
\end{equation}
dans $D(\mathcal{B})$.

\begin{definition}
\label{definition-atiyah-class-general}
Soit $\mathcal{C}$ un site. Soit $\mathcal{A} \to \mathcal{B}$ un
homomorphisme de faisceaux d'anneaux. Soit $\mathcal{F}$ un faisceau de
$\mathcal{B}$-modules. L'application $\mathcal{F} \to
L_{\mathcal{B}/\mathcal{A}} \otimes_\mathcal{B}^\mathbf{L} \mathcal{F}[1]$
de (\ref{equation-atiyah-general}) est appelée la {\it classe d'Atiyah} de
$\mathcal{F}$.
\end{definition}









\section{Le complexe cotangent d'un morphisme d'espaces annelés}
\label{section-cotangent-morphism-ringed-spaces}

\noindent
Le complexe cotangent d'un morphisme d'espaces annelés se définit au moyen du
complexe cotangent défini ci-dessus.

\begin{definition}
\label{definition-cotangent-complex-morphism-ringed-spaces}
Soit $f : (X, \mathcal{O}_X) \to (S, \mathcal{O}_S)$ un morphisme d'espaces
annelés. Le {\it complexe cotangent} $L_f$ de $f$ est
$L_f = L_{\mathcal{O}_X/f^{-1}\mathcal{O}_S}$.
Nous utiliserons également la notation
$L_f = L_{X/S} = L_{\mathcal{O}_X/\mathcal{O}_S}$.
\end{definition}

\noindent
Plus précisément, nous considérons le complexe cotangent (définition
\ref{definition-cotangent-complex-morphism-sheaves-rings}) de l'homomorphisme
$f^\sharp : f^{-1}\mathcal{O}_S \to \mathcal{O}_X$ de faisceaux d'anneaux sur
le site associé à l'espace topologique $X$ (Sites, exemple
\ref{sites-example-site-topological}).

\begin{lemma}
\label{lemma-H0-L-morphism-ringed-spaces}
Soit $f : (X, \mathcal{O}_X) \to (S, \mathcal{O}_S)$ un morphisme d'espaces
annelés. Alors $H^0(L_{X/S}) = \Omega_{X/S}$.
\end{lemma}

\begin{proof}
C'est un cas particulier du lemme \ref{lemma-H0-L-morphism-sheaves-rings}.
\end{proof}

\begin{lemma}
\label{lemma-triangle-ringed-spaces}
Soient $f : X \to Y$ et $g : Y \to Z$ des morphismes d'espaces annelés. Il
existe alors un triangle distingué canonique
$$
Lf^* L_{Y/Z} \to L_{X/Z} \to L_{X/Y} \to Lf^*L_{Y/Z}[1]
$$
dans $D(\mathcal{O}_X)$.
\end{lemma}

\begin{proof}
Posons $h = g \circ f$, de sorte que
$h^{-1}\mathcal{O}_Z = f^{-1}g^{-1}\mathcal{O}_Z$. D'après le lemme
\ref{lemma-pullback-cotangent-morphism-topoi}, nous avons
$f^{-1}L_{Y/Z} = L_{f^{-1}\mathcal{O}_Y/h^{-1}\mathcal{O}_Z}$
et celui-ci est un complexe de $f^{-1}\mathcal{O}_Y$-modules plats. Le triangle
distingué précédent est donc un cas du triangle distingué du lemme
\ref{lemma-triangle-sheaves-rings}, avec
$\mathcal{A} = h^{-1}\mathcal{O}_Z$, $\mathcal{B} = f^{-1}\mathcal{O}_Y$ et
$\mathcal{C} = \mathcal{O}_X$.
\end{proof}

\begin{lemma}
\label{lemma-compare-cotangent-complex-with-naive-ringed-spaces}
Soit $f : (X, \mathcal{O}_X) \to (Y, \mathcal{O}_Y)$ un morphisme d'espaces
annelés. Il existe une application canonique $L_{X/Y} \to \NL_{X/Y}$ qui
identifie le complexe cotangent naïf à la troncation
$\tau_{\geq -1}L_{X/Y}$.
\end{lemma}

\begin{proof}
C'est un cas particulier du lemme
\ref{lemma-compare-cotangent-complex-with-naive}.
\end{proof}






\section{Déformations d'espaces annelés et complexe cotangent}
\label{section-deformations-ringed-spaces}

\noindent
Cette section prolonge Théorie des déformations, section
\ref{defos-section-deformations-ringed-spaces}, que nous invitons vivement le
lecteur à lire d'abord. Rappelons brièvement la situation. Nous avons un
épaississement du premier ordre
$t : (S, \mathcal{O}_S) \to (S', \mathcal{O}_{S'})$ d'espaces annelés, avec
$\mathcal{J} = \Ker(t^\sharp)$, un morphisme d'espaces annelés
$f : (X, \mathcal{O}_X) \to (S, \mathcal{O}_S)$, un $\mathcal{O}_X$-module
$\mathcal{G}$ et un $f$-morphisme $c : \mathcal{J} \to \mathcal{G}$ de
faisceaux de modules. Nous cherchons si l'on peut remplacer le point
d'interrogation dans le diagramme suivant
\begin{equation}
\label{equation-to-solve-ringed-spaces}
\vcenter{
\xymatrix{
0 \ar[r] & \mathcal{G} \ar[r] & {?} \ar[r] & \mathcal{O}_X \ar[r] & 0 \\
0 \ar[r] & \mathcal{J} \ar[u]^c \ar[r] & \mathcal{O}_{S'} \ar[u] \ar[r] &
\mathcal{O}_S \ar[u] \ar[r] & 0
}
}
\end{equation}
et, si une solution existe, à quel point elle est unique. Plus précisément,
nous cherchons un épaississement du premier ordre
$i : (X, \mathcal{O}_X) \to (X', \mathcal{O}_{X'})$
et un morphisme d'épaississements $(f, f')$ comme dans Théorie des
déformations, équation (\ref{defos-equation-morphism-thickenings}), où
$\Ker(i^\sharp)$ est identifié à $\mathcal{G}$, tels que $(f')^\sharp$ induise
l'application donnée $c$. Nous dirons que $X'$ est une {\it solution} de
(\ref{equation-to-solve-ringed-spaces}).

\begin{lemma}
\label{lemma-find-obstruction-ringed-spaces}
Dans la situation précédente :
\begin{enumerate}
\item il existe un élément canonique
$\xi \in \Ext^2_{\mathcal{O}_X}(L_{X/S}, \mathcal{G})$
dont l'annulation est une condition nécessaire et suffisante pour l'existence
d'une solution de (\ref{equation-to-solve-ringed-spaces}) ;
\item s'il existe une solution, l'ensemble des classes d'isomorphisme de
solutions est un espace principal homogène sous
$\Ext^1_{\mathcal{O}_X}(L_{X/S}, \mathcal{G})$ ;
\item étant donnée une solution $X'$, l'ensemble des automorphismes de $X'$
qui s'insèrent dans (\ref{equation-to-solve-ringed-spaces}) est canoniquement
isomorphe à $\Ext^0_{\mathcal{O}_X}(L_{X/S}, \mathcal{G})$.
\end{enumerate}
\end{lemma}

\begin{proof}
Au moyen des identifications $\NL_{X/S} = \tau_{\geq -1}L_{X/S}$ (lemme
\ref{lemma-compare-cotangent-complex-with-naive-ringed-spaces}) et
$H^0(L_{X/S}) = \Omega_{X/S}$
(lemme \ref{lemma-H0-L-morphism-ringed-spaces}), nous avons déjà établi (2) et
(3) dans Théorie des déformations, lemmes
\ref{defos-lemma-huge-diagram-ringed-spaces} et
\ref{defos-lemma-choices-ringed-spaces}.

\medskip\noindent
Démonstration de (1). En gros, l'assertion découle de la discussion de Théorie
des déformations, remarque
\ref{defos-remark-parametrize-solutions-ringed-spaces}, en remplaçant le
complexe cotangent naïf par le complexe cotangent entier. Voici une explication
plus détaillée. D'après Théorie des déformations, lemme
\ref{defos-lemma-parametrize-solutions-ringed-spaces}, il existe un élément
$$
\xi' \in
\Ext^1_{\mathcal{O}_X}(Lf^*\NL_{S/S'}, \mathcal{G}) =
\Ext^1_{\mathcal{O}_X}(Lf^*L_{S/S'}, \mathcal{G})
$$
tel qu'une solution existe si et seulement si cet élément appartient à l'image
de l'application
$$
\Ext^1_{\mathcal{O}_X}(NL_{X/S'}, \mathcal{G}) =
\Ext^1_{\mathcal{O}_X}(L_{X/S'}, \mathcal{G})
\longrightarrow
\Ext^1_{\mathcal{O}_X}(Lf^*L_{S/S'}, \mathcal{G})
$$
Le triangle distingué du lemme \ref{lemma-triangle-ringed-spaces}, appliqué à
$X \to S \to S'$, donne naissance à une suite exacte longue
$$
\ldots \to
\Ext^1_{\mathcal{O}_X}(L_{X/S'}, \mathcal{G}) \to
\Ext^1_{\mathcal{O}_X}(Lf^*L_{S/S'}, \mathcal{G}) \to
\Ext^2_{\mathcal{O}_X}(L_{X/S}, \mathcal{G}) \to \ldots
$$
Il suffit donc de prendre pour $\xi$ l'image de $\xi'$.
\end{proof}










\section{Le complexe cotangent d'un morphisme de topos annelés}
\label{section-cotangent-morphism-ringed-topoi}

\noindent
Le complexe cotangent d'un morphisme de topos annelés se définit au moyen du
complexe cotangent défini ci-dessus.

\begin{definition}
\label{definition-cotangent-complex-morphism-ringed-topoi}
Soit $(f, f^\sharp) : (\Sh(\mathcal{C}), \mathcal{O}_\mathcal{C}) \to
(\Sh(\mathcal{D}), \mathcal{O}_\mathcal{D})$ un morphisme de topos annelés.
Le {\it complexe cotangent} $L_f$ de $f$ est
$L_f = L_{\mathcal{O}_\mathcal{C}/f^{-1}\mathcal{O}_\mathcal{D}}$.
Nous écrivons parfois
$L_f = L_{\mathcal{O}_\mathcal{C}/\mathcal{O}_\mathcal{D}}$.
\end{definition}

\noindent
Cette définition s'applique dans de nombreuses situations, mais elle ne donne
pas toujours l'objet attendu. Par exemple, si $f : X \to Y$ est un morphisme
de schémas, alors $f$ induit un morphisme de grands sites étales
$f_{big} : (\Sch/X)_\etale \to (\Sch/Y)_\etale$
qui est un morphisme de topos annelés (Descente, remarque
\ref{descent-remark-change-topologies-ringed}). Toutefois,
$L_{f_{big}} = 0$ puisque $(f_{big})^\sharp$ est un isomorphisme. En revanche,
si nous prenons $L_f$ en regardant $f$ comme un morphisme entre les topos
annelés de Zariski sous-jacents, alors $L_f$ coïncide bien avec le complexe
cotangent $L_{X/Y}$ (défini ci-dessous), dont le faisceau de cohomologie de
degré zéro est $\Omega_{X/Y}$.

\begin{lemma}
\label{lemma-H0-L-morphism-ringed-topoi}
Soit $f : (\Sh(\mathcal{C}), \mathcal{O}) \to
(\Sh(\mathcal{B}), \mathcal{O}_\mathcal{B})$ un morphisme de topos annelés.
Alors $H^0(L_f) = \Omega_f$.
\end{lemma}

\begin{proof}
C'est un cas particulier du lemme \ref{lemma-H0-L-morphism-sheaves-rings}.
\end{proof}

\begin{lemma}
\label{lemma-triangle-ringed-topoi}
Soient $f : (\Sh(\mathcal{C}_1), \mathcal{O}_1) \to
(\Sh(\mathcal{C}_2), \mathcal{O}_2)$ et
$g : (\Sh(\mathcal{C}_2), \mathcal{O}_2) \to
(\Sh(\mathcal{C}_3), \mathcal{O}_3)$ des morphismes de topos annelés. Il
existe alors un triangle distingué canonique
$$
Lf^* L_g \to L_{g \circ f} \to L_f \to Lf^*L_g[1]
$$
dans $D(\mathcal{O}_1)$.
\end{lemma}

\begin{proof}
Posons $h = g \circ f$, de sorte que
$h^{-1}\mathcal{O}_3 = f^{-1}g^{-1}\mathcal{O}_3$. D'après le lemme
\ref{lemma-pullback-cotangent-morphism-topoi}, nous avons
$f^{-1}L_g = L_{f^{-1}\mathcal{O}_2/h^{-1}\mathcal{O}_3}$
et celui-ci est un complexe de $f^{-1}\mathcal{O}_2$-modules plats. Le triangle
distingué précédent est donc un cas du triangle distingué du lemme
\ref{lemma-triangle-sheaves-rings}, avec
$\mathcal{A} = h^{-1}\mathcal{O}_3$, $\mathcal{B} = f^{-1}\mathcal{O}_2$ et
$\mathcal{C} = \mathcal{O}_1$.
\end{proof}

\begin{lemma}
\label{lemma-compare-cotangent-complex-with-naive-ringed-topoi}
Soit $f : (\Sh(\mathcal{C}), \mathcal{O}) \to
(\Sh(\mathcal{B}), \mathcal{O}_\mathcal{B})$ un morphisme de topos annelés. Il
existe une application canonique $L_f \to \NL_f$ qui identifie le complexe
cotangent naïf à la troncation $\tau_{\geq -1}L_f$.
\end{lemma}

\begin{proof}
C'est un cas particulier du lemme
\ref{lemma-compare-cotangent-complex-with-naive}.
\end{proof}








\section{Déformations de topos annelés et complexe cotangent}
\label{section-deformations-ringed-topoi}

\noindent
Cette section prolonge Théorie des déformations, section
\ref{defos-section-deformations-ringed-topoi}, que nous invitons vivement le
lecteur à lire d'abord. Rappelons brièvement la situation. Nous avons un
épaississement du premier ordre
$t : (\Sh(\mathcal{B}), \mathcal{O}_\mathcal{B}) \to
(\Sh(\mathcal{B}'), \mathcal{O}_{\mathcal{B}'})$ de topos annelés, avec
$\mathcal{J} = \Ker(t^\sharp)$, un morphisme de topos annelés
$f : (\Sh(\mathcal{C}), \mathcal{O}) \to
(\Sh(\mathcal{B}), \mathcal{O}_\mathcal{B})$, an $\mathcal{O}$-module
$\mathcal{G}$ et une application $f^{-1}\mathcal{J} \to \mathcal{G}$ de
faisceaux de $f^{-1}\mathcal{O}_\mathcal{B}$-modules. Nous cherchons si l'on
peut remplacer le point d'interrogation dans le diagramme suivant
\begin{equation}
\label{equation-to-solve-ringed-topoi}
\vcenter{
\xymatrix{
0 \ar[r] & \mathcal{G} \ar[r] & {?} \ar[r] & \mathcal{O} \ar[r] & 0 \\
0 \ar[r] & f^{-1}\mathcal{J} \ar[u]^c \ar[r] &
f^{-1}\mathcal{O}_{\mathcal{B}'} \ar[u] \ar[r] &
f^{-1}\mathcal{O}_\mathcal{B} \ar[u] \ar[r] & 0
}
}
\end{equation}
et, si une solution existe, à quel point elle est unique. Plus précisément,
nous cherchons un épaississement du premier ordre
$i : (\Sh(\mathcal{C}), \mathcal{O}) \to (\Sh(\mathcal{C}'), \mathcal{O}')$
et un morphisme d'épaississements $(f, f')$ comme dans Théorie des
déformations, équation
(\ref{defos-equation-morphism-thickenings-ringed-topoi}), où
$\Ker(i^\sharp)$ est identifié à $\mathcal{G}$, tels que $(f')^\sharp$ induise
l'application donnée $c$. Nous dirons que
$(\Sh(\mathcal{C}'), \mathcal{O}')$ est une {\it solution} de
(\ref{equation-to-solve-ringed-topoi}).

\begin{lemma}
\label{lemma-find-obstruction-ringed-topoi}
Dans la situation précédente :
\begin{enumerate}
\item il existe un élément canonique
$\xi \in \Ext^2_\mathcal{O}(L_f, \mathcal{G})$
dont l'annulation est une condition nécessaire et suffisante pour l'existence
d'une solution de (\ref{equation-to-solve-ringed-topoi}) ;
\item s'il existe une solution, l'ensemble des classes d'isomorphisme de
solutions est un espace principal homogène sous
$\Ext^1_\mathcal{O}(L_f, \mathcal{G})$ ;
\item étant donnée une solution $X'$, l'ensemble des automorphismes de $X'$
qui s'insèrent dans (\ref{equation-to-solve-ringed-topoi}) est canoniquement
isomorphe à $\Ext^0_\mathcal{O}(L_f, \mathcal{G})$.
\end{enumerate}
\end{lemma}

\begin{proof}
Au moyen des identifications $\NL_f = \tau_{\geq -1}L_f$ (lemme
\ref{lemma-compare-cotangent-complex-with-naive-ringed-topoi}) et
$H^0(L_f) = \Omega_f$
(lemme \ref{lemma-H0-L-morphism-ringed-topoi}), nous avons déjà établi (2) et
(3) dans Théorie des déformations, lemmes
\ref{defos-lemma-huge-diagram-ringed-topoi} et
\ref{defos-lemma-choices-ringed-topoi}.

\medskip\noindent
Démonstration de (1). Pour accorder les notations à celles de Théorie des
déformations, section \ref{defos-section-deformations-ringed-topoi}, nous
écrirons
$\NL_f = \NL_{\mathcal{O}/\mathcal{O}_\mathcal{B}}$ et
$L_f = L_{\mathcal{O}/\mathcal{O}_\mathcal{B}}$, et de même pour les morphismes
$t$ et $t \circ f$. D'après Théorie des déformations, lemme
\ref{defos-lemma-parametrize-solutions-ringed-topoi}, il existe un élément
$$
\xi' \in
\Ext^1_\mathcal{O}(
Lf^*\NL_{\mathcal{O}_\mathcal{B}/\mathcal{O}_{\mathcal{B}'}}, \mathcal{G}) =
\Ext^1_\mathcal{O}(
Lf^*L_{\mathcal{O}_\mathcal{B}/\mathcal{O}_{\mathcal{B}'}}, \mathcal{G})
$$
tel qu'une solution existe si et seulement si cet élément appartient à l'image
de l'application
$$
\Ext^1_\mathcal{O}(
\NL_{\mathcal{O}/\mathcal{O}_{\mathcal{B}'}}, \mathcal{G}) =
\Ext^1_\mathcal{O}(
L_{\mathcal{O}/\mathcal{O}_{\mathcal{B}'}}, \mathcal{G})
\longrightarrow
\Ext^1_\mathcal{O}(
Lf^*L_{\mathcal{O}_\mathcal{B}/\mathcal{O}_{\mathcal{B}'}}, \mathcal{G})
$$
Le triangle distingué du lemme \ref{lemma-triangle-ringed-topoi}, appliqué à
$f$ et $t$, donne naissance à une suite exacte longue
$$
\ldots \to
\Ext^1_\mathcal{O}(
L_{\mathcal{O}/\mathcal{O}_{\mathcal{B}'}}, \mathcal{G}) \to
\Ext^1_\mathcal{O}(
Lf^*L_{\mathcal{O}_\mathcal{B}/\mathcal{O}_{\mathcal{B}'}}, \mathcal{G})
\to
\Ext^1_\mathcal{O}(
L_{\mathcal{O}/\mathcal{O}_\mathcal{B}}, \mathcal{G})
$$
Il suffit donc de prendre pour $\xi$ l'image de $\xi'$.
\end{proof}












\section{Le complexe cotangent d'un morphisme de schémas}
\label{section-cotangent-morphism-schemes}

\noindent
Comme annoncé plus haut, nous définissons comme suit le complexe cotangent d'un
morphisme de schémas.

\begin{definition}
\label{definition-cotangent-morphism-schemes}
Soit $f : X \to Y$ un morphisme de schémas. Le {\it complexe cotangent
$L_{X/Y}$ de $X$ sur $Y$} est le complexe cotangent de $f$ considéré comme un
morphisme d'espaces annelés (définition
\ref{definition-cotangent-complex-morphism-ringed-spaces}).
\end{definition}

\noindent
En particulier, les résultats de la section
\ref{section-cotangent-morphism-ringed-spaces} s'appliquent aux complexes
cotangents des morphismes de schémas. Le lemme suivant montre que cette
définition est compatible avec celle donnée pour les homomorphismes d'anneaux ;
il implique aussi que $L_{X/Y}$ est un objet de $D_\QCoh(\mathcal{O}_X)$.

\begin{lemma}
\label{lemma-morphism-affine-schemes}
Soit $f : X \to Y$ un morphisme de schémas. Soient
$U = \Spec(B) \subset X$ et $V = \Spec(A) \subset Y$ des ouverts affines tels
que $f(U) \subset V$. Il existe une application canonique
$$
\widetilde{L_{B/A}} \longrightarrow L_{X/Y}|_U
$$
de complexes qui est un isomorphisme dans $D(\mathcal{O}_U)$. Cette application
est compatible à la restriction à des ouverts affines plus petits de $X$ et
de $Y$.
\end{lemma}

\begin{proof}
D'après la remarque \ref{remark-map-sections-over-U}, il existe une
application canonique de complexes
$L_{\mathcal{O}_X(U)/f^{-1}\mathcal{O}_Y(U)} \to L_{X/Y}(U)$
de $B = \mathcal{O}_X(U)$-modules, compatible aux restrictions ultérieures. En
utilisant l'application canonique $A \to f^{-1}\mathcal{O}_Y(U)$, nous
obtenons une application canonique
$L_{B/A} \to L_{\mathcal{O}_X(U)/f^{-1}\mathcal{O}_Y(U)}$
de complexes de $B$-modules. La propriété universelle du foncteur
$\widetilde{\ }$ (voir Schémas, lemme
\ref{schemes-lemma-compare-constructions}) fournit l'application de l'énoncé.
Pour vérifier qu'elle est un isomorphisme sur les faisceaux de cohomologie, il
suffit de vérifier qu'elle induit des isomorphismes sur les fibres. Cela
résulte immédiatement des lemmes \ref{lemma-stalk-cotangent-complex} et
\ref{lemma-localize} (et de la description des fibres de
$\mathcal{O}_X$ et $f^{-1}\mathcal{O}_Y$
en un point $\mathfrak p \in \Spec(B)$ comme $B_\mathfrak p$ et
$A_\mathfrak q$, où $\mathfrak q = A \cap \mathfrak p$ ; les références
utilisées sont Schémas, lemme \ref{schemes-lemma-spec-sheaves}, et Faisceaux,
lemme \ref{sheaves-lemma-stalk-pullback}).
\end{proof}

\begin{lemma}
\label{lemma-scheme-over-ring}
Soit $\Lambda$ un anneau. Soit $X$ un schéma sur $\Lambda$. Alors
$$
L_{X/\Spec(\Lambda)} = L_{\mathcal{O}_X/\underline{\Lambda}}
$$
où $\underline{\Lambda}$ est le faisceau constant de valeur $\Lambda$ sur $X$.
\end{lemma}

\begin{proof}
Soit $p : X \to \Spec(\Lambda)$ le morphisme structural. Soit
$q : \Spec(\Lambda) \to (*, \Lambda)$ le morphisme évident. D'après le triangle
distingué du lemme \ref{lemma-triangle-ringed-spaces}, il suffit de montrer que
$L_q = 0$. Pour cela, il suffit de montrer, pour
$\mathfrak p \in \Spec(\Lambda)$, que
$$
(L_q)_\mathfrak p =
L_{\mathcal{O}_{\Spec(\Lambda), \mathfrak p}/\Lambda} =
L_{\Lambda_\mathfrak p/\Lambda}
$$
(lemme \ref{lemma-stalk-cotangent-complex}) est nul, ce qui résulte du lemme
\ref{lemma-when-zero}.
\end{proof}






\section{Le complexe cotangent d'un schéma sur un anneau}
\label{section-cotangent-schemes-variant}

\noindent
Soit $\Lambda$ un anneau et soit $X$ un schéma sur $\Lambda$. Écrivons
$L_{X/\Spec(\Lambda)} = L_{X/\Lambda}$, ce qui est justifié par le lemme
\ref{lemma-scheme-over-ring}. Dans cette section, nous donnons de
$L_{X/\Lambda}$ une description analogue à celle du lemme
\ref{lemma-compute-cotangent-complex}. Plus précisément, nous construisons une
catégorie $\mathcal{C}_{X/\Lambda}$ fibrée au-dessus de $X_{Zar}$ et la
munissons d'un faisceau $\mathcal{O}$ de $\Lambda$-algèbres (polynomiales) tel
que
$$
L_{X/\Lambda} =
L\pi_!(\Omega_{\mathcal{O}/\underline{\Lambda}} \otimes_\mathcal{O}
\underline{\mathcal{O}}_X).
$$
Nous utiliserons plus loin la catégorie $\mathcal{C}_{X/\Lambda}$ pour
construire une théorie naïve des obstructions pour le champ des faisceaux
cohérents.

\medskip\noindent
Soit $\Lambda$ un anneau. Soit $X$ un schéma sur $\Lambda$. Soit
$\mathcal{C}_{X/\Lambda}$ la catégorie dont les objets sont les diagrammes
commutatifs
\begin{equation}
\label{equation-object}
\vcenter{
\xymatrix{
X \ar[d] & U \ar[l] \ar[d] \\
\Spec(\Lambda) & \mathbf{A} \ar[l]
}
}
\end{equation}
de schémas tels que
\begin{enumerate}
\item $U$ soit un sous-schéma ouvert de $X$,
\item il existe un isomorphisme $\mathbf{A} = \Spec(P)$, où $P$ est une
algèbre polynomiale sur $\Lambda$ (en un certain ensemble de variables).
\end{enumerate}
Autrement dit, $\mathbf{A}$ est un espace affine (de dimension infinie) sur
$\Spec(\Lambda)$. Les morphismes sont donnés par des diagrammes commutatifs.
Rappelons que $X_{Zar}$ désigne le petit site de Zariski de $X$. Il existe un
foncteur d'oubli
$$
u : \mathcal{C}_{X/\Lambda} \to X_{Zar},\ (U \to \mathbf{A}) \mapsto U
$$
Remarquons que la catégorie fibre au-dessus de $U$ est canoniquement
équivalente à la catégorie $\mathcal{C}_{\mathcal{O}_X(U)/\Lambda}$ introduite
dans la section \ref{section-compute-L-pi-shriek}.

\begin{lemma}
\label{lemma-category-fibred}
Dans la situation précédente, la catégorie $\mathcal{C}_{X/\Lambda}$ est
fibrée au-dessus de $X_{Zar}$.
\end{lemma}

\begin{proof}
Étant donnés un objet $U \to \mathbf{A}$ de $\mathcal{C}_{X/\Lambda}$ et un
morphisme $U' \to U$ de $X_{Zar}$, considérons l'objet
$U' \to \mathbf{A}$ de $\mathcal{C}_{X/\Lambda}$, où
$U' \to \mathbf{A}$ est la composée de $U' \to U$ et de
$U \to \mathbf{A}$. Le morphisme
$(U' \to \mathbf{A}) \to (U \to \mathbf{A})$ of
de $\mathcal{C}_{X/\Lambda}$ est fortement cartésien au-dessus de $X_{Zar}$.
\end{proof}

\noindent
Munissons $\mathcal{C}_{X/\Lambda}$ de la topologie héritée de $X_{Zar}$ (voir
Champs, section \ref{stacks-section-topology}). Le foncteur $u$ définit un
morphisme de topos
$\pi : \Sh(\mathcal{C}_{X/\Lambda}) \to \Sh(X_{Zar})$. Le site
$\mathcal{C}_{X/\Lambda}$ est muni de plusieurs faisceaux d'anneaux :
\begin{enumerate}
\item le faisceau $\mathcal{O}$ donné par la règle
$(U \to \mathbf{A}) \mapsto \Gamma(\mathbf{A}, \mathcal{O}_\mathbf{A})$.
\item le faisceau $\underline{\mathcal{O}}_X = \pi^{-1}\mathcal{O}_X$ donné
par la règle $(U \to \mathbf{A}) \mapsto \mathcal{O}_X(U)$ ;
\item le faisceau constant $\underline{\Lambda}$.
\end{enumerate}
Nous obtenons des morphismes de topos annelés
\begin{equation}
\label{equation-pi-schemes}
\vcenter{
\xymatrix{
(\Sh(\mathcal{C}_{X/\Lambda}), \underline{\mathcal{O}}_X) \ar[r]_i \ar[d]_\pi &
(\Sh(\mathcal{C}_{X/\Lambda}), \mathcal{O}) \\
(\Sh(X_{Zar}), \mathcal{O}_X)
}
}
\end{equation}
Le morphisme $i$ est l'identité sur les topos sous-jacents, et
$i^\sharp : \mathcal{O} \to \underline{\mathcal{O}}_X$ est l'homomorphisme
évident. Le morphisme $\pi$ est un cas particulier de Cohomologie sur les
sites, situation \ref{sites-cohomology-situation-fibred-category}. Les
foncteurs dérivés suivants joueront un rôle important :
$
Li^* : D(\mathcal{O}) \longrightarrow D(\underline{\mathcal{O}}_X)
$
adjoint à gauche de
$Ri_* = i_* : D(\underline{\mathcal{O}}_X) \to D(\mathcal{O})$, et
$
L\pi_! : D(\underline{\mathcal{O}}_X) \longrightarrow D(\mathcal{O}_X)
$
adjoint à gauche de
$\pi^* = \pi^{-1} : D(\mathcal{O}_X) \to D(\underline{\mathcal{O}}_X)$. Nos
travaux précédents permettent de calculer $L\pi_!$.

\begin{remark}
\label{remark-compute-L-pi-shriek}
Dans la situation précédente, pour tout ouvert $U \subset X$, soit
$P_{\bullet, U}$ la résolution standard de $\mathcal{O}_X(U)$ sur $\Lambda$.
Posons $\mathbf{A}_{n, U} = \Spec(P_{n, U})$. Alors
$\mathbf{A}_{\bullet, U}$
est un objet cosimplicial de la catégorie fibre
$\mathcal{C}_{\mathcal{O}_X(U)/\Lambda}$ de $\mathcal{C}_{X/\Lambda}$ au-dessus
de $U$. De plus, comme expliqué dans la remarque \ref{remark-resolution},
$\mathbf{A}_{\bullet, U}$ est un objet cosimplicial de
$\mathcal{C}_{\mathcal{O}_X(U)/\Lambda}$ comme dans Cohomologie sur les sites,
lemme \ref{sites-cohomology-lemma-compute-by-cosimplicial-resolution}. La
construction $U \mapsto \mathbf{A}_{\bullet, U}$ étant fonctorielle en $U$,
tout faisceau (abélien) $\mathcal{F}$ sur $\mathcal{C}_{X/\Lambda}$ donne un
complexe de préfaisceaux
$$
U \longmapsto \mathcal{F}(\mathbf{A}_{\bullet, U})
$$
dont les groupes de cohomologie calculent l'homologie de $\mathcal{F}$ sur la
catégorie fibre. D'après Cohomologie sur les sites, lemme
\ref{sites-cohomology-lemma-compute-left-derived-pi-shriek}, nous en concluons
que la faisceautisation calcule $L_n\pi_!(\mathcal{F})$. Autrement dit, le
complexe de faisceaux dont le terme de degré $-n$ est la faisceautisation de
$U \mapsto \mathcal{F}(\mathbf{A}_{n, U})$ calcule $L\pi_!(\mathcal{F})$.
\end{remark}

\noindent
Nous pouvons maintenant énoncer le résultat principal de cette section.

\begin{lemma}
\label{lemma-cotangent-morphism-schemes}
Dans la situation précédente, il existe un isomorphisme canonique
$$
L_{X/\Lambda} = 
L\pi_!(Li^*\Omega_{\mathcal{O}/\underline{\Lambda}}) =
L\pi_!(i^*\Omega_{\mathcal{O}/\underline{\Lambda}}) =
L\pi_!(\Omega_{\mathcal{O}/\underline{\Lambda}}
\otimes_\mathcal{O} \underline{\mathcal{O}}_X)
$$
dans $D(\mathcal{O}_X)$.
\end{lemma}

\begin{proof}
Observons d'abord que, pour tout objet $(U \to \mathbf{A})$ de
$\mathcal{C}_{X/\Lambda}$, la valeur du faisceau $\mathcal{O}$ est une algèbre
polynomiale sur $\Lambda$. Ainsi,
$\Omega_{\mathcal{O}/\underline{\Lambda}}$ est un $\mathcal{O}$-module plat,
et les deuxième et troisième égalités de l'énoncé sont vérifiées.

\medskip\noindent
D'après la remarque \ref{remark-compute-L-pi-shriek}, l'objet
$L\pi_!(\Omega_{\mathcal{O}/\underline{\Lambda}}
\otimes_\mathcal{O} \underline{\mathcal{O}}_X)$
se calcule comme la faisceautisation du complexe de préfaisceaux
$$
U \mapsto
\left(\Omega_{\mathcal{O}/\underline{\Lambda}}
\otimes_\mathcal{O} \underline{\mathcal{O}}_X\right)(\mathbf{A}_{\bullet, U})
=
\Omega_{P_{\bullet, U}/\Lambda} \otimes_{P_{\bullet, U}} \mathcal{O}_X(U) =
L_{\mathcal{O}_X(U)/\Lambda}
$$
avec les notations de la remarque \ref{remark-compute-L-pi-shriek}. La remarque
\ref{remark-map-sections-over-U} montre alors que
$L\pi_!(\Omega_{\mathcal{O}/\underline{\Lambda}}
\otimes_\mathcal{O} \underline{\mathcal{O}}_X)$
calcule le complexe cotangent de l'homomorphisme d'anneaux
$\underline{\Lambda} \to \mathcal{O}_X$ sur $X$. C'est le résultat voulu
d'après le lemme \ref{lemma-scheme-over-ring}.
\end{proof}








\section{Le complexe cotangent d'un morphisme d'espaces algébriques}
\label{section-cotangent-morphism-spaces}

\noindent
Nous définissons le complexe cotangent d'un morphisme d'espaces algébriques au
moyen du morphisme associé entre les petits sites étales.

\begin{definition}
\label{definition-cotangent-morphism-spaces}
Soit $S$ un schéma. Soit $f : X \to Y$ un morphisme d'espaces algébriques sur
$S$. Le {\it complexe cotangent $L_{X/Y}$ de $X$ sur $Y$} est le complexe
cotangent du morphisme de topos annelés $f_{small}$ entre les petits sites
étales de $X$ et de $Y$ (voir Propriétés des espaces, lemme
\ref{spaces-properties-lemma-morphism-ringed-topoi}, et définition
\ref{definition-cotangent-complex-morphism-ringed-topoi}).
\end{definition}

\noindent
En particulier, les résultats de la section
\ref{section-cotangent-morphism-ringed-topoi} s'appliquent aux complexes
cotangents des morphismes d'espaces algébriques. Les lemmes suivants montrent
que cette définition est compatible avec celles données pour les
homomorphismes d'anneaux et pour les schémas, et que $L_{X/Y}$ est un objet de
$D_\QCoh(\mathcal{O}_X)$.

\begin{lemma}
\label{lemma-etale-localization}
Soit $S$ un schéma. Considérons un diagramme commutatif
$$
\xymatrix{
U \ar[d]_p \ar[r]_g & V \ar[d]^q \\
X \ar[r]^f & Y
}
$$
d'espaces algébriques sur $S$, où $p$ et $q$ sont étales. Il existe alors une
identification canonique $L_{X/Y}|_{U_\etale} = L_{U/V}$ dans
$D(\mathcal{O}_U)$.
\end{lemma}

\begin{proof}
La formation du complexe cotangent commute à l'image inverse (lemme
\ref{lemma-pullback-cotangent-morphism-topoi}), et nous avons
$p_{small}^{-1}\mathcal{O}_X = \mathcal{O}_U$ et
$g_{small}^{-1}\mathcal{O}_{V_\etale} =
p_{small}^{-1}f_{small}^{-1}\mathcal{O}_{Y_\etale}$
parce que $q_{small}^{-1}\mathcal{O}_{Y_\etale} =
\mathcal{O}_{V_\etale}$
(Propriétés des espaces, lemme
\ref{spaces-properties-lemma-etale-exact-pullback}).
En suivant les définitions, nous en concluons que
$L_{X/Y}|_{U_\etale} = L_{U/V}$.
\end{proof}

\begin{lemma}
\label{lemma-compare-spaces-schemes}
Soit $S$ un schéma. Soit $f : X \to Y$ un morphisme d'espaces algébriques sur
$S$. Supposons que $X$ et $Y$ soient représentables par des schémas $X_0$ et
$Y_0$. Il existe alors une identification canonique
$L_{X/Y} = \epsilon^*L_{X_0/Y_0}$ in $D(\mathcal{O}_X)$
où $\epsilon$ est celui de Catégories dérivées des espaces, section
\ref{spaces-perfect-section-derived-quasi-coherent-etale}, et où
$L_{X_0/Y_0}$ est celui de la définition
\ref{definition-cotangent-morphism-schemes}.
\end{lemma}

\begin{proof}
Soit $f_0 : X_0 \to Y_0$ le morphisme de schémas correspondant à $f$. Il
existe une application canonique
$\epsilon^{-1}f_0^{-1}\mathcal{O}_{Y_0} \to f_{small}^{-1}\mathcal{O}_Y$
compatible à
$\epsilon^\sharp : \epsilon^{-1}\mathcal{O}_{X_0} \to \mathcal{O}_X$
parce qu'il existe un diagramme commutatif
$$
\xymatrix{
X_{0, Zar} \ar[d]_{f_0} & X_\etale \ar[l]^\epsilon \ar[d]^f \\
Y_{0, Zar} & Y_\etale \ar[l]_\epsilon
}
$$
voir Catégories dérivées des espaces, remarque
\ref{spaces-perfect-remark-match-total-direct-images}.
Nous obtenons ainsi une application canonique
$$
\epsilon^{-1}L_{X_0/Y_0} =
\epsilon^{-1}L_{\mathcal{O}_{X_0}/f_0^{-1}\mathcal{O}_{Y_0}} =
L_{\epsilon^{-1}\mathcal{O}_{X_0}/\epsilon^{-1}f_0^{-1}\mathcal{O}_{Y_0}}
\longrightarrow
L_{\mathcal{O}_X/f^{-1}_{small}\mathcal{O}_Y} = L_{X/Y}
$$
par la fonctorialité étudiée dans la section
\ref{section-cotangent-complex} et le lemme
\ref{lemma-pullback-cotangent-morphism-topoi}. Pour voir que l'application
induite $\epsilon^*L_{X_0/Y_0} \to L_{X/Y}$ est un isomorphisme, nous pouvons
le vérifier sur les fibres aux points géométriques (Propriétés des espaces,
théorème \ref{spaces-properties-theorem-exactness-stalks}). Nous utiliserons
le lemme \ref{lemma-stalk-cotangent-complex} pour calculer les fibres. Soit
$\overline{x} : \Spec(k) \to X_0$ un point géométrique au-dessus de
$x \in X_0$, et soit $\overline{y} = f \circ \overline{x}$ au-dessus de
$y \in Y_0$. Alors
$$
L_{X/Y, \overline{x}} =
L_{\mathcal{O}_{X, \overline{x}}/\mathcal{O}_{Y, \overline{y}}}
$$
et
$$
(\epsilon^*L_{X_0/Y_0})_{\overline{x}} =
L_{X_0/Y_0, x} \otimes_{\mathcal{O}_{X_0, x}}
\mathcal{O}_{X, \overline{x}} =
L_{\mathcal{O}_{X_0, x}/\mathcal{O}_{Y_0, y}}
\otimes_{\mathcal{O}_{X_0, x}} \mathcal{O}_{X, \overline{x}}
$$
Nous omettons quelques détails (indication : utiliser le fait que la fibre
d'une image inverse est la fibre au point image ; voir Sites, lemme
\ref{sites-lemma-point-morphism-sites}, ainsi que le résultat correspondant
pour les modules, Modules sur les sites, lemme
\ref{sites-modules-lemma-pullback-stalk}). Remarquons que
$\mathcal{O}_{X, \overline{x}}$ est le hensélisé strict de
$\mathcal{O}_{X_0, x}$, et de même pour $\mathcal{O}_{Y, \overline{y}}$
(Propriétés des espaces, lemme
\ref{spaces-properties-lemma-describe-etale-local-ring}). Le résultat découle
donc du lemme \ref{lemma-cotangent-complex-henselization}.
\end{proof}

\begin{lemma}
\label{lemma-space-over-ring}
Soit $\Lambda$ un anneau. Soit $X$ un espace algébrique sur $\Lambda$. Alors
$$
L_{X/\Spec(\Lambda)} = L_{\mathcal{O}_X/\underline{\Lambda}}
$$
où $\underline{\Lambda}$ est le faisceau constant de valeur $\Lambda$ sur
$X_\etale$.
\end{lemma}

\begin{proof}
Soit $p : X \to \Spec(\Lambda)$ le morphisme structural. Soit
$q : \Spec(\Lambda)_\etale \to (*, \Lambda)$ le morphisme évident. D'après le
triangle distingué du lemme \ref{lemma-triangle-ringed-topoi}, il suffit de
montrer que $L_q = 0$. Pour cela, il suffit de montrer (Propriétés des espaces,
théorème \ref{spaces-properties-theorem-exactness-stalks}) que, pour tout point
géométrique $\overline{t} : \Spec(k) \to \Spec(\Lambda)$,
$$
(L_q)_{\overline{t}} =
L_{\mathcal{O}_{\Spec(\Lambda)_\etale, \overline{t}}/\Lambda}
$$
(lemme \ref{lemma-stalk-cotangent-complex}) est nul. Comme
$\mathcal{O}_{\Spec(\Lambda)_\etale, \overline{t}}$ est le hensélisé strict
d'un anneau local de $\Lambda$ (Propriétés des espaces, lemme
\ref{spaces-properties-lemma-describe-etale-local-ring}), cela résulte du lemme
\ref{lemma-when-zero}.
\end{proof}










\section{Le complexe cotangent d'un espace algébrique sur un anneau}
\sectionmark{Complexe cotangent sur un anneau}
\label{section-cotangent-spaces-variant}

\noindent
Soit $\Lambda$ un anneau et soit $X$ un espace algébrique sur $\Lambda$.
Écrivons $L_{X/\Spec(\Lambda)} = L_{X/\Lambda}$, ce qui est justifié par le
lemme \ref{lemma-space-over-ring}. Dans cette section, nous donnons de
$L_{X/\Lambda}$ une description analogue à celle du lemme
\ref{lemma-compute-cotangent-complex}. Plus précisément, nous construisons une
catégorie $\mathcal{C}_{X/\Lambda}$ fibrée au-dessus de $X_\etale$ et la
munissons d'un faisceau $\mathcal{O}$ de $\Lambda$-algèbres (polynomiales) tel
que
$$
L_{X/\Lambda} =
L\pi_!(\Omega_{\mathcal{O}/\underline{\Lambda}} \otimes_\mathcal{O}
\underline{\mathcal{O}}_X).
$$
Nous utiliserons plus loin la catégorie $\mathcal{C}_{X/\Lambda}$ pour
construire une théorie naïve des obstructions pour le champ des faisceaux
cohérents.

\medskip\noindent
Soit $\Lambda$ un anneau. Soit $X$ un espace algébrique sur $\Lambda$. Soit
$\mathcal{C}_{X/\Lambda}$ la catégorie dont les objets sont les diagrammes
commutatifs
\begin{equation}
\label{equation-object-space}
\vcenter{
\xymatrix{
X \ar[d] & U \ar[l] \ar[d] \\
\Spec(\Lambda) & \mathbf{A} \ar[l]
}
}
\end{equation}
de schémas tels que
\begin{enumerate}
\item $U$ soit un schéma,
\item $U \to X$ soit étale,
\item il existe un isomorphisme $\mathbf{A} = \Spec(P)$, où $P$ est une
algèbre polynomiale sur $\Lambda$ (en un certain ensemble de variables).
\end{enumerate}
Autrement dit, $\mathbf{A}$ est un espace affine (de dimension infinie) sur
$\Spec(\Lambda)$. Les morphismes sont donnés par des diagrammes commutatifs.
Rappelons que $X_\etale$ désigne le petit site étale de $X$, dont les objets
sont les schémas étales sur $X$. Il existe un foncteur d'oubli
$$
u : \mathcal{C}_{X/\Lambda} \to X_\etale,
\quad
(U \to \mathbf{A}) \mapsto U
$$
Remarquons que la catégorie fibre au-dessus de $U$ est canoniquement
équivalente à la catégorie $\mathcal{C}_{\mathcal{O}_X(U)/\Lambda}$ introduite
dans la section \ref{section-compute-L-pi-shriek}.

\begin{lemma}
\label{lemma-category-fibred-space}
Dans la situation précédente, la catégorie $\mathcal{C}_{X/\Lambda}$ est
fibrée au-dessus de $X_\etale$.
\end{lemma}

\begin{proof}
Étant donnés un objet $U \to \mathbf{A}$ de $\mathcal{C}_{X/\Lambda}$ et un
morphisme $U' \to U$ de $X_\etale$, considérons l'objet
$U' \to \mathbf{A}$ de $\mathcal{C}_{X/\Lambda}$, où
$U' \to \mathbf{A}$ est la composée de $U' \to U$ et de
$U \to \mathbf{A}$. Le morphisme
$(U' \to \mathbf{A}) \to (U \to \mathbf{A})$ de
$\mathcal{C}_{X/\Lambda}$ est fortement cartésien au-dessus de $X_\etale$.
\end{proof}

\noindent
Munissons $\mathcal{C}_{X/\Lambda}$ de la topologie héritée de $X_\etale$
(voir Champs, section \ref{stacks-section-topology}). Le foncteur $u$ définit
un morphisme de topos
$\pi : \Sh(\mathcal{C}_{X/\Lambda}) \to \Sh(X_\etale)$. Le site
$\mathcal{C}_{X/\Lambda}$ est muni de plusieurs faisceaux d'anneaux :
\begin{enumerate}
\item le faisceau $\mathcal{O}$ donné par la règle
$(U \to \mathbf{A}) \mapsto \Gamma(\mathbf{A}, \mathcal{O}_\mathbf{A})$.
\item le faisceau $\underline{\mathcal{O}}_X = \pi^{-1}\mathcal{O}_X$ donné
par la règle $(U \to \mathbf{A}) \mapsto \mathcal{O}_X(U)$ ;
\item le faisceau constant $\underline{\Lambda}$.
\end{enumerate}
Nous obtenons des morphismes de topos annelés
\begin{equation}
\label{equation-pi-spaces}
\vcenter{
\xymatrix{
(\Sh(\mathcal{C}_{X/\Lambda}), \underline{\mathcal{O}}_X) \ar[r]_i \ar[d]_\pi &
(\Sh(\mathcal{C}_{X/\Lambda}), \mathcal{O}) \\
(\Sh(X_\etale), \mathcal{O}_X)
}
}
\end{equation}
Le morphisme $i$ est l'identité sur les topos sous-jacents, et
$i^\sharp : \mathcal{O} \to \underline{\mathcal{O}}_X$ est l'homomorphisme
évident. Le morphisme $\pi$ est un cas particulier de Cohomologie sur les
sites, situation \ref{sites-cohomology-situation-fibred-category}. Les
foncteurs dérivés suivants joueront un rôle important :
$
Li^* : D(\mathcal{O}) \longrightarrow D(\underline{\mathcal{O}}_X)
$
adjoint à gauche de
$Ri_* = i_* : D(\underline{\mathcal{O}}_X) \to D(\mathcal{O})$, et
$
L\pi_! : D(\underline{\mathcal{O}}_X) \longrightarrow D(\mathcal{O}_X)
$
adjoint à gauche de
$\pi^* = \pi^{-1} : D(\mathcal{O}_X) \to D(\underline{\mathcal{O}}_X)$. Nos
travaux précédents permettent de calculer $L\pi_!$.

\begin{remark}
\label{remark-compute-L-pi-shriek-spaces}
Dans la situation précédente, pour tout objet $U \to X$ de $X_\etale$, soit
$P_{\bullet, U}$ la résolution standard de $\mathcal{O}_X(U)$ sur $\Lambda$.
Posons $\mathbf{A}_{n, U} = \Spec(P_{n, U})$. Alors
$\mathbf{A}_{\bullet, U}$ est un objet cosimplicial de la catégorie fibre
$\mathcal{C}_{\mathcal{O}_X(U)/\Lambda}$ de $\mathcal{C}_{X/\Lambda}$ au-dessus
de $U$. De plus, comme expliqué dans la remarque \ref{remark-resolution},
$\mathbf{A}_{\bullet, U}$ est un objet cosimplicial de
$\mathcal{C}_{\mathcal{O}_X(U)/\Lambda}$ comme dans Cohomologie sur les sites,
lemme \ref{sites-cohomology-lemma-compute-by-cosimplicial-resolution}. La
construction $U \mapsto \mathbf{A}_{\bullet, U}$ étant fonctorielle en $U$,
tout faisceau (abélien) $\mathcal{F}$ sur $\mathcal{C}_{X/\Lambda}$ donne un
complexe de préfaisceaux
$$
U \longmapsto \mathcal{F}(\mathbf{A}_{\bullet, U})
$$
dont les groupes de cohomologie calculent l'homologie de $\mathcal{F}$ sur la
catégorie fibre. D'après Cohomologie sur les sites, lemme
\ref{sites-cohomology-lemma-compute-left-derived-pi-shriek}, nous en concluons
que la faisceautisation calcule $L_n\pi_!(\mathcal{F})$. Autrement dit, le
complexe de faisceaux dont le terme de degré $-n$ est la faisceautisation de
$U \mapsto \mathcal{F}(\mathbf{A}_{n, U})$ calcule $L\pi_!(\mathcal{F})$.
\end{remark}

\noindent
Nous pouvons maintenant énoncer le résultat principal de cette section.

\begin{lemma}
\label{lemma-cotangent-morphism-spaces}
Dans la situation précédente, il existe un isomorphisme canonique
$$
L_{X/\Lambda} = 
L\pi_!(Li^*\Omega_{\mathcal{O}/\underline{\Lambda}}) =
L\pi_!(i^*\Omega_{\mathcal{O}/\underline{\Lambda}}) =
L\pi_!(\Omega_{\mathcal{O}/\underline{\Lambda}}
\otimes_\mathcal{O} \underline{\mathcal{O}}_X)
$$
dans $D(\mathcal{O}_X)$.
\end{lemma}

\begin{proof}
Observons d'abord que, pour tout objet $(U \to \mathbf{A})$ de
$\mathcal{C}_{X/\Lambda}$, la valeur du faisceau $\mathcal{O}$ est une algèbre
polynomiale sur $\Lambda$. Ainsi,
$\Omega_{\mathcal{O}/\underline{\Lambda}}$ est un $\mathcal{O}$-module plat,
et les deuxième et troisième égalités de l'énoncé sont vérifiées.

\medskip\noindent
D'après la remarque \ref{remark-compute-L-pi-shriek-spaces}, l'objet
$L\pi_!(\Omega_{\mathcal{O}/\underline{\Lambda}}
\otimes_\mathcal{O} \underline{\mathcal{O}}_X)$
se calcule comme la faisceautisation du complexe de préfaisceaux
$$
U \mapsto
\left(\Omega_{\mathcal{O}/\underline{\Lambda}}
\otimes_\mathcal{O} \underline{\mathcal{O}}_X\right)(\mathbf{A}_{\bullet, U})
=
\Omega_{P_{\bullet, U}/\Lambda} \otimes_{P_{\bullet, U}} \mathcal{O}_X(U) =
L_{\mathcal{O}_X(U)/\Lambda}
$$
avec les notations de la remarque \ref{remark-compute-L-pi-shriek-spaces}. La
remarque \ref{remark-map-sections-over-U} montre alors que
$L\pi_!(\Omega_{\mathcal{O}/\underline{\Lambda}}
\otimes_\mathcal{O} \underline{\mathcal{O}}_X)$
calcule le complexe cotangent de l'homomorphisme d'anneaux
$\underline{\Lambda} \to \mathcal{O}_X$ sur $X_\etale$. C'est le résultat
voulu d'après le lemme \ref{lemma-space-over-ring}.
\end{proof}






\section{Produits fibrés d'espaces algébriques et complexe cotangent}
\sectionmark{Produits fibrés et complexe cotangent}
\label{section-fibre-product}

\noindent
Soit $S$ un schéma. Soient $X \to B$ et $Y \to B$ des morphismes d'espaces
algébriques sur $S$. Considérons le produit fibré $X \times_B Y$, muni des
morphismes de projection $p : X \times_B Y \to X$ et
$q : X \times_B Y \to Y$. Dans cette section, nous étudions
$L_{X \times_B Y/B}$. L'essentiel des renseignements recherchés est contenu
dans le diagramme suivant
\begin{equation}
\label{equation-fibre-product}
\vcenter{
\xymatrix{
Lp^*L_{X/B} \ar[r] &
L_{X \times_B Y/Y} \ar[r] &
E \\
Lp^*L_{X/B} \ar[r] \ar@{=}[u] &
L_{X \times_B Y/B} \ar[r] \ar[u] &
L_{X \times_B Y/X} \ar[u] \\
 &
Lq^*L_{Y/B} \ar[u] \ar@{=}[r] &
Lq^*L_{Y/B} \ar[u]
}
}
\end{equation}
Explication : la ligne médiane est le triangle fondamental du lemme
\ref{lemma-triangle-ringed-topoi} associé aux morphismes
$X \times_B Y \to X \to B$. La colonne médiane est le triangle fondamental
associé aux morphismes $X \times_B Y \to Y \to B$. Ensuite, $E$ est un objet
de $D(\mathcal{O}_{X \times_B Y})$ qui « complète » le coin supérieur droit,
c'est-à-dire qui fait de la ligne supérieure et de la colonne droite des
triangles distingués. Un tel $E$ existe d'après Catégories dérivées,
proposition \ref{derived-proposition-9}, appliquée au carré inférieur gauche
(en plaçant $0$ à la position manquante). Plus explicitement, nous pourrions
par exemple définir $E$ comme le cône (Catégories dérivées, définition
\ref{derived-definition-cone}) de l'application de complexes
$$
Lp^*L_{X/B} \oplus Lq^*L_{Y/B} \longrightarrow L_{X \times_B Y/B}
$$
et obtenir les deux applications de but $E$ en appliquant TR3. Dans le cas
Tor-indépendant, l'objet $E$ est nul.

\begin{lemma}
\label{lemma-fibre-product-tor-independent}
Dans la situation précédente, si $X$ et $Y$ sont Tor-indépendants sur $B$,
alors l'objet $E$ de (\ref{equation-fibre-product}) est nul. Dans ce cas, nous
avons
$$
L_{X \times_B Y/B} = Lp^*L_{X/B} \oplus Lq^*L_{Y/B}
$$
\end{lemma}

\begin{proof}
Choisissons un schéma $W$ et un morphisme étale surjectif $W \to B$.
Choisissons un schéma $U$ et un morphisme étale surjectif
$U \to X \times_B W$. Choisissons un schéma $V$ et un morphisme étale
surjectif $V \to Y \times_B W$. Alors
$U \times_W V \to X \times_B Y$ est également étale surjectif. Il suffit donc
de montrer que la restriction de $E$ à $U \times_W V$ est nulle. D'après le
lemme \ref{lemma-compare-spaces-schemes} et Catégories dérivées des espaces,
lemme \ref{spaces-perfect-lemma-tor-independent}, nous sommes ramenés au cas
des schémas. En choisissant des ouverts affines convenables, nous nous ramenons
au cas des schémas affines. Le lemme \ref{lemma-morphism-affine-schemes} nous
ramène alors au cas d'un produit tensoriel d'anneaux, c'est-à-dire au lemme
\ref{lemma-tensor-product-tor-independent}.
\end{proof}

\noindent
Dans le cas général, nous pouvons dire ce qui suit de l'objet $E$.

\begin{lemma}
\label{lemma-fibre-product}
Soit $S$ un schéma. Soient $X \to B$ et $Y \to B$ des morphismes d'espaces
algébriques sur $S$. L'objet $E$ de (\ref{equation-fibre-product}) vérifie
$H^i(E) = 0$ pour $i = 0, -1$, et, pour tout point géométrique
$(\overline{x}, \overline{y}) : \Spec(k) \to X \times_B Y$, nous avons
$$
H^{-2}(E)_{(\overline{x}, \overline{y})} =
\text{Tor}_1^R(A, B) \otimes_{A \otimes_R B} C
$$
où $R = \mathcal{O}_{B, \overline{b}}$, $A = \mathcal{O}_{X, \overline{x}}$,
$B = \mathcal{O}_{Y, \overline{y}}$ et
$C = \mathcal{O}_{X \times_B Y, (\overline{x}, \overline{y})}$.
\end{lemma}

\begin{proof}
La formation du complexe cotangent commute au passage aux fibres et aux images
inverses ; voir les lemmes \ref{lemma-stalk-cotangent-complex} et
\ref{lemma-pullback-cotangent-morphism-topoi}. Remarquons que $C$ est un
hensélisé de $A \otimes_R B$.
$L_{C/R} = L_{A \otimes_R B/R} \otimes_{A \otimes_R B} C$
d'après les résultats de la section \ref{section-localization}. Ainsi, la fibre
de $E$ en notre point géométrique est le cône de l'application
$L_{A/R} \otimes C \to L_{A \otimes_R B/R} \otimes C$. Les assertions du
lemme résultent donc du cas des anneaux, c'est-à-dire du lemme
\ref{lemma-tensor-product}.
\end{proof}








\begin{multicols}{2}[\section{Autres chapitres}]
\noindent
Préliminaires
\begin{enumerate}
\item \hyperref[introduction-section-phantom]{Introduction}
\item \hyperref[conventions-section-phantom]{Conventions}
\item \hyperref[sets-section-phantom]{Théorie des ensembles}
\item \hyperref[categories-section-phantom]{Catégories}
\item \hyperref[topology-section-phantom]{Topologie}
\item \hyperref[sheaves-section-phantom]{Faisceaux sur les espaces}
\item \hyperref[sites-section-phantom]{Sites et faisceaux}
\item \hyperref[stacks-section-phantom]{Champs}
\item \hyperref[fields-section-phantom]{Corps}
\item \hyperref[algebra-section-phantom]{Algèbre commutative}
\item \hyperref[brauer-section-phantom]{Groupes de Brauer}
\item \hyperref[homology-section-phantom]{Algèbre homologique}
\item \hyperref[derived-section-phantom]{Catégories dérivées}
\item \hyperref[simplicial-section-phantom]{Méthodes simpliciales}
\item \hyperref[more-algebra-section-phantom]{Compléments d'algèbre}
\item \hyperref[smoothing-section-phantom]{Lissage des morphismes d'anneaux}
\item \hyperref[modules-section-phantom]{Faisceaux de modules}
\item \hyperref[sites-modules-section-phantom]{Modules sur les sites}
\item \hyperref[injectives-section-phantom]{Objets injectifs}
\item \hyperref[cohomology-section-phantom]{Cohomologie des faisceaux}
\item \hyperref[sites-cohomology-section-phantom]{Cohomologie sur les sites}
\item \hyperref[dga-section-phantom]{Algèbre différentielle graduée}
\item \hyperref[dpa-section-phantom]{Algèbre à puissances divisées}
\item \hyperref[sdga-section-phantom]{Faisceaux différentiels gradués}
\item \hyperref[hypercovering-section-phantom]{Hyperrecouvrements}
\end{enumerate}
Schémas
\begin{enumerate}
\setcounter{enumi}{25}
\item \hyperref[schemes-section-phantom]{Schémas}
\item \hyperref[constructions-section-phantom]{Constructions de schémas}
\item \hyperref[properties-section-phantom]{Propriétés des schémas}
\item \hyperref[morphisms-section-phantom]{Morphismes de schémas}
\item \hyperref[coherent-section-phantom]{Cohomologie des schémas}
\item \hyperref[divisors-section-phantom]{Diviseurs}
\item \hyperref[limits-section-phantom]{Limites de schémas}
\item \hyperref[varieties-section-phantom]{Variétés}
\item \hyperref[topologies-section-phantom]{Topologies sur les schémas}
\item \hyperref[descent-section-phantom]{Descente}
\item \hyperref[perfect-section-phantom]{Catégories dérivées de schémas}
\item \hyperref[more-morphisms-section-phantom]{Compléments sur les morphismes}
\item \hyperref[flat-section-phantom]{Compléments sur la platitude}
\item \hyperref[groupoids-section-phantom]{Schémas en groupoïdes}
\item \hyperref[more-groupoids-section-phantom]{Compléments sur les schémas en groupoïdes}
\item \hyperref[etale-section-phantom]{Morphismes étales de schémas}
\end{enumerate}
Thèmes de théorie des schémas
\begin{enumerate}
\setcounter{enumi}{41}
\item \hyperref[chow-section-phantom]{Homologie de Chow}
\item \hyperref[intersection-section-phantom]{Théorie de l'intersection}
\item \hyperref[pic-section-phantom]{Schémas de Picard des courbes}
\item \hyperref[weil-section-phantom]{Théories cohomologiques de Weil}
\item \hyperref[adequate-section-phantom]{Modules adéquats}
\item \hyperref[dualizing-section-phantom]{Complexes dualisants}
\item \hyperref[duality-section-phantom]{Dualité pour les schémas}
\item \hyperref[discriminant-section-phantom]{Discriminants et différentes}
\item \hyperref[derham-section-phantom]{Cohomologie de de Rham}
\item \hyperref[local-cohomology-section-phantom]{Cohomologie locale}
\item \hyperref[algebraization-section-phantom]{Géométrie algébrique et formelle}
\item \hyperref[curves-section-phantom]{Courbes algébriques}
\item \hyperref[resolve-section-phantom]{Résolution des surfaces}
\item \hyperref[models-section-phantom]{Réduction semi-stable}
\item \hyperref[functors-section-phantom]{Foncteurs et morphismes}
\item \hyperref[equiv-section-phantom]{Catégories dérivées de variétés}
\item \hyperref[pione-section-phantom]{Groupes fondamentaux de schémas}
\item \hyperref[etale-cohomology-section-phantom]{Cohomologie étale}
\item \hyperref[crystalline-section-phantom]{Cohomologie cristalline}
\item \hyperref[proetale-section-phantom]{Cohomologie pro-étale}
\item \hyperref[relative-cycles-section-phantom]{Cycles relatifs}
\item \hyperref[more-etale-section-phantom]{Compléments de cohomologie étale}
\item \hyperref[trace-section-phantom]{La formule des traces}
\end{enumerate}
Espaces algébriques
\begin{enumerate}
\setcounter{enumi}{64}
\item \hyperref[spaces-section-phantom]{Espaces algébriques}
\item \hyperref[spaces-properties-section-phantom]{Propriétés des espaces algébriques}
\item \hyperref[spaces-morphisms-section-phantom]{Morphismes d'espaces algébriques}
\item \hyperref[decent-spaces-section-phantom]{Espaces algébriques décents}
\item \hyperref[spaces-cohomology-section-phantom]{Cohomologie des espaces algébriques}
\item \hyperref[spaces-limits-section-phantom]{Limites d'espaces algébriques}
\item \hyperref[spaces-divisors-section-phantom]{Diviseurs sur les espaces algébriques}
\item \hyperref[spaces-over-fields-section-phantom]{Espaces algébriques sur des corps}
\item \hyperref[spaces-topologies-section-phantom]{Topologies sur les espaces algébriques}
\item \hyperref[spaces-descent-section-phantom]{Descente et espaces algébriques}
\item \hyperref[spaces-perfect-section-phantom]{Catégories dérivées d'espaces}
\item \hyperref[spaces-more-morphisms-section-phantom]{Compléments sur les morphismes d'espaces}
\item \hyperref[spaces-flat-section-phantom]{Platitude sur les espaces algébriques}
\item \hyperref[spaces-groupoids-section-phantom]{Groupoïdes dans les espaces algébriques}
\item \hyperref[spaces-more-groupoids-section-phantom]{Compléments sur les groupoïdes dans les espaces}
\item \hyperref[bootstrap-section-phantom]{Amorçage}
\item \hyperref[spaces-pushouts-section-phantom]{Sommes amalgamées d'espaces algébriques}
\end{enumerate}
Thèmes de géométrie
\begin{enumerate}
\setcounter{enumi}{81}
\item \hyperref[spaces-chow-section-phantom]{Groupes de Chow des espaces}
\item \hyperref[groupoids-quotients-section-phantom]{Quotients de groupoïdes}
\item \hyperref[spaces-more-cohomology-section-phantom]{Compléments sur la cohomologie des espaces}
\item \hyperref[spaces-simplicial-section-phantom]{Espaces simpliciaux}
\item \hyperref[spaces-duality-section-phantom]{Dualité pour les espaces}
\item \hyperref[formal-spaces-section-phantom]{Espaces algébriques formels}
\item \hyperref[restricted-section-phantom]{Algébrisation des espaces formels}
\item \hyperref[spaces-resolve-section-phantom]{Résolution des surfaces revisitée}
\end{enumerate}
Théorie des déformations
\begin{enumerate}
\setcounter{enumi}{89}
\item \hyperref[formal-defos-section-phantom]{Théorie formelle des déformations}
\item \hyperref[defos-section-phantom]{Théorie des déformations}
\item \hyperref[cotangent-section-phantom]{Le complexe cotangent}
\item \hyperref[examples-defos-section-phantom]{Problèmes de déformation}
\end{enumerate}
Champs algébriques
\begin{enumerate}
\setcounter{enumi}{93}
\item \hyperref[algebraic-section-phantom]{Champs algébriques}
\item \hyperref[examples-stacks-section-phantom]{Exemples de champs}
\item \hyperref[stacks-sheaves-section-phantom]{Faisceaux sur les champs algébriques}
\item \hyperref[criteria-section-phantom]{Critères de représentabilité}
\item \hyperref[artin-section-phantom]{Axiomes d'Artin}
\item \hyperref[quot-section-phantom]{Espaces de Quot et de Hilbert}
\item \hyperref[stacks-properties-section-phantom]{Propriétés des champs algébriques}
\item \hyperref[stacks-morphisms-section-phantom]{Morphismes de champs algébriques}
\item \hyperref[stacks-limits-section-phantom]{Limites de champs algébriques}
\item \hyperref[stacks-cohomology-section-phantom]{Cohomologie des champs algébriques}
\item \hyperref[stacks-perfect-section-phantom]{Catégories dérivées de champs}
\item \hyperref[stacks-introduction-section-phantom]{Introduction aux champs algébriques}
\item \hyperref[stacks-more-morphisms-section-phantom]{Compléments sur les morphismes de champs}
\item \hyperref[stacks-geometry-section-phantom]{La géométrie des champs}
\end{enumerate}
Thèmes de théorie des espaces de modules
\begin{enumerate}
\setcounter{enumi}{107}
\item \hyperref[moduli-section-phantom]{Champs de modules}
\item \hyperref[moduli-curves-section-phantom]{Espaces de modules de courbes}
\end{enumerate}
Divers
\begin{enumerate}
\setcounter{enumi}{109}
\item \hyperref[examples-section-phantom]{Exemples}
\item \hyperref[exercises-section-phantom]{Exercices}
\item \hyperref[guide-section-phantom]{Guide bibliographique}
\item \hyperref[desirables-section-phantom]{Desiderata}
\item \hyperref[coding-section-phantom]{Style de programmation}
\item \hyperref[obsolete-section-phantom]{Obsolète}
\item \hyperref[fdl-section-phantom]{Licence de documentation libre GNU}
\item \hyperref[index-section-phantom]{Index généré automatiquement}
\end{enumerate}
\end{multicols}


\bibliography{my}
\bibliographystyle{amsalpha}

\end{document}
